% concatenated from GM_new.tex, Appendix_new.tex
\documentclass[12pt]{article}
\usepackage{comment}
\usepackage{rotating}
\usepackage[utf8]{inputenc}
\usepackage{mathtools}
\usepackage{minted}
%\newcommand{\indep}{\perp \!\!\! \perp}
\usepackage{fullpage}
\usepackage{amsmath,amssymb,amsfonts,pdfpages}
\usepackage{amsthm}
\usepackage{enumerate}
\usepackage[utf8]{inputenc}
%\usepackage{algorithmic}
\usepackage{minted}
\usepackage{color}
\usepackage{graphicx}
\PassOptionsToPackage{hyphens}{url}\usepackage{hyperref}
\usepackage{caption,subcaption}
\usepackage[symbol]{footmisc}
\usepackage[useregional]{datetime2}
\usepackage{verbatim}
\usepackage{comment}
\usepackage[export]{adjustbox}
\usepackage{tikz}
\usepackage{algorithm}
\usepackage{algpseudocode}
\usepackage[sort]{natbib}
\usepackage[margin=0.6in]{geometry}
\usepackage{makecell} 



\newcommand{\RR}{\mathbb{R}}
\newcommand{\EE}{\mathbb{E}}
\newcommand{\PP}{\mathbb{P}}
\newcommand{\NN}{\mathbb{N}}

\newtheorem{lemma}{Lemma}
\newtheorem{thm}{Theorem}
\newtheorem{corollary}{Corollary}
\newtheorem{proposition}{Proposition}

\theoremstyle{definition}
\newtheorem{example}{Example}
\newtheorem{model}{Model} % Tianyi's check
\newtheorem{property}{Property} % Tianyi's check

% Define two environments with independent counters
\newtheorem{lproperty}{Property}
\newtheorem{aproperty}{Property}

% Redefine the numbering for each to include a prefix letter
\renewcommand{\thelproperty}{L.\arabic{lproperty}}
\renewcommand{\theaproperty}{A.\arabic{aproperty}}

\newtheorem{conjecture}{Conjecture} % Tianyi's check
\newtheorem{definition}{Definition}
\newtheorem{remark}{Remark}

\newcommand{\R}{\mathbb{R}}
\newcommand{\Z}{\mathbb{Z}}
\newcommand{\E}{\mathbb{E}}
\newcommand{\bA}{\mathbf{A}}
\newcommand{\bB}{\mathbf{B}}
\newcommand{\bD}{\mathbf{D}}
\newcommand{\bE}{\mathbf{E}}
\newcommand{\bJ}{\mathbf{J}}
\newcommand{\bH}{\mathbf{H}}
\newcommand{\bI}{\mathbf{I}}
\newcommand{\bL}{\mathbf{L}}
\newcommand{\bM}{\mathbf{M}}
\newcommand{\bN}{\mathbf{N}}
\newcommand{\bO}{\mathbf{O}}
\newcommand{\bP}{\mathbf{P}}
\newcommand{\bR}{\mathbf{R}}
\newcommand{\bS}{\mathbf{S}}
\newcommand{\bT}{\mathbf{T}}
\newcommand{\bQ}{\mathbf{Q}}
\newcommand{\bU}{\mathbf{U}}
\newcommand{\bV}{\mathbf{V}}
\newcommand{\bW}{\mathbf{W}}
\newcommand{\bX}{\mathbf{X}}
\newcommand{\bY}{\mathbf{Y}}
\newcommand{\bZ}{\mathbf{Z}}

\newcommand{\be}{\mathbf{e}}
\newcommand{\bv}{\mathbf{v}}
\newcommand{\bx}{\mathbf{x}}
\newcommand{\by}{\mathbf{y}}

\newcommand{\calB}{\mathcal{B}}
\newcommand{\calG}{\mathcal{G}}
\newcommand{\calL}{\mathcal{L}}
\newcommand{\calN}{\mathcal{N}}
\newcommand{\calO}{\mathcal{O}}
\newcommand{\calP}{\mathcal{P}}
\newcommand{\calW}{\mathcal{W}}
\newcommand{\mX}{\mathcal{X}}
\newcommand{\mY}{\mathcal{Y}}
\newcommand{\mZ}{\mathcal{Z}}
\newcommand{\mC}{\mathcal{C}}

\newcommand{\Xhat}{\hat{\bX}}
\newcommand{\dhat}{\hat{d}}
\newcommand{\Khat}{\hat{K}}
\newcommand{\Yhat}{\hat{\bY}}
\newcommand{\Zhat}{\hat{\bZ}}
\newcommand{\secmomd}{\lambda_d({\bf \Delta})}
\newcommand{\omni}{\operatorname{OMNI}}
\newcommand{\zeromx}{\mathbf{0}}
\newcommand{\Xomni}{\bX_{\omni}}
\newcommand{\supp}{\operatorname{supp}}
\newcommand{\tr}{\operatorname{tr}}

\newcommand{\Ptilde}{\tilde{\bP}}
\newcommand{\Qtilde}{\tilde{\bQ}}
\newcommand{\Wtilde}{\tilde{\bW}}
\newcommand{\Wn}{\bw_m}
\newcommand{\Wntilde}{\Wtilde_n}
\newcommand{\Va}{\bV_1}
\newcommand{\Vb}{\bV_2}
%\newcommand{\Wnprime}{\bW^\prime_n}
\newcommand{\Xtilde}{\tilde{\bX}}
\newcommand{\Xstar}{\bX^*}
\newcommand{\Xtildesi}{\Xtilde_i^{(s)}}
\newcommand{\Zstar}{\bZ^*}
\newcommand{\utilde}{\tilde{u}}
\newcommand{\tti}{2 \rightarrow \infty}
\newcommand{\Abar}{\bar{\bA}}
\newcommand{\Xbar}{\bar{\bX}}
\newcommand{\jbar}{\bar{j}}
\newcommand{\rbar}{\bar{r}}
\newcommand{\tbar}{\bar{t}}
\newcommand{\CX}{C_{\bX}}
\renewcommand*{\thefootnote}{\fnsymbol{footnote}}
\setcounter{footnote}{0}

\newcommand{\balpha}{\mathbf{\alpha}}
\newcommand{\bDelta}{\mathbf{\Delta}}
\newcommand{\bSigma}{\mathbf{\Sigma}}
\newcommand{\Sigmatilde}{\tilde{\bSigma}}

\newcommand{\alphavec}{\vec{\balpha}}

\newcommand{\UM}{\bU_{\bM}}
\newcommand{\UA}{\bU_{\bA}}
\newcommand{\ULA}{\bU_{\calL(\bA)}}
\newcommand{\SLA}{\bS_{\calL(\bA)}}
\newcommand{\SLP}{\bS_{\calL(\bP)}}
\newcommand{\ULP}{\bU_{\calL(\bP)}}
\newcommand{\SM}{\bS_{\bM}}
\newcommand{\SA}{\bS_{\bA}}
\newcommand{\UPt}{\bU_{\Ptilde}}
\newcommand{\SPt}{\bS_{\Ptilde}}
\newcommand{\UP}{\bU_{\bP}}
\newcommand{\SP}{\bS_{\bP}}

\newcommand{\SBM}{\operatorname{SBM}}
\newcommand{\ASE}{\operatorname{ASE}}
\newcommand{\Dir}{\operatorname{Dir}}
\newcommand{\HSBM}{\operatorname{HSBM}}
\newcommand{\RDPG}{\operatorname{RDPG}}
\newcommand{\JRDPG}{\operatorname{JRDPG}}
\newcommand{\indic}{\mathbb{I}}
\newcommand{\rank}{\operatorname{rank}}
\newcommand{\Bern}{\operatorname{Bernoulli}}
\newcommand{\Var}{\operatorname{Var}}
\newcommand{\Cov}{\operatorname{Cov}}

\newcommand{\inlaw}{\xrightarrow{\calL}}
\newcommand{\inprob}{\xrightarrow{P}}
\newcommand{\iid}{\stackrel{\text{i.i.d.}}{\sim}}

\newcommand{\tM}{t_M}
\newcommand{\tDqi}{t_{\mxD,q,i}}
\newcommand{\tDqri}{t_{\mxD,q,r,i}}
\newcommand{\tZqij}{t_{Z,q,i,j}}
\newcommand{\tZqrij}{t_{Z,q,r,i,j}}

\newcommand{\mxE}{\mathbf{E}}
\newcommand{\mxD}{\mathbf{D}}

\newcommand{\ErdosRenyi}{Erd\H{o}s-R\'{e}nyi }
\newcommand{\Furedi}{F\"{u}redi }
\newcommand{\Erdos}{Erd\H{o}s}
\newcommand{\Renyi}{R\'{e}nyi}
\newcommand{\Komlos}{Koml\'{o}s }
\newcommand{\mclustase}{GMM \circ ASE}
\newcommand{\smclustase}{SemiparGMM \circ ASE}
\newcommand{\thetahat}{\widehat{\theta}}
\newcommand{\pihat}{\widehat{\pi}}
\newcommand{\rhohat}{\widehat{\rho}}
\newcommand{\Chat}{\widehat{\mathcal{C}}_{KC}}
\newcommand{\CKC}{\mathcal{C}_{KC}}
\newcommand{\synthMB}{\texttt{synthMB}}
\newcommand{\semiparSBM}{\texttt{StructuredSBM}}
\newcommand{\synthSemiparMB}{\texttt{synthSemiparMB}}

\newcommand{\Bhat}{\widehat{B}}
\newcommand{\Phat}{\widehat{P}}
\newcommand{\ind}{\perp\!\!\!\!\perp} 

\newcommand{\pre}{\mathrm{pre}}
\newcommand{\post}{\mathrm{post}}

\newcommand{\inddmv}[1][]{\mathop{\mbox{ind-$d_{MV}^{#1}$}}}
\newcommand{\procwp}[1][2]{\mathop{\mbox{proc-$W_#1$}}}
\newcommand{\procwphat}[1][2]{\mathop{\mbox{proc-$\hat{W}_#1$}}}

\newcommand{\psim}{\psi^{(m)}}
\newcommand{\tpsim}{\tilde{\psi}^{(m)}}
\newcommand{\hpsim}{\hat{\psi}^{(m)}}
\newcommand{\thpsim}{\tilde{\hat{\psi}}^{(m)}}
\makeatletter
\let\@fnsymbol\@arabic
\makeatother

%\usepackage{endfloat}
\usepackage{marginnote}

\begin{document}

\title{Vertex misalignment and changepoint localization in network time series}

\author{Tianyi Chen\footnote{Co-first author}, Mohammad Sharifi Kiasari\footnotemark[1], Sijing Yu, Youngser Park, \\ Avanti Athreya, Vince Lyzinski, Carey Priebe, Zachary Lubberts}

% \affil[1]{Department of Applied Mathematics and Statistics, Johns Hopkins University}
% \affil[2]{Department of Statistics, University of Virginia}
% \affil[3]{Center for Imaging Science, Johns Hopkins University}

\date{}

\maketitle


\begin{abstract}

%AA: Changepoint localization in a time series of networks generally relies on accurate vertex correspondence between network realizations at different times. However, such vertex alignments are often misspecified or even unknown. We prove results on the relationship between graph matching and optimal transport, two mechanisms for mitigating errors from misalignment, and we construct two illustrative models for network evolution, each with a similar changepoint. We compare techniques for changepoint localization, ranging from the simple network statistic of average degree to the more involved and recently developed procedure of Euclidean mirrors. In one model, vertex misalignment causes comparatively little error, and in the other, it seriously impairs localization, although the Euclidean mirror procedure can still extract meaningful signal. We show how misalignment between network realizations at different times can effectively weaken their underlying correlation, impeding inference procedures that rely on accurate correlation estimates. We present simulations to illustrate these effects on approaches to localization.

Inference for time series of networks often relies on accurate vertex correspondence between network realizations at different times. In practice, however, such vertex alignments can be misspecified or unknown. We study the impact of vertex alignment on changepoint localization for dynamic networks through two illustrative models, each with a similar changepoint, with
%, but with a key distinction on the informational content between joint and marginal distributions. 
the key distinction being whether changepoint information is contained in marginal or joint distributions of the time-varying latent positions.
We compare localization techniques ranging from the simple network statistic of average degree to the modern procedure of Euclidean mirrors. 
In one model, vertex misalignment causes little error, and in the other, it impairs localization in ways that cannot be corrected through graph matching or optimal transport, which we show are closely related in this setting. 
Our results demonstrate that robust network inference necessitates reckoning with the subtle interplay of marginal and joint information in the observed network time series.


%with respect to the marginal and joint distributions, and utilizing appropriately robust estimation procedures while coming from an informed perspective about the limitations of shuffled data.



%reckoning with subtle distinctions in how informational content depends on joint and marginal distributions, 

%are core to robust network inference.
%Despite this, the Euclidean mirror procedure can still extract meaningful signal when the proportion of misaligned vertices is small. 
%We examine how relevant information about the network time series may depend on both the marginal and joint distributions across times. 
%We analyze the relationship between graph matching and optimal transport, two approaches that mitigate vertex misalignment, and present simulations to illustrate these contrasting effects on approaches to localization.



%We describe changepoints in network evolutionary processes that can be successfully localized without knowledge of vertex alignment and, conversely, other changepoints that cannot be localized at all.
% This study addresses the challenge of changepoint localization in evolving time series of networks, particularly in the context of vertex misalignment across different time points. We explore two models of network evolution: one where misalignment has minimal impact on inference, and another where it significantly hampers changepoint detection. We determine conditions under which changepoints can be successfully localized without vertex alignment, as well as scenarios where localization fails. Additionally, we discuss graph matching methods, which can be viewed as a form of noisy optimal transport, that can sometimes mitigate errors from misalignment but may also fail to provide any benefit under certain model. We present simulations that illustrate these varying effects on localization efficacy.

  
% Identifying changepoints in time series of graphs is crucial. The Euclidean mirror, a recently developed method, can effectively represent the dynamics of such time series. In particular, we are interested in solving the problem of localizing the first-order changepoint in a time series graph (TSG) using the Euclidean mirror. In this paper we mainly study how losing vertex alignment between graphs impact the changepoint localization task. To investigate this, we propose two models: the London model and the Atlanta model as extreme cases. Both models have first-order changepoints. The London model is robust to missing vertex alignment, while the Atlanta model requires vertex alignment for the first-order changepoint. Next, we investigate whether graph matching can recover the true alignment or help localize changepoints when the graphs are shuffled. We find that the $d_{MV}$ distance with the graph matching solution are similar to the Wasserstein distance. We observe that the London model using Wasserstein distance reveals the first-order changepoint more effectively than using $d_{MV}$ with the true alignment. Conversely, the Atlanta model using Wasserstein distance does not provide any signal about the first-order changepoint. Finally we can further prove the changepoint in Atlanta model is asymptotically nonidentifiable after shuffling the observed graphs. What's more, we can show any model with parameter that is only identifiable with joint distribution and nonidentifiable with marginal distribution will be asymptotically nonidentifiable after shuffling. These findings inspire potential applications of real data graph matching. 

\end{abstract}

%* Office of Naval Research (ONR) Science of Autonomy award number N00014-24-1-2278
%* Air Force Office of Scientific Research (AFOSR) Complex Networks award number FA9550-25-1-0128
%* Defense Advanced Research Projects Agency (DARPA) Artifical Intelligence Quantified award number HR00112520026
%* Microsoft Research


% \section{To-do}

% \begin{itemize}
%     \item use $\EE[]$ instead of $\EE()$
%     \item add the theorem about percentage shuffling
% %    \item add the in-fill asym square 1-d Euclidean realizable for Atlanta model.

    
% \end{itemize}

% Outline
% \begin{enumerate}[(a)]
%     \item Overall theme: robustness to absent or errorful vertex alignment for changepoint localization in evolving networks
%     \item Standard intro \& background
%     \item Simulated data story first (swarm data)
%     \item We need to formalize this---background and definitions, LPP TSG, dissimil, CMDS, mirror,changepoint. no estimation stuff.
%     \item Shuffling $\Rightarrow$ independence
%     \item Tale of two models (contextualize)
%     \item Theorem 1: Distances for London and Atlanta, with and without shuffling.
%     \item Figure: Table of mirrors for London and Atlanta, with and without shuffling
%     \item Theorem 2: Mirrors for London with and without shuffling (asymp pw linear); Atlanta with and without shuffling (asymp pw linear, useless orthogonal points)
%     \item Corollary: $d_{MV}$ with partial shuffling = convex combo of $d_{MV}$ and shuffled-$d_{MV}$
%     \item \textcolor{blue}{Atlanta nonidentifiability/hardness: \emph{Do we even need this?}}
%     \item Wasserstein distance and optimal transport, theorems, this is an alternative method for measuring distance when you don't know the alignment
%     \item Figure: Mirrors for London and Atlanta with Wasserstein distance
%     \item Theorem: Mirrors for London and Atlanta
%     \item Simulation results 


%     \begin{comment}
%     \item Summary of London model and way to localize changepoint
%     \item Shuffling does not hurt you in London: simulation
%     \item In London model, labels are a nuisance semiparametric function
%     \item We should use Wasserstein \& talk about Zach's theorem: exact 1-d realizability vs dMV with asymptotic 1-d realizability; simulations; 
%     \item In real life shuffling does hurt you (TC has ONR SofA from Youngser yesterday!)
%     \item Atlanta model as a concrete one in which vertex correspondence is essential
%     \item Partial shuffling vs total derangements
%     \item Thresholds for information retained in partial shuffling and impact on localization
%     \item Ways to compute a "restricted" dMV based on unrestricted spectral embedding and then use only the portion associated to vertices that are not shuffled
%     \item Approaches to fixing this: graph matching, wasserstein
%     \item House political data
%     \end{comment}

% \end{enumerate}



% \newpage


\begin{comment}
% \textit{words for practitioner:

% When you have a time series of graphs and you are looking for a first-order changepoint, and if you decide you want to embed the TSG into a low dimension Euclidean time series. Even if you know the 1-1 correspondence, firstly check the average degree and Wasserstein distance induced mirror. We learned from London model, if the 1-1 correspondence is a nuisance parameter to your task, ignore it can lead to better performance. If these doesn't give you meaningful structure or just flat, you can try the $d_{MV}$ induced mirror looking for finer structure. 
% Secondly if you don't know 1-1 correspondence, but you know 1-1 correspondence for a specific common subset of vertices across the time, then after trying average degree and Wasserstein and find nothing, you can try construct $d_{MV}$ mirror in the following way: first get ASE results using all vertices, then discard the vertices where you don't know then 1-1 correspondence, construct distance matrix only with the vertices with known 1-1 correspondence, this potentially gives as good as you known all 1-1 correspondence. In this case you can also try graph matching on the whole graph, maybe it works. We have empirical evidence of it.
% Thirdly when 1-1 correspondence is lost for all, the 
% }
% Revised practitioner notes:
%
% When analyzing a time series of graphs to detect a first-order changepoint,
% you can first embed the TSG into a low-dimensional Euclidean time series.
% Even if the vertex correspondence is known, start by examining the average degree
% and the Wasserstein distance–induced mirror. As observed in the London model,
% when the one-to-one correspondence acts merely as a nuisance, ignoring it may boost performance.
%
% If these summaries do not reveal meaningful or non-flat structure, consider the 
% d_{MV}–induced mirror to capture finer changes.
%
% In situations where the full one-to-one correspondence is unknown but you have 
% it for a common subset of vertices across time, proceed as follows:
% first, obtain the ASE embedding using all vertices;
% then, discard vertices without known correspondence and construct the distance 
% matrix using only those vertices with confirmed one-to-one alignment.
% This approach can potentially perform as well as using full correspondence.
%
% Finally, if one-to-one correspondence is entirely lost for all vertices, alternative 
% strategies such as graph matching or other alignment methods should be considered.


$t^*$ is identifiable in the Atlanta model, but in the shuffled Atlanta model, that is at each time step it jumps to the other vertex's latent position(i will describe later.) $t^*$ is not identifiable, that is no matter where $t^*$ is, it always yields the same joint distribution for TSG. Thus as a nuisance parameter the 1-1 correspondence in Atlanta model if not the true one, will cause $t^*$ non identifiable. 

This shows if we lost the 1-1 correspondence in the Atlanta model, then it is the shuffled Atlanta model and estimate $t^*$ is not feasible. 

However in London model, and shuffled London model, $t^*$ are both identifiable. And as a nuisance parameter, it is ``orthogonal" to $t^*$ in London model. 


%The ratio is not a test for whether your 1-1 correspondence is true or not, it is merely a quantity to show how important the true 1-1 correspondence is for the model. For practitioners, you got the real data and if you really don’t know whether the graphs are aligned, then you can just look at the quantity. 
%1, If the empirical quantity is distributed away from 1(this is really the hard part, how do we define away and close), then this means your data is really different from the completely shuffled case, so you somehow do have to some extend some correspondence. 
%2, On the other hand, if you data has the quantity distributed close to 1, it means even if you really have the alignment, it doesn’t matter, the true alignment already looks like the shuffled case. Or you just have the shuffled data and you don’t have true alignment. 

%And then further, you do the subsequent inference task on mirror, and find what you want/expect, then it means the signal is strong enough and you don’t need alignment, and you don’t care whether your alignment is true. Then if the alignment is actually expensive to get, we can suggest for this task, in the future you don’t collect this info anymore.  
%However if the practitioner is not satisfied with the subsequent inference result on mirror, then we could say, maybe, just one possible reason, is because your data is actually not aligned, i.e the data is actually shuffled already, your numerator is shuffled-dMV. Then you should go back to look at the data again see if we can find the alignment. 



If you know you are under Atlanta model, and you know some part of 1-1 correspondence, you really should do restriced dMV, meaning completely discard the shuffling part. 



\end{comment}



\section{Introduction}



Changepoint localization and detection in time series data is a classical problem in statistical inference \cite{shumway2000time,box2015time,kedem2005regression, verzelen2023optimal, padilla2021optimal}. The analogue for network data---namely, identifying if and when a time series of networks undergoes a distributional change---is a rapidly developing subfield of network inference, with theoretical and practical relevance across disciplines. However, statistical changepoint analysis for networks presents distinct challenges, because network data is inherently non-Euclidean in its structure and because network evolution involves dependencies between underlying network features as well as correspondences across vertices over time.  Theoretical facets of the analysis of time-varying networks encompass random matrix theory and eigenvector fluctuations \cite{cape2019signal,agterberg2022entrywise,tang2025eigenvector}, tensor analysis 
\cite{macdonald2022latent,pensky2025signed,chen2022analysis}, minimax bounds and thresholds of detectability \cite{cai2021optimal,zhang2018tensor,lei2024computational}, and optimization and graph matching \cite{fishkind2019seeded,arroyo2021graph,saad2021graph, BergmannHerzogJasa:2024}. On a practical level, the problem of identifying shifts in network dynamics has broad applicability, for example in organizational networks for corporations \cite{zuzul2021dynamic}; neural development for organoids \cite{chen_worm}; and transportation and logistical networks \cite{agterberg2022joint}. 

The non-Euclidean nature of network data can be addressed by considering lower-dimensional spectral embeddings of multiple networks or a network adjacency tensor \cite{paul2020spectral,jing2021community,pantazis2022importance,macdonald2022latent,wang2023multilayer,athreya2025euclidean,agterberg2022joint,padilla2022change,gallagher2021spectral}, but this requires accurate alignment between vertices across networks.
In \cite{athreya2025euclidean}, the authors develop the methodology of Euclidean mirrors, which involves the construction of a dissimilarity measure between time-indexed pairs of adjacency matrices $\bA_{t}$, $\bA_{t'}$ of a network time series. 
For certain types of dissimilarity measures and under certain model assumptions, a low-dimensional embedding of the matrix of dissimilarities---known as {\em Euclidean mirror} or {\em iso-mirror}---can help identify when the sequence of underlying networks undergoes a significant change \cite{chen2023discovering, chen2024euclidean}. Existing approaches to dynamic networks, including that of the Euclidean mirror, often presume a fixed vertex set and a known correspondence across vertices. In practice, both of these assumptions can fail. 
As networks evolve, vertex sets may themselves change (see, for example, the cases in \cite{kivela2014multilayer} that do not have equal size across layers); and even when vertex sets remain the same, correct alignments may be unknown, errorfully observed, or corrupted \cite{priebe2015statistical, lyzinski16:_infoGM, arroyo2021maximum}.
For example, in the analysis of neural organoid time series, the vertex set of neural cells expands as the organoid grows, while constant agitation due to in-vitro culturing makes tracking vertex correspondence over time impossible.
In \cite{chen2023discovering}, the authors mitigate the lack of vertex correspondence in organoids by using graph matching on the padded adjacency matrices and applying Euclidean mirrors for exploratory data analysis on organoid time series. 
The empirical insights so obtained lead to further theoretical questions.
%An initial step toward a deeper understanding is to explore, the impact of vertex misalignment on changepoint localization.
%[AA: ADD PARAGRAPH HERE ABOUT VERTEX CORRESPONDENCE AND ITS IMPORTANCE IN INFERENCE, WITH DISCUSSION OF GRAPH MATCHING, WASSERSTEIN, TRADEOFFS. VL: HELP ME! ZZZ]
There is a large body of work devoted to recovering an unknown vertex correspondence across networks---this is the well-studied graph matching problem \cite{gms1,gms2,gms3}---
and the effect of label errors has been studied in the context of graph and vertex classification
\cite{vogelstein2015shuffled,chen_worm}, joint graph clustering \cite{racz2021correlated,racz2,chai2024efficient}, and two-sample graph hypothesis testing \cite{saxena2025lost,qi2025asymptotically}.
The present work is, to our knowledge, the first to systematically consider the effect of label error in time-series analysis and change-point localization.

To formalize these ideas, we consider a collection of random networks on the same vertex set, indexed by time. For this network time series, we assume that associated to each vertex $i$ in each network at time $t$ is a (typically unobserved) low-dimensional vector, called the {\em latent position}, and that the probability of connections between nodes in the network depends on these latent positions. We assume these latent positions evolve according to some underlying stochastic process. A central question is whether, if the underlying stochastic process---called a {\em latent position process} (LPP)---
changes in some significant way, we can detect this by appropriate analysis of the observed network data itself. The observed network data consists of a time-indexed collection of adjacency matrices, each entry of which is a Bernoulli random variable with success probability determined by the latent positions of the respective nodes. The observed network adjacency $\bA_t$ at time $t$ is a noisy version of the matrix of connection probabilities $\bP_t$ at time $t$. The matrix of connection probabilities $\bP_t$ is itself a compression of the information in the underlying realizations of the latent position processes. We do not observe the $\bP_t$ matrices, nor the realizations of the underlying latent position process. However, because the edge distribution of a latent position network is governed by the latent position process, changepoints for the distribution of the network can be associated to changepoints in the latent position process.

Our goal is to examine the impact of vertex misalignments for changepoint localization in a conditionally-independent time series of networks governed by a latent position process.
To that end, we prove several results on the impact of vertex shuffling on latent position processes, specifically describing how this weakens dependence across time. We further prove key results relating graph matching and optimal transport in this setting, and the impact of these two procedures for rectifying information loss from vertex mislabeling. We then  describe two distinct latent position processes, dubbed {\em London} and {\em Atlanta}. In both latent position processes, there is a concrete distributional change at a specific timepoint.
In each case, the changepoint in the underlying respective LPP gives rise to a corresponding  changepoint in the associated time series of networks. 
In the London model, the vertex labels are effectively uninformative and this changepoint can be localized even when vertex correspondence is lost. 
Moreover, several network statistics or functionals, ranging from simple ones such as average degree to more complicated ones involving Euclidean mirrors, can localize the changepoint. In the Atlanta model, by contrast, a shuffled vertex correspondence can severely degrade changepoint localization, sometimes irrecoverably.
Here, the Euclidean mirror methodology can correctly localize the changepoint under a known correspondence, and retains some signal even under partial vertex shuffling. 
However, in the presence of substantial shuffling, we show that these localization approaches will generally fail. 
Further, efforts to reconstruct the latent alignment, such as via graph matching or optimal transport, provide scant improvement.  
% We analyze the stark differences between the London and Atlanta models in terms of robustness of localization to vertex shuffling. 

While both models are simplifications, they present novel edge cases. 
In London, the vertex labels do not carry information relevant to changepoint localization, while in Atlanta, changepoint recovery cannot proceed without the vertex correspondence. This distinction can be characterized in terms of the marginal and joint distributions across time, for each network time series and its underlying LPP. We prove that vertex relabeling can weaken correlation between latent position processes at different times; thus, for cases in which changepoint localization depends on accurately gauging correlation, a loss of vertex correspondence can degrade downstream inference. In the London model, changepoint localization can proceed via accurate estimation of marginal properties of the underlying latent position process, and there is no information loss from vertex mislabelling. In the Atlanta model, changepoint localization depends on understanding joint distributions, and the correlation-weakening effect of vertex mislabelling ranges from navigable (when a small number of vertices are mislabelled) to insurmountable (once a nontrivial fraction of vertices are incorrectly matched).
The importance of vertex correspondence in real data lies, of course, between these two extremes. Understanding the nuances of models lends insight into representations for real data and changepoint detection methodology in the presence of vertex mislabeling.



 
{\bf Outline of paper}. We structure the paper as follows. In Section \ref{sec:Background_and_notation}, we encapsulate notation, definitions, and background on latent position process time series of graphs (LPPTSG) and Euclidean mirrors. In Section \ref{sec:Main_results_GM_OT}, we prove our main theorems on the impact of vertex misalignments for the latent position process and the relationship between graph matching and optimal transport in this context.  We define our two principal latent position process models, the London and Atlanta models, and then establish results on changepoint localization for the associated LPPTSGs. In Section \ref{sec:Simulations}, we demonstrate simulation evidence of the impact of vertex misalignment in both London and Atlanta models, as well as an exploratory data analysis of simulated data on swarms. In Section \ref{sec:Conclusion}, we provide a summary and discussion of open problems and ongoing work.  





% \begin{enumerate}
% \item  Why is changepoint localization in dynamic networks important?
% \item What are core practical applications? Organoids (CITE ourselves), organizational networks (CITE carey \& pahnke, ourselves), social networks (Clauset \& Peel), infrastructure networks (CITE Yu Yi, Runbing Zheng, Lubberts and Agterberg). 
% \item Primary approaches---methodology includes spectral decompositions, tensor decompositions, etc. etc (see first-order change point paper for further citations)
% \item Vertex misalignment and its challenges in general network inference. Ways to address misalignment: graph matching (CITE VL and collaborators): model-dependent, downstream inference-task dependent
% \item cite shuffled regression 

% %it is complicated and subtle
% \item Evolving networks and localization---also model dependent and inference-task dependent. We have a tale of two models: London,  where misalignment doesn't hurt; Atlanta, where it does.
% localization case (average degree gives us sufficient information can really be a problem 
% \item Graph matching and optimal transport---introductory description
% \item  Verbal description of LPP, verbal description of first order change point, verbal description of mirror and localization procedure 
% \end{enumerate}




% LPP is the canonical/natural way to study TSG. Thus losing correspondence in LPP is exactly losing correspondence in a joint distribution and this paper quantify the lost, it will give you product measure. 

% Also vertex correspondence is not observable in many real data applications, for example for neural organoid. 


% What's more the idea of this paper lie in the intersection of graph inference setting and one of the oldest topic in the statistics: pair versus unpair. 
% (Also if we only consider the first moment, pair is the same with the unpair, however the second moments makes real difference, paired has smaller variance: 
% %when $\Cov[X,Y]>0$,
% %%$$
% %\Var[X-Y] \leq \Var[X'-Y']
% %\quad \text{where } X' \overset{\mathcal{L}}{=}X, Y' \overset{\mathcal{L}}%{=}Y, X' \ind Y'.$$
% %%)

% the London model and Atlanta model serves not only as LPP for generative TSG and give a changepoint inference problem, the inference problem of $(p,q,t^*)$ is meaningful even without the TSG setting. ANd it provides a unique perspective to the pair versus unpair question, which is usually hidden under the rug in real data application. 





\section{Background and notation}\label{sec:Background_and_notation}

\subsection{Notation}

For $n \in \NN$, the natural numbers, let $[n] : = \{1,2,\cdots,n\}$. Let $e_i$ denote the $i$th standard basis vector in $\mathbb{R}^d$, and let $v=\left(v_1, v_2, \cdots, v_d\right)^\top$ denote the column vector $v=\sum_{i=1}^d v_i e_i$.  The vector Euclidean norm is denoted by $\|\cdot\|$. For a matrix $A \in \R^{k \times l}$, denote its $i, j$th entry by $A_{ij}$; its $i$th row by $A_{i.}$; and its $j$th column by $A_{.j}$. Where there is no ambiguity, we will use $A_i$ or $A_j$ to denote the $i$th row or $j$th column of $A$. 
%Let $\hat{\mu}_A$ to denote the empirical measure on $\mathbb{R}^k$ induced by the rows of $A$: $\{(A)_{1.},(A)_{2.},\cdots,(A)_{m.} \}$. 
%, a $p_1 \times p_2$ matrix . 
The spectral norm of $A$ is denoted by $\|A\|_2$, the Frobenius norm by $\|A\|_F$, the maximum Euclidean row norm by $\|A\|_{2 \to \infty}$, and the maximum absolute row sum by $\|A\|_{\infty}$. For a square matrix $A \in \R^{n \times n}$, $\lambda_1(A), \cdots, \lambda_n(A)$ denote the eigenvalues in decreasing order. We denote the set of matrices $W\in \RR^{n\times d}$ with orthonormal columns by $\mathcal{O}^{n\times d}$, and the set of permutation matrices $P\in \RR^{n\times n}$ by $\mathcal{P}_n$.
In our asymptotic analysis, we consider the asymptotic order of various quantities, so we introduce the order notation here. The relevant terminology and definitions below are reproduced directly from \cite{chen2024euclidean}.
\begin{definition}[Order Notation]
Let $\omega(m)$, $\alpha(m)$ be two quantities depending on $m$. We say that $\omega$ is of order $\alpha(m)$ and use the notation $\omega(m) \sim \Theta(\alpha(m))$ to mean that there exist positive constants such that for $m$ sufficiently large, $c\alpha(m)\leq \omega(m) \leq C \alpha(m)$. We write $\omega(m) \sim O(\alpha(m))$ if there exists a constant $C$ such that for sufficient large $m$ $\omega(m) \leq C \alpha(m)$.
\end{definition}

%Additionally we will consider the probabilities of certain events decay at specific rate, so we introduce the notion of convergence with high probability.

%\begin{definition}[convergence with high probability]
%Given a sequence of events $\left\{F_n\right\} \in \mathcal{F}$, where $n=1,2, \cdots$,  We say that $F_n$ occurs with high probability, and write $F_n$ w.h.p. , if for any $c_0>1$, there exists finite positive constant $C_0$ depending on $c_0$ such that $\operatorname{Pr}\left[F_n^c\right] \leq C_0 n^{-c_0}$ for all $n$. 
%We note that $F_n$ occurring w.h.p. is stronger than $F_n$ occurring asymptotically almost surely. Morever, $F_n$ occurring with high probability implies, by the Borel-Cantelli Lemma [23], that with probability 1 there exists an $n_0$ such that $F_n$ holds for all $n \geq n_0$.
%\end{definition}





\begin{definition}[Random dot product graph]
We say that the undirected random graph $G$ with adjacency matrix $\bA\in\RR^{n\times n}$ is a \emph{random dot product graph (RDPG)} with latent position matrix $\bX\in\RR^{n\times d}$, whose rows are the vectors $X^1,\ldots,X^n\in\mathcal{X}\subseteq\RR^d$, if
$$\PP[\bA|\bX]=\prod_{i<j}\langle X^i,X^j \rangle^{A_{i,j}}(1- \langle X^i,X^j \rangle)^{1-A_{i,j}},$$
where $\langle x,y\rangle= x^\top y$ is the inner product between two column vectors. We may denote this as $\bA\sim\mathrm{RDPG}(\bX)$. We call $\bP=\bX \bX^T$ the connection probability matrix. 
%In this case, each $A_{ij}$ is marginally distributed (conditionally on $X^i, X^j$) as Bernoulli($\langle X^i,X^j\rangle$) and independent of each other. 
If instead $\bP = \bX I_{p,q}\bX^\top$, where $I_{p,q}=I_p\oplus(-I_q)$ for some $p+q=d$, we call $\bA$ a \emph{generalized random dot product graph (GRDPG).}
\end{definition}

\begin{remark}[Orthogonal nonidentifiability in RDPGs] \label{rem:nonid}
Note that if $\bX \in \R^{n \times d}$ is a latent position matrix
and $\bW \in \R^{d \times d}$ is orthogonal,
$\bX$ and $\bX\bW$ give rise to the same distribution over graphs.
Thus, the RDPG model has a nonidentifiability up to orthogonal transformation. In the case of a GRDPG, $\bX$ and $\bX\bW$ give rise to the same distribution whenever $\bW I_{p,q} \bW^\top=I_{p,q},$ meaning that $\bW$ is an indefinite orthogonal transformation.
\end{remark}

\begin{definition}[Inner product distribution]
\label{def:innerprod}
Let $F$ be a probability distribution on $\R^d$. we call $F$ a \emph{$d$-dimensional inner product distribution}
if $0 \leq x^{\top} y \leq 1$ for all $x,y \in \supp F$. We call $F$ a \emph{$(p,q)$-dimensional generalized inner product distribution} if $0\leq x^\top I_{p,q} y \leq 1$ for all $x,y\in\supp F$, where $I_{p,q}=I_p\oplus (-I_q)$. 
%We will suppose throughout this work that for a $d$-dimensional inner product distribution $F$ and $X\sim F$, $\EE[XX^\top]$ has rank $d$.
\end{definition}

\begin{definition}[Latent position process]\label{def:latent_pos_proc}
A {\em latent position process} $\varphi(t)$ is a map $\varphi:[0,T]\rightarrow L^2(\Omega, \mathbb{P})$ such that for each $t\in [0,T]$, $\varphi(t)=X_t$, a random vector in $\R^d$ which has a (generalized) inner product distribution.
\end{definition}

%Once we have the latent position process, we can construct a time series of RDPGs whose vertices have independent, identically distributed latent positions given by $X_t$.

\begin{definition}[Latent position process network time series]
\label{def:tsg}
Let $\varphi$ be a latent position process, and fix a given number of vertices $n$ and collection of times $\mathcal{T}\subseteq[0,T]$. We draw an i.i.d.\ sample $\omega_j\in \Omega$ for $1\leq j\leq n$, and obtain the latent position matrices $\bX_{t}\in\RR^{n\times d}$ for $t\in\mathcal{T}$ by appending the rows $X_t(\omega_j)$, $1\leq j\leq n$. The \emph{time series of graphs (TSG)} $\{G_t: t\in\mathcal{T}\}$ are conditionally independent RDPGs with latent position matrices $\bX_{t}, t\in\mathcal{T}$.
\end{definition}

Existing literature on changepoint detection for networks has primarily focused on the marginal distributions of the networks. On the other hand, for a TSG where edges are assumed to arise conditionally independently (at each time and between times), and with a fixed small rank $d$ (or $p,q$ with $p+q=d$) for the mean matrices at each time, there is implicitly a (generalized) latent position process governing network evolution. As such, it is natural to define changepoints through properties of the LPP, rather than the distributions of the graphs themselves. This leads us to the notion of changepoints of different orders, below.

\begin{definition}
\label{def:highchangepoints}
A latent position process $\varphi(t)$ is said to have a \emph{zeroth-order changepoint} if there is some $t^*\in[0,T]$ such that
$$ X_t - X_{t'} \overset{\mathcal{L}}{=} \begin{cases}\delta_0 &\text{if }t,t'\leq t^*\text{ or }t^*\leq t,t'\\ F \neq \delta_0&\text{otherwise.}\end{cases}$$
Here $\delta_0$ is point mass at $0\in\RR^d$, and $F$ is some other distribution on $\RR^d$.

Given a partition of $[0,T]$ of mesh $\delta$, define the $k$th order difference of $\varphi$ at mesh $\delta>0$ recursively as follows: 
\begin{align*}
\Delta_1^\delta(t) &= X_t - X_{t-\delta},\\ 
\Delta_k^\delta(t) &= \Delta_{k-1}^\delta(t)-\Delta_{k-1}^\delta(t-\delta).
\end{align*}

A latent position process $\varphi(t)$ is said to have a \emph{$k$th-order changepoint of mesh }$\delta$ if there is some $t^*\in[0,T]$ such that for some distributions $F_1$, $F_2$, $F_3(t)$, $t^{*}<t\leq t^{*}+\delta$, not all equal, we have 
$$ \Delta_{k}^\delta(t) \overset{\mathcal{L}}{=} \begin{cases} F_1 &\text{if }t\leq t^*,\\
F_2&\text{if }t^*+\delta<  t,\\ F_3(t)&\text{if }t^*< t\leq t^*+\delta.\end{cases}$$ 
\end{definition}

Note that a first-order changepoint exactly corresponds to the setting in which 
$$\{\varphi(t): t\in [0,t^*)\}, \{\varphi(t): t\in(t^*+\delta,T]\}$$ are stationary stochastic processes (at least with respect to increments of length $\delta$), but with a loss of stationarity at the point $t^*$.

We emphasize in \ref{def:tsg} that each vertex in the TSG corresponds to a single $\omega\in\Omega$, which induces dependence between the latent positions for that vertex across times, but the latent position trajectories of any two distinct vertices are independent of one another across all times. Since these trajectories form an i.i.d.\ sample from the latent position process, it is natural to measure their evolution over time using the metric on the corresponding random variables, as described in \cite{athreya2025euclidean}. 
%We reproduce the maximum directional variation metric, $d_{MV}$, a metric on the space of square-integrable random variables, as follows.
\begin{definition}
The \emph{maximum directional variation metric} $d_{MV}:L^2(\RR^d)^2\rightarrow [0,+\infty)$ is given by
\begin{equation}
d_{MV}(X_t,X_{t'}):=\min_{W \in \mathcal{O}^{d \times d}} \left\|\EE[(X_t-WX_{t'})(X_t-WX_{t'})^\top]\right\|_2^{1/2},
\end{equation}
where $\mathcal{O}^{d\times d}$ is the set of all $d \times d$ orthogonal matrices and $\| \cdot \|_2$ is the spectral norm.
\end{definition}
%Given a map $\varphi:[0,T]\rightarrow L^2(\Omega)$ that assigns times $t$ to random variables $\varphi(t)=X_t$, we may write $d_{MV}(\varphi(t),\varphi(t')):=d_{MV}(X_t,X_{t'}).$

%In the definition of this distance, the expectation is over $\omega\in\Omega$, which means that it depends on the joint distribution of $X_t$ and $X_{t'}$. In particular
The $d_{MV}$ measure depends on more than just the marginal distributions of the random vectors $X_t$ and $X_{t'}$ individually; it also takes into account the dependence inherited from the latent position process $\varphi$. Because it depends on an expectation, it is a functional of the distribution of the latent position process. To estimate it from observed values of the latent position, we consider $\hat{d}_{MV}$, defined as follows.

\begin{definition}\label{hat-dmv}
Fix $t, s \in [0, T]$. Let $(X^{i}_t, X^{i}_s)$ for $1 \leq i \leq n$, be an observed pair of latent positions drawn independently and identically from a $d$-dimensional joint latent position distribution $F_{t,s}$. Let $\bX_t, \bX_s$ be the matrices of latent positions with rows $X^{i}_t$ and $X^i_s$, respectively. Define $\hat{d}_{MV}^2$ by
$$\hat{d}_{MV}^2(\bX_t,\bX_s)=\min_W \frac{1}{n} \left\|\sum_{i=1}^n (X^{i}_t-WX^{j}_s)(X^i_t-WX^j_s)^\top\right\|_2= \min_W\frac{1}{n}\left\|(\bX_t-\bX_s W)^\top(\bX_t-\bX_s W)\right\|_2,$$
where the minimization is over real orthogonal matrices $W\in\mathcal{O}^{d\times d}$.
\end{definition}

Next we consider the geometric properties of the image $\varphi([0,T])$ when equipped with the metric $d_{MV}$. Often, the map $\varphi$ admits a Euclidean analogue, called a {\em mirror}, which is a finite-dimensional curve that retains important signal from the generating LPP for the network time series. 

\begin{definition}[Exact Euclidean realizability with mirror $\psi$]
Let $\varphi$ be a latent position process. We say that the LPP $\varphi$ is \emph{exactly Euclidean $c$-realizable} with \emph{mirror} $\psi$ if there exists a Lipschitz continuous curve $\psi:[0,T]\rightarrow\RR^c$ such that:
$$
d_{MV}\left(\varphi(t), \varphi(t')\right):=d_{MV}\left(X_t, X_{t'}\right)=\|\psi(t)- \psi(t')\| \text{ for all } t,t' \in [0,T]. 
$$
\end{definition}

\begin{remark}
Exact Euclidean realizability implies the pairwise $d_{MV}$ distances between the latent position process at $t$ and $t'$ coincide exactly with Euclidean distances along the curve $\psi$ at $t$ and $t'$. This curve reflects, in Euclidean space, the $d_{MV}$ distances, and hence we call $\psi$ the {\em Euclidean mirror}.
% Sometimes we also use mirror to refer to the image of the curve, and when there are finite times, we refer mirror as the points representing the graphs in Euclidean space, which is also called as \emph{zero-skeleton mirror} in \cref{def:0-skeleton-dMV-mirror}.
\end{remark}

We similarly define approximate Euclidean realizability.
\begin{definition}[Approximate Euclidean realizability with mirror $\psi$]
Let $\varphi$ be a latent position process. For a fixed $\alpha \in (0, 1)$, we say that $\varphi$ is \emph{approximately \(\alpha\)-H\"older Euclidean \(c\)-realizable} if $\psi$ is \(\alpha\)-H\"older continuous, and there is some $C > 0$ such that
\[
\biggl| d_{MV}(\varphi(t), \varphi(t')) - \|\psi(t) - \psi(t')\| \biggr| \leq C|t - t'|^\alpha \quad \text{for all } t, t' \in [0, T].
\]
\end{definition}

If the LPP is exactly $c$-Euclidean realizable with mirror $\psi$, and we sample $m$ time points
$$\mathcal{T}=\{t_1,t_2,\cdots,t_m\} \subseteq [0,T],$$ then the corresponding distance (or dissimilarity) matrix $$\left(\mathcal{D}_\varphi\right)_{i,j}:=d_{MV}(X_{t_i},X_{t_j})$$ is an \emph{exactly $c$-Euclidean realizable distance (dissimilarity) matrix}: that is, there exist $m$ points $$\psi(t_1),\psi(t_2),\cdots,\psi(t_m) \in \mathbb{R}^c$$ such that $$\left(\mathcal{D}_\varphi\right)_{i,j}=\|\psi(t_i)-\psi(t_j)\| \text{~for all~} i,j \in \{1,2,\cdots,m\}.$$ 

In this case, applying classical multidimensional scaling (CMDS) to $\mathcal{D}_{\varphi}$ will recover the mirror $\psi(t_1),\psi(t_2),\cdots,\psi(t_m)$ exactly up to an orthogonal transformation. Recall that the CMDS embedding is defined as follows.

\begin{definition}[CMDS embedding to dimension $c$]
\label{def:CMDS}
Let $\mathcal{D} \in \mathbb{R}^{m \times m}$ be a distance matrix and define the centering matrix $P$ by $P:=I_m-\frac{J_m}{m}$ where $I_m$ is the $m \times m$ identity matrix and $J_m$ is the $m \times m$ all ones matrix. Consider $B:=-\frac{1}{2}P\mathcal{D}^{(2)}P$ where $\mathcal{D} ^{(2)}$ is the entrywise square of the distance matrix $\mathcal{D} $. Denote the $c$ largest eigenvalues and corresponding orthogonal eigenvectors of $B$ as $\lambda_1,\cdots,\lambda_c$ and $u_1,\cdots,u_c$. Set $\max(\lambda_i,0)=\sigma^2_i$ and $S$ a diagonal matrix with $S_{ii}:=\sigma^2_i$ for $i \in \{1,\dots,c\}$; let $U \in \R^{m \times c} := [ u_1 | \cdots | u_c ]$. The classical multidimensional scaling of $\mathcal{D}$ into dimension $c$, where $1\leq c\leq m-1$, is the set of $m$ points $\psi(t_1),...,\psi(t_m) \in \mathbb{R}^c$ defined by the rows of $\Psi$ as 
%For each $1 \leq j \leq c$, let $\tilde{\psi}_j$ denote the column vector in $\mathbb{R}^m$ defined by $\tilde{\psi}_j=[\psi_j(t_1) \cdots \psi_j(t_m)]^{\top}$ We have
$$
\Psi:=[\psi_{j}(t_i)]_{i=1,j=1}^{m,c}
=\left[
\begin{array}{ccc}
\psi_1(t_1) & \cdots & \psi_c(t_1) \\ 
\vdots &  & \vdots \\ 
\psi_1(t_m) & \cdots & \psi_c(t_m)
\end{array}\right]=
%\left[
%\tilde{\psi}_1 | \dots | \tilde{\psi}_c \right]=
 \left[
 \begin{array}{c}
  \psi^{\top}(t_1) \\ 
 \vdots \\ 
 \psi^{\top}(t_m) 
 \end{array}
 \right]=
US^{\frac{1}{2}}=
\left[ \sigma_1 u_1|\cdots| \sigma_c u_c\right].
$$
\end{definition}


\begin{remark}
If there are no positive eigenvalues for \( B \), then \( S=0 \) and thus \( \Psi=0 \). This implies that the distance matrix has no meaningful Euclidean representation. We note that \( B \) must be positive semidefinite for \( \mathcal{D} \) to be exactly Euclidean realizable.
\end{remark}
Embedding a distance or dissimilarity matrix through CMDS provides a constellation of points in Euclidean space whose interpoint Euclidean distances approximate those in the dissimilarity matrix. In the case of a latent position process, this embedding can provide key signal about the underlying process itself. We call this the {\em zero-skeleton} mirror. 
\begin{definition}[Zero-skeleton mirror]\label{def:0-skeleton-dMV-mirror} Suppose $\mathcal{D}_{\varphi}$ is the $d_{MV}$ distance matrix associated with an  LPP $\varphi$. The output of CMDS applied to $\mathcal{D}_{\varphi}$ produces $m$ points in $\R^c$, denoted $\{\psi(t_1),...,\psi(t_m)\}$, called the \emph{zero-skeleton $c$-dimensional ($d_{MV}$) mirror}. 
\end{definition}



% \begin{remark}
% Throughout the paper, $\psi(\cdot):\mathbb[0,T] \rightarrow \R^c$, where $c$ is the chosen CMDS embedding dimension, denotes the $c$-dimensional zero-skeleton Euclidean mirror, which we call the ``mirror" for simplicity. The image of the function $\psi$ is a $c$ dimensional curve, and $\psi(t) \in \mathbb{R}^c$ provides representation of the graph at time $t$ in $c$ dimensional Euclidean space. 
% % Meanwhile, it is also the image of the function $\psi: \R \to \R^c$ at $t$ and the function $\psi$ describes a curve lives in $c$ dimensional Euclidean space. 
% When $c>1$, we use $\psi_k:\R \to \R$ to denote the $k$th component of $\psi$ and $[\psi_k(t_i)]^m_{i=1} =\sigma_k u_k$.
% When we consider $m$ time points, we stack $\{\psi^{\top}(t_1),\psi^{\top}(t_2),\cdots,\psi^{\top}(t_m)\}$ row-wise, as a matrix, denoted by $\Psi \in \R^{m \times c}$.
% \end{remark}



Our network time series consists of time-indexed graphs on a common vertex set, and for any given time, each vertex has a corresponding time-dependent latent position. The latent positions are typically unknown, but can be consistently estimated through spectral decompositions of the observed adjacency matrices \cite{STFP-2011}.  Suppose that $G$ is a random dot product graph with latent position matrix $\bX \in \R^{n \times d}$, where the rows of $\bX$ are independent, identically distributed draws from a latent position distribution $F$ on $\mathbb{R}^d$. Let $\bP=\bX \bX^T$ be the connection probability matrix and $\bA$ be the adjacency matrix for this graph. The {\em adjacency spectral embedding} (ASE) is a rank $d$ eigendecomposition of the adjacency matrix.

\begin{definition}[Adjacency Spectral Embedding]
Given an adjacency matrix $\bA$, we define the \emph{adjacency spectral embedding} (ASE) with dimension $d$ as $\hat{\bX}=\hat{\bU}|\hat{\bS}|^{1/2}$, where $\hat{\bS}\in\R^{d\times d}$ is the diagonal matrix with the $d$ largest-magnitude eigenvalues of $\bA$ on its diagonal, arranged in decreasing order. $\hat{\bU}\in\R^{n\times d}$ is the matrix of $d$ corresponding orthonormal eigenvectors arranged in the same order. $\hat{\bX}$ is the estimated latent position matrix.  
\end{definition}
% A central consideration in this work is the effect of vertex permutations. When the vertices of a graph with adjacency matrix $\bA$ are shuffled by a permutation matrix $P$, the new adjacency matrix is given by the similarity transformation $P\bA P^{\top}$. A key property of Adjacency Spectral Embedding (ASE) is that it is equivariant with respect to vertex permutation group action.
% \begin{lemma}[ASE is permutation equivariant.] \label{lem:ASE_equivariant}
% Denote $\hat{\bX}_{\bA}$ as the ASE of $\bA$ with $d$ dimension. Then for any undirected adjacency matrix $\bA$ with $n$ vertices any $d<n$ and any $n \times n$ permutation matrix $P$, 
% $$
% \hat{\bX}_{P\bA P^{\top}} = P \hat{\bX}_{\bA} .
% $$ 
% \end{lemma}


%Under some regularity conditions on RDPG(X), as shown in \cite{athreya2017statistical}, with high probability, $\bA$ has top $d$ eigenvalues nonnegative. Thus ASE is well-defined. 
When the true dimension $d$ of the latent positions---or equivalently, the rank of $\bP$---is unknown, we can infer it using the scree plot \cite{zhu2006automatic}. The ASEs of the observed adjacency matrices in our TSG  at times $t$ and $s$, denoted $\hat{\bX}_t$ and $\hat{\bX}_s$, are estimates of the latent position matrices $\bX_t$ and $\bX_s$, from which we obtain the estimate $\hat{d}_{MV}$ of the $d_{MV}$ distance between the latent position random variables over time by applying Definition~\ref{hat-dmv} to these estimated matrices of latent positions.

%\begin{enumerate}
%    \item LPP (done)
%    \item TSG (done)
 %   \item CMDS (done)
%    \item mirror (done)
%    \item estimated dMV (done)
%\end{enumerate}






% \begin{definition}
% \label{def:highchangepoints}
% A latent position process $\varphi(t)$ is said to have a \emph{zeroth-order changepoint} if there is some $t^*\in[0,T]$ such that
% $$ X_t - X_{t'} \overset{\mathcal{L}}{=} \begin{cases}\delta_0 &\text{if }t,t'\leq t^*\text{ or }t^*\leq t,t'\\ F \neq \delta_0&\text{otherwise.}\end{cases}$$
% Here $\delta_0$ is point mass at $0\in\RR^d$, and $F$ is some other distribution on $\RR^d$.

% Given a partition of $[0,T]$ of mesh $\delta$, define the $k$th order difference of $\varphi$ at mesh $\delta>0$ recursively as follows: 
% \begin{align*}
% \Delta_1^\delta(t) &= X_t - X_{t-\delta},\\ 
% \Delta_k^\delta(t) &= \Delta_{k-1}^\delta(t)-\Delta_{k-1}^\delta(t-\delta).
% \end{align*}

% A latent position process $\varphi(t)$ is said to have a \emph{$k$th-order changepoint of mesh }$\delta$ if there is some $t^*\in[0,T]$ such that for some distributions $F_1$, $F_2$, $F_3(t)$, $t^{*}<t\leq t^{*}+\delta$, not all equal, we have 
% $$ \Delta_{k}^\delta(t) \overset{\mathcal{L}}{=} \begin{cases} F_1 &\text{if }t\leq t^*,\\
% F_2&\text{if }t^*+\delta<  t,\\ F_3(t)&\text{if }t^*< t\leq t^*+\delta.\end{cases}$$ 
% \end{definition}




% We recall the $d_{MV}$ metric defined in \cite{athreya2025euclidean}. 
% \begin{definition}
% The \emph{maximum directional variation metric} $d_{MV}:L^2(\RR^d)^2\rightarrow [0,+\infty)$ is given by
% \begin{equation}
% d_{MV}(X_t,X_{t'}):=\min_{W \in \mathcal{O}^{d \times d}} \left\|\EE[(X_t-WX_{t'})(X_t-WX_{t'})^\top]\right\|_2^{1/2},
% \end{equation}
% where $\mathcal{O}^{d\times d}$ is the set of all $d \times d$ orthogonal matrices and $\| \cdot \|_2$ is the spectral norm.
% \end{definition}


% \begin{definition}\label{def:hat_dmv}
% The {\em estimated} $d_{MV}$ distance $\hat{d}_{MV}$ is defined as 
% \begin{equation}\label{equ:hat_dmv}
% \hat{d}_{MV}(\hat{\bX}_{t},\hat{\bX}_{s}):=\min_{W\in\mathcal{O}^{d\times d}} \frac{1}{\sqrt{n}}\|\hat{\bX}_{t}-\hat{\bX}_{s}W\|_2.
% \end{equation}
% \end{definition}

% \begin{definition}\label{def:shuffling}
% [$\alpha$-shuffling, Shuffling, and Derangement]
% For a positive integer \( n \) with \( n \ge 2 \), a \emph{permutation} \(\sigma: [n] \to [n]\) is a one-to-one map from \([n]\) to \([n]\). The set of all permutations is denoted by \( S_n \), and \( |S_n| = n! \). To every \(\sigma \in S_n\), there corresponds a permutation matrix \( P_{\sigma} \in \mathbb{R}^{n \times n} \). A permutation is called a \emph{derangement} if for all \( i \in [n] \), \( \sigma(i) \neq i \). Denote the number of fixed points of any permutation $\sigma_0$ by $N_{\text{fix}}^{\sigma_0}$. Denote by $\cdot^{-1}:S_n \to S_n$ the map such that for any $\sigma_0 \in S_n$, $\cdot^{-1}(\sigma_0) =\sigma_0^{-1}$. Define a composite map of two permutations $\circ:S_n \times S_n \to S_n$ as follows: for any $(\sigma^1_0,\sigma^2_0) \in S_n \times S_n$, $ \circ(\sigma_0^1 , \sigma^2_0) =\sigma_0^1(\sigma^2_0)$.
% \emph{shuffling} is a random permutation drawn from \( S_n \) according to the uniform distribution. That is, $\sigma$ is a shuffling if $\PP(\sigma = \pi) = \frac{1}{n!}$ for any $\pi \in S_n$.
% \textbf{We denote a shuffling on $S_n$'s corresponding permutation matrix as $P_n$.}
%Formally, consider the probability space $(S_n, 2^{S_n} , \PP_{S_n} )$ where $\PP_{S_n}$ is the uniform measure on $S_n$, namely, for any $\sigma_0 \in S_n$, $\PP_{S_n}(\sigma_0) = \frac{1}{n!}$. A shuffling $\sigma$ is any random permutation
%$\sigma$ defined on $(S_n, 2^{S_n} , \PP_{S_n} )$ that follows a uniform distribution $\PP_{S_n}$ on $S_n$.
%TC: i don't care it maps from which space to S_n and its measure on the original space, as long as its push forward measure on Sn is the uniform distirbution, it is a shuffling. 
%, where $(\mathcal{X}, \mathcal{F}, \mu)$ is an arbitrary probability space, but we will only use $\mathcal{X} = S_n$ or $\mathcal{X} = S^2_n$.
% For a given $0\le\alpha \le 1$ and $n$, we also define \emph{fixed partial shuffling with percentage $\alpha$}, or in short $\alpha-$shuffling.
% Denote $\alpha_n = \lfloor \alpha n \rfloor$.
% $\sigma_{\alpha_n}$ is a shuffling from $S_{\alpha_n}$ and \textbf{its corresponding permutation matrix is $P_{n, \alpha}$.} 
% Then $P_{n, \alpha} := I_{n-\alpha_n} \oplus P_{\alpha_n}$ is the matrix corresponding to the so-called $\alpha$-shuffling. That is, an $\alpha$-shuffling only shuffles the last $\alpha_n$ indices in $[n]$. 
% \end{definition}



% \begin{lemma}[Properties of shuffling]\label{lem:prop of shuffling}
% \begin{itemize}

%     \item If $\sigma$ is a shuffling, then $\EE[N_{\text{fix}}^\sigma] = 1$.

%     \item If \(\sigma\) is a shuffling, then its inverse \(\sigma^{-1}:=\cdot^{-1}(\sigma)\) is also uniformly distributed on \( S_n \), namely a shuffling.

%     \item Suppose $\sigma_0 \in S_n$ is a permutation, and $\sigma$ is a shuffling, then $\sigma_0 \circ \sigma$ and  $\sigma \circ \sigma_0$ are both shufflings.
    
%     \item Suppose \(\sigma\) and \(\sigma'\) are two independent shufflings. 
    %Consider the product probability space $(S_n \times S_n, 2^{S_n \times S_n}, \PP_{S_n \times S_n})$ with $\PP_{S_n \times S_n}$ as the product measure of $\PP_{S_n} \otimes \PP_{S_n}$.
%     Then  $\sigma \circ\sigma' := \circ(\sigma,\sigma')$ is also a shuffling.  

%     \item If \(\sigma_1\) and \(\sigma_2\) are two independent $\alpha-$shufflings, then $\sigma_1\circ \sigma_2$ is an $\alpha-$shufflings.
    
% \end{itemize}
% \end{lemma}

%Also we can define \emph{random partial shuffling with percentage $\alpha$}, that is we firstly draw an element from $[n]$ uniformly without replacement $\alpha_n$ times, that is we get a subset of $[n]$ with $\alpha_n$ elements. Then shuffle the drawn set. Namely for the drawn set $ A = \{a_1,...,a_{\alpha_n}\}$, draw a shuffling $\sigma_{\alpha_n}$ from $S_{\alpha_n}$, then the random partial shuffling with percentage $\alpha$ $\sigma$ satisfy that: when $x \in A$, assume $x = a_i$ then  $\sigma(x) = a_{\sigma_{\alpha_n}(i)}$;when $x \notin A$ $\sigma(x) = x$.

%\item If \(\sigma_1\) and \(\sigma_2\) are two independent random $\alpha-$shufflings, then $\sigma_1\circ \sigma_2$ has expected unaffected size $n(2\alpha-\alpha^2)$.

\section{Vertex misalignment, graph matching, and optimal transport}\label{sec:Main_results_GM_OT}
In this section, we consider consequences of vertex misalignment on latent position process networks, specifically for correlation and dissimilarity measures, and we explore the relationship between graph matching and optimal transport, which are two approaches for mitigating vertex misalignment.
%second, to describe two different edge-case LPPTSG models, London and Atlanta, both of which undergo a changepoint; third, to examine how different statistics can be used for changepoint localization for these processes; and fourth, to see how vertex misalignment impacts this localization. 

\subsection{Vertex shuffling, dependence, and dissimilarity}\label{sec:D^3}

The $d_{MV}$ distance captures dependence between latent positions across time by computing an expectation with respect to their joint distribution. To estimate this from data therefore requires the vertex correspondence to be known, and when it is not, the measured dependence between the random variables is weakened, leading to the following notion of distance:

% ordering of the rows of the adjacency matrices can be partially or completely uninformative, and the correspondence between rows of the true and estimated . The $d_{MV}$ distanceSuch a shuffling also impacts the correlation between the true latent position at time $t$ and time $s$The impact of vertex shuffling on a network is represented as similarity transformation by a uniformly-distributed permutation matrix, by which the matrix $A_t$  is mapped to  $PA_t P^\top$. Since the eigendecomposition is permutation-equivariant, a shuffling of the vertices between times $t$ and $s$ on the $\hat{d}_{MV}$ distance transforms 
% $\hat{d}_{MV}(\hat{\bX}_t,\hat{\bX}_s)$ into $\hat{d}_{MV}(\hat{\bX}_t,P\hat{\bX}_s).
% $
%Theorem~\ref{thm:shuffleddMV_convergence} shows that $\hat{d_{MV}(\$ converges in probability to a notion of distance on latent position random variables called shuffled-$d_{MV}$, where we replace the expectation in the definition of $d_{MV}$, which is defined with respect to the joint measure of the random variables, with an expectation with respect to the product of the marginal measures, essentially treating the variables as independent. This gives theoretical justification for our interpreting vertex misalignment as weakening the temporal dependence of latent positions.

% Assume we have observations of a latent position process time series of graphs (LPPTSG) on $n$ correctly aligned vertices at $m$ distinct time points. Denote the realized adjacency matrices of this time series by
% $
% \bA_1,\bA_2,....,\bA_m.
% $
% %To study or represent the dynamics of the TSG, the simplest statistics is to consider graph statistics change over time, for example, average degree of each graphs. 
% %Further, to utilize the low rankness followed from assuming graphs from RDPG, we can ASE each graph and obtain
% %$$
% %\{\hat{\bX}_{1}, \hat{\bX}_{2} , ... , \hat{\bX}_{m}  \}.
% %$$ 
% %Most methods, such as Euclidean mirror will capture the dynamics using the above ASE results, they are essentially a functional or statistics based on  above ASE results. 
% Let $\{\sigma_1,\sigma_2,...\sigma_m\}$, with corresponding permutation matrices be drawn independently and uniformly from $S_n$ with corresponding permutation matrices
% $\{P_1,P_2,...,P_m\}$. The adjacency matrices shuffled by these permutations may be written as $P_1 \bA_1 P_1^{\top} ,  P_2 \bA_2 P_2^\top , ...,  P_m \bA_m P^\top_m$, and their corresponding ASE estimates are given by $P_t \hat{\bX}_t$. 



% And by Lemma \ref{lem:ASE_equivariant} the corresponding ASE results will be row wise shuffle the estimated latent position matrices:
% $
% P_1\hat{\bX}_{1}, P_{2}\hat{\bX}_{2} , ... , P_{m}\hat{\bX}_{m} $.
% By \cite{athreya_survey} as $n$ grows, the above will converge up to a orthogonal transformation to the true latent position matrices:
% $
% P_1\bX_{1}, P_2\bX_{2} , ... , P_m\bX_{m}.
% $
% Thus for simplicity, we consider the case when the true latent positions matrices are known.
% %and the dynamic representation as a function on shuffled true latent position matrices,
% %Recall as assumed in Definition \ref{def:tsg} the $i$th row of those $m$ matrices forms an i.i.d sample path drawn from the LPP. Namely 
% %$$\biggl(  
% %\left(\bX_1\right)_{i}, \left(\bX_2\right)_{i},..., \left(\bX_m\right)_{i} \biggr) \overset{\mathcal{L}}{=} (X_1,X_2,...,X_m), ~~ \forall ~ i \in[n].
% %$$
% %Thus by well-known Glivenko–Cantelli theorem, the empirical distribution formed by i.i.d samples  converges to joint distribution of $(X_1,X_2,...,X_m)$. Thus any subsequent inference statistics as a function of $\{\bX_{1}, \bX_{2} , ... , \bX_{m}  \}$ can be studied using function on $(X_1,X_2,...,X_m)$.
% Equivalently, $P_1\bX_{1}, P_2\bX_{2} , ... , P_m\bX_{m}$ can be seen as the shuffled samples:
% $\biggl\{
% \biggl(  
% \left(\bX_1\right)_{\sigma^{-1}_1(i)}, \left(\bX_2\right)_{\sigma^{-1}_2(i)},..., \left(\bX_m\right)_{\sigma^{-1}_m(i)} \biggr) \biggr\}^n_{i=1}.
% $
% That is, for vertex $i$, at time $1$, instead of observing $(\bX_1)_i$, we observe $(P_1\bX_1)_i = (\bX_1)_{\sigma^{-1}_1(i)}$ --- the realized latent position of the other vertex $\sigma^{-1}_1(i)$. 
% %it is a vertex misalignment, then what is the joint distribution of such shuffled samples. 


% The next step in mirror method is to compute $\hat{d}_{MV}$ on $P_t\mathbf{X}_t$ and $P_s\mathbf{X}_s$ for any $t, s \in [m]$.
% This quantity is a function of the sample formed by pairing the rows of the two permuted matrices.
% Let the rows of the true latent positions $\mathbf{X}_t$ and $\mathbf{X}_s$ be $\{x_i\}_{i=1}^n$ and $\{y_i\}_{i=1}^n$, respectively.
% Pairing the $i$-th row of $P_t\mathbf{X}_t$ (which is $x_{\sigma_t^{-1}(i)}$) with the $i$-th row of $P_s\mathbf{X}_s$ (which is $y_{\sigma_s^{-1}(i)}$) gives the sample:
% $
% \left\{ \left( x_{\sigma_t^{-1}(i)}, y_{\sigma_s^{-1}(i)} \right) \right\}_{i=1}^n.
% $
% By re-indexing we see that it is equivalent in distribution to a sample of the form $\{(x_j, y_{\sigma(j)}) \}_{j=1}^n$, where $\sigma = \sigma_s^{-1} \circ \sigma_t$ is also a uniform random permutation by Lemma \ref{lem:prop of shuffling}.

% This structure motivates a more fundamental statistical question.
% Assume we have i.i.d. samples $\{ (x_i,y_i) : i\in [n] \}$ from a joint distribution $\mu_{X,Y}$. If $\sigma$ is a random permutation, what is the empirical distribution of the shuffled sample $\{ (x_i,y_{\sigma(i)}) : i\in [n] \}$?
% \begin{comment}
%     Next note for mirror method, we only consider $\hat{d}_{MV}$ on the pair of time points, which is based on $$\biggl\{
% \biggl(  
% \left(\bX_t\right)_{\sigma_t(i)}, \left(\bX_s\right)_{\sigma_s(i)} \biggr) \biggr\}^n_{i=1},
% $$
% whose empirical distribution is the same as 
% $$\biggl\{
% \biggl(  
% \left(\bX_t\right)_i, \left(\bX_s\right)_{\sigma_t^{-1} \circ\sigma_s(i)} \biggr) \biggr\}^n_{i=1},
% $$
% by Lemma \ref{lem:prop of shuffling}, $\sigma^{-1}_t \circ \sigma_s$ is still a shuffling. Thus the question of interest amounts to consider i.i.d samples $\{ (x_i,y_i) : i\in [n] \}$ from a joint distribution $\mu_{X,Y}$, if $\sigma$ is a shuffling, what is the empirical distribution of $\{ (x_i,y_{\sigma(i)}) : i\in [n] \}$? 
% \end{comment}
% Intuitively, applying a shuffling breaks the correlation and weaken the dependence while preserving the marginal distributions.
% This concept is central to permutation-based tests of independence \cite{albert2015tests}.
% To see it, observe that because $(x_i,y_i)$ are independent to each other,  $x_i \ind y_{\sigma(i)}$ if $i \neq\sigma(i)$, thus \( (x_i, y_{\sigma(i)}) \sim   \mu_X \otimes \mu_Y\). 
% %namely, the shuffled pair distributed as if it is drawn from $X$, $Y$ with original marginal but now independent.
% The theorem \ref{thm:shuffling-2d} formalizes this intuition by showing that the applying a test function on shuffling samples is ``close" to its expected value with product measure \(\mu_X\otimes \mu_Y\) with high probability for large $n$.
% Recall a formal definition of product measure: 
% \begin{definition}[product measure]
% Let $(\Omega_1, \mathcal{F}_1, \PP_1)$ and $(\Omega_2, \mathcal{F}_2, \PP_2)$ be two probability spaces. The product measure $\PP_{\Omega_1 \times \Omega_2}$ is a probability measure on the product space $(\Omega_1 \times \Omega_2, \mathcal{F}_1 \otimes \mathcal{F}_2)$ defined by the property that for all measurable sets $A_1 \in \mathcal{F}_1$ and $A_2 \in \mathcal{F}_2$,
% $$
% (\PP_{\Omega_1 \times \Omega_2})(A_1 \times A_2) = \PP_1(A_1) \PP_2(A_2).
% $$
% Here, $\mathcal{F}_1 \otimes \mathcal{F}_2$ is the $\sigma$-algebra on $\Omega_1 \times \Omega_2$ generated by subsets of the form $A_1 \times A_2$, with $\times$ being Cartesian product.

% Because probability spaces are always $\sigma$ finite, then the product measure $\PP_{\Omega_1 \times \Omega_2}$ is unique. For any measurable set $E \in \mathcal{F}_1 \otimes \mathcal{F}_2$, its measure is given by
% $$
% (\PP_{\Omega_1 \times \Omega_2})(E) = \int_{\Omega_1} \PP_2(E_x) \, d \PP_1(x) = \int_{\Omega_2} \PP_1(E_y) \, d \PP_2(y),
% $$
% where $E_x = \{y \in \Omega_2 : (x,y) \in E\}$ is the $x$-section of $E$, and $E_y = \{x \in \Omega_1 : (x,y) \in E\}$ is the $y$-section of $E$.
% \end{definition}


% %maybe we can prove joint distribution of $(x_i, y_{\sigma(i)})$ is still exchangeable and then use exchangeable Hoeffding to prove a better bound. 

% \begin{thm}[Shuffling weakens dependence]\label{thm:shuffling-2d}
% Consider two random variables $X,Y\in\RR^1$ with joint distribution $\mu_{X,Y}$ on $\RR^2$. 
% %defined on a probability space $(\RR^2, \mathcal{B}(\RR^2))$.
% For any $n>2$, denote $\{(x_i,y_i)^n_{i=1}\}$ as $n$ i.i.d samples from $\mu_{X,Y}$ defined in the probability space 
% %$((\RR^2)^{\times n}, (\mathcal{B}(\RR^2))^{\times n}) = : 
% $(\Omega, \mathcal{B} ).$ 
% Denote $\sigma$ as a shuffling on $S_n$.
% %and $\PP_{S_n}$ as the uniform distribution on $S_n$.
% Assume $\sigma$ is independent of the sample $\{(x_i,y_i)^n_{i=1}\}$.
% Then $\{ \left(x_i, y _{\sigma(i)} \right); i \in [n]\}$ is the shuffled sample. 
% %Assume the product space of $\Omega$ and $S_n$ as $S_n \times \Omega$ and the product measure $\PP_{\Omega} \otimes \PP_{S_n}$ on it is denoted as $\PP_{\Omega\times S_n}$. 
% For a function $g: \RR^d \to \RR$, denote $\EE_{\mu}[g] : =\EE_{(x_1,x_2,\cdots,x_d) \sim \mu}[g(x_1,x_2,\cdots,x_d)]$.
% Then for every function $f: \R^2 \to \R^1$ bounded by $B_{f, \mathcal{S}}$ on the support of $\mathcal{S}: = \text{support of } \mu_X \otimes \mu_Y $, and any $\epsilon > 0 $,
% $$
% \PP_{\Omega \times S_n} \left( \left| \frac{1}{n} \sum^n_{i=1} f\left(x_i, y_{\sigma(i)}\right)  -  \EE_{ \mu_X \otimes \mu_Y}[f]   \right| > \epsilon \right) \leq \frac{ 12B_{f,\mathcal{S}}^2}{n \epsilon ^2} + \frac{4B_{f,\mathcal{S}}}{ n\epsilon }.
% $$
% \end{thm}

% \begin{remark}
% %$\sigma$ as a prefixed constant derangement is independent of the sample $\{(x_i,y_i); i \in [n]\}$. 
% The independence assumption is essential because it guarantees that if $\sigma(i) \neq i $, $x_i \ind y_{\sigma(i)}$. For example, consider an ``adversarial" permutation $\sigma_a$ with property that for all $n$, for any $i$ with $\sigma(i) \neq i$, we still have $y_{\sigma(i)} = y_i$ (this can happen when $\mu_Y$ is supported on a discrete set). Then clearly as $n$ increases $\frac{1}{n} \sum^n_{i=1} f\left(x_i, y_{\sigma(i)}\right)$ will converge to $\EE_{\mu_{X,Y}}[f]$ instead of $\EE_{\mu_{X}\otimes\mu_{Y}}[f]$. 
% But clearly in this case $\sigma$ is not independent of the samples since $\sigma$ is a function of $\{y_i\}^n_{i=1}$.
% %Meanwhile if derangement $\sigma$ is fixed and we happen to have a sample so that $\forall i, y_{\sigma(i)} = y_i$. This doesn't violate the theorem because although in this case $ \EE_{\hat{\mu}_n}[f] -  \EE_{ \mu_X \otimes \mu_Y}[f]$ might be big, but the probability of this happening is exponentially decreasing with $n$, which is allowed by the theorem.

% %We should be able to extend this result to the case when $\sigma$ is independent of   $\{(x_i,y_i); i \in [n]\}$ and is drawn randomly from the set of all possible permutations. 
% \end{remark}


% Theorem \ref{thm:shuffling-2d} shows for every $f$ bounded on the support, as $n$ grows, the empirical mean of $f$ on shuffled samples can be arbitrarily close to the expected value of $f$ on the product measure with high probability.
% From a statistics perspective, this means after shuffling, the statistics $f$ applied on the shuffled samples drawn from $\mu_{X,Y}$, behave as if the samples were from $\mu_X \otimes \mu_Y$. 
% %That is, the information that was in the joint distribution and can be captured by $f$ is now washed out by shuffling.
% This answers the question: what is lost if we loose the true correspondence among samples?
% The answer is we loose the information in the joint distribution $\mu_{X,Y}$ and only the information remained in the marginal distribution $\mu_{X}\otimes\mu_{Y}$ gets preserved.
% Thus if the parameter of interest is in the joint distribution and not the marginal distribution, then misalignment will undermine the inference of the parameter.
% On the other hand when the marginal distribution is already informative for the parameter, then the inference procedure is robust to misalignment. 
%Further we can extend the result from derangement to shuffling $\sigma$. 
%Although now there could be several fixed points in $\sigma$, we will show they are negligible and the above result still holds. 
%Theorem \ref{thm:derangement-2d} and \ref{thm:shuffling-2d} illustrates that samples with derangement or shuffling $\{(x_i,  y_{\sigma(i)})\}$ behave like samples from $X$, $Y$ with original marginal distribution but jointly distributed as if they are independently. 
%\textcolor{red}{to understand theorem 6's implications we need to consider shuffling in the context of our inference task...Shuff DMV is next then(DONE!)}
%Specifically, we consider the lost of vertex correspondence in the context of time series of graphs.
%Next, to estimate the mirror for representation of dynamics of the TSG, we calculate $\hat{d}_{MV}(\hat{\bX}_t, \hat{\bX}_s)$ for all pairs $t, s \in [m]$ and construct a $d_{MV}$ distance matrix. In the vertex misalignment context, we are actually considering $\hat{d}_{MV}(P_t\hat{\bX}_t, P_s\hat{\bX}_s)$. 
%Thus when applying the mirror method, the impact of vertex misalignment can be quantifies as difference between $\hat{d}_{MV}(\hat{\bX}_t, \hat{\bX}_s)$ and $\hat{d}_{MV}(P_t\hat{\bX}_t, P_{s}\hat{\bX}_s)$. Note 
%$$\hat{d}^2_{MV}(P_t\hat{\bX}_t, P_s\hat{\bX}_s) = \min_{W \in \mathcal{O}^{d \times d}} \|P_t\hat{\bX}_t -  P_s\hat{\bX}_sW \|_2 = \min_{W \in \mathcal{O}^{d \times d}} \|\hat{\bX}_t -  P^{-1}_tP_s\hat{\bX}_sW \|_2.$$ 
%And by Lemma \ref{lem:prop of shuffling}, $ P^{-1}_tP_s$ is still a shuffling. Thus it is enough to consider the quantity $\hat{d}_{MV} ( \bX_t, P\bX_s )$ where $P$ is shuffling drawn from $S_n$. 
%To further simplify, because \cite{athreya2025euclidean} has shown that $\hat{d}_{MV}(\hat{\bX}_t, P\hat{\bX}_s) \to \hat{d}_{MV}(\bX_t, P\bX_s)$ with high probability as $n$ grows, we only focus on $\hat{d}_{MV}(\bX_t, P\bX_s)$, where $\bX_{t}$ is the true latent position matrix with rows as $n$ i.i.d samples drawn from $X_{t}$.
%
% Now we consider the impact of vertex misalignment specifically for Euclidean mirror.
% Recall mirror is the CMDS result of $\hat{d}_{MV}$ distance matrix whose $t,s$-th entry in the misalignment case is $\hat{d}_{MV}(\hat{\bX}_t, P_n\hat{\bX}_s)$ with $P_n$ as a shuffling on $S_n$. 
% And this quantity converges to $\hat{d}_{MV}(\bX_t, P\bX_s)$ because of consistency of ASE.  
% To analyze this quantity, we focus on a specific class of LPPs.
% \begin{definition}[Positive scalar LPP]\label{def:simple_class_LPP}
% A LPP $\{X_t ; t \in [m]\}$ is a positive scalar LPP if for all $t \in [m]$, $X_t$ is a 1 dimension random variable and is supported on positive real numbers. 
% \end{definition}

% For a positive, scalar-valued LPP, $d_{MV}$ distance is always minimized by choosing othogonal $W=1$.
%  $d_{MV}$ and $\hat{d}_{MV}$ simplifies to:
% $$
% d^2_{MV} \left(X_t, X_{t'}\right) = \EE [(X_t-X_{t'})^2], ~~~~
% \hat{d}^2_{MV} \left(\bX_t, \bX_{t'}\right) = \frac{1}{n}\| \bX_t -  \bX_{t'}\|^2_F. 
% $$   

% Applying this simplified form to the misaligned case yields $\hat{d}^2_{MV} \left(\bX_t, P\bX_{s}\right) = \frac{1}{n}\| \bX_t -  P\bX_{s}\|^2_F = \frac{1}{n} \sum^n_{i=1} ((\bX_t)_ i - (\bX_s)_{\sigma(i)} )^2$.
% The behavior of this expression is quantified by Theorem \ref{thm:shuffling-2d} with $f(x_1,x_2) = (x_1-x_2)^2$. 
%Theorem \ref{thm:shuffleddMV_convergence} shows instead of converging to $d_{MV}(X_t,X_s)$, in the misalignment case $\hat{d}^2_{MV}(\mathbf{X}_{t}, P_{\sigma} \mathbf{X}_{s})$ will converge to shuffled-$d_{MV}(X_t,X_s)$ defined as below:
%To formalize the impact of vertex permutations on the $d_{MV}$ dissimilarity, we define its shuffled variant below.
\begin{definition}[independent-$d_{MV}$ Distance]
\label{def:independent_d_MV}
Let $X,Y\in L^2(\Omega)$ be latent position random variables. The independent-$d_{MV}$ distance is defined as
$$\inddmv(X,Y) = \min_{W\in\mathcal{O}^{d\times d}} \left\|\EE_{(X',Y')\sim \mu_X \otimes \mu_Y}[(X'-WY')(X'-WY')^\top]\right\|,$$
where $\mu_{X}\otimes \mu_{Y}$ is the product of the marginal measures of $X$ and $Y$.
\end{definition}
The impact of vertex shuffling on a network is represented as similarity transformation by a uniformly-distributed permutation matrix, by which the matrix $A_t$  is mapped to  $PA_t P^\top$. Since the eigendecomposition is permutation-equivariant, a shuffling of the vertices between times $t$ and $s$ on the $\hat{d}_{MV}$ distance transforms 
$\hat{d}_{MV}(\hat{\bX}_t,\hat{\bX}_s)$ into $\hat{d}_{MV}(\hat{\bX}_t,P\hat{\bX}_s).
$

The following result concerns the case of positive scalar LPPs, for which the calculation of $d_{MV}$ is especially simple since the minimizing rotation is always given by $W=1$ and the spectral norm becomes equivalent to the Frobenius norm. 

\begin{thm}\label{thm:independentdMV_convergence}
Consider a positive scalar LPP $\{X_t ; t\in [m]\}$. 
Let $\{\bX_{t}\}_{t\in[m]}$ be latent position matrices with $n$ rows drawn i.i.d from this LPP.
Let $\sigma$ be a random permutation uniformly distributed on $S_n$, and $P_n$ its corresponding permutation matrix. Then for any $\epsilon > 0$,  
$$ \PP \left( \left| \hat{d}^2_{MV}(\mathbf{X}_{t}, P_{n} \mathbf{X}_{s})-  \inddmv[2](X_{t}, X_{s}) \right| \geq \epsilon \right) \leq  \frac{12}{n \epsilon^2} + \frac{4}{n \epsilon}.$$
\end{thm}
We can also consider the case where only a portion of the vertices are misaligned, or where only a portion of the vertex alignments are known, as in the following.
\begin{definition}[$\alpha$-shuffling]\label{def:alpha-shuffling}
Let $\alpha\in(0,1]$, and define $\alpha_n=\lfloor \alpha n\rfloor$. Let $\sigma_{\alpha_n}$ be a uniformly distributed random permutation on $\{n-\alpha_n+1,\ldots,n\}$ with corresponding permutation matrix $P_{\alpha_n}$. Define the {\em $\alpha$-shuffling} by $\sigma\in S_n$ satisfying $\sigma(i)=i$ for $i=1,\ldots, n-\alpha_n$ and $\sigma(j)=\sigma_{\alpha_n}(j)$ for $j=n-\alpha_n+1,\ldots, n$. The corresponding permutation matrix of such an $\alpha$-shuffling is given by 
$P_{n,\alpha}:=I_{n-\alpha_n}\oplus P_{\alpha_n}$. When $\alpha=1$, we call this a \emph{shuffling}.
\end{definition}

\begin{definition}[$\alpha$-shuffled TSG]
\label{def:alpha-shuffled-tsg}
Let $\bA_1,\ldots,\bA_T$ be the adjacency matrices of a TSG on a fixed, aligned vertex set with $n$ vertices. Given $\alpha\in(0,1]$, draw $T$ independent $\alpha$-shufflings as in Definition \ref{def:alpha-shuffling}, with permutation matrices $\{P^1_{n, \alpha},...,P^T_{n, \alpha}\}$.
We call the TSG with adjacency matrices given by 
$$P^1_{n, \alpha} \bA_1 (P^1_{n, \alpha})^\top ,\ldots,  P^T_{n, \alpha} \bA_T (P^T_{n, \alpha})^\top$$ an $\alpha$\emph{-shuffled TSG}. When $\alpha = 1$, we shorten this to a \emph{shuffled TSG}.
\end{definition}
For positive scalar latent positions, when we consider the $\hat{d}_{MV}$ distance applied to $\bX_t$ and an $\alpha$-shuffled $\bX_s$, $\hat{d}_{MV}$ satisfies
\begin{equation}\label{equ:decompose_partial_shuffling}
\hat{d}_{MV}(\bX_t,P_{n,\alpha}\bX_s)^2= \frac{1}{n}\| \bX_t - (I_{n - \alpha_n} \oplus P_{\alpha_n} )   \bX_{s}\|^2_F 
= \frac{n - \alpha_n}{n} \hat{d}^2_{MV}(\bX_t^{(1)},\bX^{(1)}_s) + \frac{\alpha_n}{n}\hat{d}^2_{MV}( \bX_t^{(2)} , P_{\alpha_n} \bX_s^{(2)} ),
\end{equation}
where $\bX_t = [(\bX_t^{(1)})^\top, (\bX_t^{(2)})^\top]^\top$. Namely, $\bX_t^{(1)}$ is the first $n-\alpha_n$ rows of $\bX_t$, which are aligned with $\bX_s^{(1)}$, while $\bX_t^{(2)}$ and $\bX_s^{(2)}$ contain the remaining $\alpha_n$ rows, which have been shuffled. 
\begin{corollary}\label{Cor:partial shuffling}
Let $\alpha\in(0,1]$ and consider a positive scalar LPP. 
For any $\epsilon > 0$, we have 
$$ \mathbb{P} \left( \left| \hat{d}^2_{MV}(\mathbf{X}_{t}, P_{n,\alpha} \mathbf{X}_{s})-  
\left( \left(1-\frac{\alpha_n}{n}\right) d^2_{MV}(X_{t},X_{s}) +  \frac{\alpha_n}{n}\, \inddmv[2](X_{t}, X_{s})\right)
\right| \geq \epsilon \right) \leq  \frac{1}{(n-\alpha_n) \epsilon^2} + \frac{12}{\alpha_n \epsilon^2} + \frac{4}{\alpha_n \epsilon}.$$
\end{corollary}

Inspired by this result, we define $\alpha$-$d_{MV}$ to be a convex combination of $d_{MV}$ and ind-$d_{MV}$:
\begin{definition}[$\alpha$-$d_{MV}$]
$$
\alpha-d^2_{MV}(X_t,X_s) := (1-\alpha) ~ d^2_{MV}(X_t,X_s) + \alpha\cdot \inddmv[2](X_t,X_s). 
$$    
\end{definition}

Besides the extreme cases of perfect vertex alignment and random shuffling, it is not uncommon to have a set of seed vertices where the alignment is known or can be correctly estimated. 
%In these intermediate cases between full knowledge of the correspondence and complete misalignment, we see that the dependency-weakening effect of shuffling actually behaves in a reasonable way. I
In settings where the alignment is important, Corollary \ref{Cor:partial shuffling} motivates improving the alignment even when it cannot be perfectly estimated.



%This quantity can be used to understand the $\hat{d}_{MV}$ applied to $\alpha-$shuffled-TSG when number of vertices $n$ tends to infinity.  
%And we see that applying Euclidean mirror with $d_{MV}$ to $\alpha$-shuffled TSG amounts to CMDS on $\alpha$-shuffled-$d_{MV}$ distance matrix as $n$ tends to infinity. 
%(\textcolor{red}{why the extreme behaviors of the models? this motivates sections on A and L, DONE!})
% In practice, when the TSG are shuffled or the true alignment are unknown, we usually estimate the alignment among graph by graph matching.
% Specifically, given two graphs $\bA_t$ and $\bA_s$ with the same number of vertices, and assuming the existence of a true alignment matrix $P_0$, graph matching seeks $\hat{P}$ that solves the optimization problem:
% $$
% \hat{P}:= \arg \min_{P \in S_n}\|\bA_t - P\bA_sP^{\top}\|_F.
% $$
% Once $\hat{P}$ is obtained, the graph $\hat{P}\bA_s\hat{P}^{\top}$ is used for further operations.
% However, despite its practical use, graph matching presents both computational challenges and conceptual limitations.
% Graph matching is equivalent to a quadratic assignment problem which is NP-hard \cite{vogelstein:_fast}, and more critically, even if the global minimizer is found, the objective function does not guarantee that $\hat{P}$ approximates the true alignment $P_0$.
% This is because the objective aims to minimize the Frobenius norm difference between $A$ and $PBP^{\top}$, without ensuring consistency between $\hat{P}$ and $P_0$.
% If there is no assumption that the true alignment makes $A$ and $B$ similar, there is no guarantee that $\hat{P}$ will be close to $P_0$.  
% In such cases, a phenomenon known as "phantom alignment" may occur, where $\hat{P}$ aligns $A$ and $B$ too similar, resulting in a deviation from the true alignment.
% As we will demonstrate later, this issue is particularly severe in the Atlanta model.
% To understand this, we first establish a conceptual connection between $\hat{d}_{MV}$, graph matching, and optimal transport.

\subsection{Graph matching and optimal transport}\label{sec:Mirrors_Matching_Transport}
When vertex alignments between networks are unknown, it is natural to seek methods for comparing networks which mitigate the lack of correspondence. We now consider two such mitigating methods, optimal transport on the latent position distributions and graph matching for the adjacency matrices, which have significant parallels in this setting. 

Consider comparing the marginal distributions of the latent positions at two different times. One approach to a such a comparison is \emph{optimal transport}, whose objective in this setting to find an alignment that minimizes the Wasserstein-$p$ distance $W_p$, between these random variables, defined for a pair of measures $\mu$ and $\nu$ on a complete, separable metric space as follows:

\begin{definition}\label{def:Wasserstein_distance_measures}
Let $\mu, \nu$ be two measures with finite $p$-moments on $\RR^d$. Let $\gamma(\mu, \nu)$ be a \emph{coupling} between these measures, namely a measure on $X\times X$ with the product $\sigma$-algebra, with marginals $\mu$ and $\nu$. Let $\Gamma(\mu, \nu)$ be the set of all such couplings.  The Wasserstein $p$-distance between $\mu$ and $\nu$ is defined by
$$W_p(\mu, \nu)=\inf_{\gamma \in \Gamma(\mu, \nu)
} \left[\mathbb{E}_{(x,y) \sim \gamma} \|x-y\|_p^p\right]^{1/p}.$$ When we have random vectors $X,Y$ supported on $\RR^d$, we will also write $W_p(X,Y)$ to denote the Wasserstein $p$-distance between their induced measures, $W_p(\mu_X,\mu_Y)$.
\end{definition}

Because of the random dot product graph's inherent nonidentifability of latent positions with respect to orthogonal transformations, we consider minimizing the Wasserstein $p$ distance over orthogonally-transformed latent positions, which we denote by $\procwp[p]$. 

\begin{definition}\label{def:procwp-distance}
Consider two random vectors $X \in \RR^d$, $X \sim \mu$ and $Y\in \RR^d$, $Y \sim \nu$. Let $W \in \mathcal{O}^{d \times d}$ be given, and let $\nu_{W}$ denote the probability measure of $Z=WY$. We define 
\begin{align*}
\procwp[p] (X,Y) : = \min_{W \in \mathcal{O}^{d\times d} } W_p(X,WY) &= \min_{W\in\mathcal{O}^{d\times d}} \inf_{\gamma\in\Gamma(\mu,\nu_W)}[\EE_{(x,z)\sim\gamma}\|x-z
\|_p^p]^{1/p}\\
&= \min_{W\in\mathcal{O}^{d\times d}}\inf_{\gamma\in\Gamma(\mu,\nu)}[\EE_{(x,y)\sim\gamma}\|x-Wy\|_p^p]^{1/p}.
\end{align*}
\end{definition}

When we have two point clouds of vectors in $\RR^d$, $\{x_1,\ldots,x_n\}$ and $\{y_1,\ldots,y_n\}$, we can arrange these as the rows of the matrices $\bX,\bY\in\RR^{n\times d}$. The empirical measure for one of these point clouds is given by $$\mu_{\bX}=\frac{1}{n}\sum_{i=1}^n \delta_{x_i},$$ where $\delta_x$ is the point mass at $x$. Now the Wasserstein distance between these empirical measures gives us a distance between matrices via 
$\hat{W}_p(\bX,\bY):= W_p(\mu_\bX,\mu_\bY).$ It turns out that the minimization over couplings is always achieved by a permutation matrix in this setting, which prompts the following definition for pairs of matrices.

\begin{definition}\label{def:Wasserstein_empirical}
Let $\bX\in \RR^{n\times d}$ and $\bY\in\RR^{n\times d}$ have rows $x_1,\ldots,x_n$ and $y_1,\ldots,y_n$, respectively. Then 

\begin{equation}
    \label{eq:wassp}
    \hat{W}_p(\bX,\bY):=\min_{\sigma\in S_n} \left(\frac{1}{n}\sum_{i=1}^n \|x_i-y_{\sigma(i)}\|_p^p\right)^{1/p},
\end{equation}
\end{definition}
%so in the case $p=2$, this is simply $W_2(\bX,\bY):=\min_{P\in \mathcal{P}_n}\frac{1}{\sqrt{N}}\|\bX-P\bY\|_F$. Of course, in the case that the sampled vectors are not directly observed, but rather estimated from the ASEs of adjacency matrices, we need to incorporate a Procrustes alignment, giving us



% Because of the of random dot product graph's inherent nonidentifability of latent positions with respect to orthogonal transformations, we consider minimizing the Wasserstein $p$ distance over orthogonally-transformed latent positions, which we denote by $\procwp[p]$. 
Again, because of orthogonal nonidentifiability for random dot product graphs, we adapt Definition \ref{def:procwp-distance} for the case of latent position matrices as follows.
\begin{definition}\label{def:procwp_matrix_distance}
% For two random variables $X \in \RR^d$ and $Y\in \RR^d$:
% $$
% \procwp[p] (X,Y) : = \min_{W \in \mathcal{O}^{d\times d} } W_p(X,WY).
% $$
Given two latent position matrices $\bX,\bY \in \RR^{n \times d}$, we define 
$$
\procwphat[p] (\bX,\bY) : = \min_{W \in \mathcal{O}^{d\times d} } W_p(\bX,\bY W) = \min_{W \in \mathcal{O}^{d\times d}}\min_{\sigma \in S_n} \left( \frac{1}{n} \sum^n_{i=1}\| (\bX)_i -  (\bY W)_{\sigma(i)} \|_p^p \right)^{\frac{1}{p}},
$$
so when $p = 2$, this is simply
$$
\procwphat (\bX,\bY) =  \frac{1}{\sqrt{n}} \min_{W \in \mathcal{O}^{d\times d}}\min_{\sigma \in S_n} \left(\sum^n_{i=1}\| (\bX)_i -  (\bY W)_{\sigma(i)} \|^2\right)^{\frac{1}{2}} = \frac{1}{\sqrt{n}} \min_{W \in \mathcal{O}^{d\times d}}\min_{P \in \mathcal{P}_n} \|\bX-P\bY W\|_F  .
$$
\end{definition}



%If we don't know the true alignment in the TSG, as mentioned before, we will apply graph matching to the graphs and use the graph matched TSG for the mirror procedure: 
 
 Since we generally do not observe the matrices of latent positions themselves, we can consider $\procwphat$ evaluated at the estimated latent position matrices $\Xhat$ and $\Yhat$. The minimization over permutations in $\procwphat(\Xhat, \Yhat)$ amounts to aligning the estimated latent positions directly. Alternatively, we can determine a vertex relabeling by aligning the observed adjacency matrices themselves. This is the {\em graph matching} problem \cite{conte04:_thirt, lyzinski15:_relax, sgm_jofc} in which, given two $n \times n$ adjacency matrices $A$ and $B$, the goal is to find a permutation matrix $P$ that minimizes disagreements between the two adjacency matrices:
\begin{equation}\label{eq:graph_matching_objective}
P_{GM}=\arg\min_{P\in \mathcal{P}_n} \|A-PBP^{\top}\|_F.
\end{equation}
Suppose we have adjacency matrices $\bA_s$ and $\bA_t$ from an LPPTSG at times $s$ and $t$, and suppose $P_{GM}$ is given by Eq. \eqref{eq:graph_matching_objective}. To understand pairwise dissimilarities for the latent position process while still accounting for and mitigating vertex misalignments, we can use graph matching between adjacency matrices to induce an alignment between the estimated latent positions, and then compute the resulting $d_{MV}$ dissimilarity. For a positive scalar latent position process, this reduces to  
$$
\frac{1}{\sqrt{n}} \| \hat{\bX}_t - P_{GM}\hat{\bX}_s\|_2 .
$$
% For a positive scalar latent position process, this reduces to 
% To gain insight of the above quantity consider a denoise of above: $\hat{d}_{MV}$ with ``graph matching" using estimated probability matrix.
% $$
% \frac{1}{\sqrt{n}} \min_{W \in \mathcal{O}^{d \times d}}\| \hat{\bX}_t - P_2\hat{\bX}_sW\|_2, \text{ where }
% P_2 = \arg \min_{P} \| \hat{\bX}_t \hat{\bX}^{\top}_t- P \hat{\bX}_s \hat{\bX}^{\top}_s  P^{\top} \|_F = \arg\max_{P}  \sum_i \lambda^2_i\left( \hat{\bX}^{\top}_t P \hat{\bX}_s \right) .
% $$
% Now by considering a positive scalar LPP, the above reduces to:
% $$
% \frac{1}{\sqrt{n}} \| \hat{\bX}_t - P_2\hat{\bX}_s\|_F, \text{ where }
% P_2 = \arg \min_{P} \| \hat{\bX}_t \hat{\bX}^{\top}_t- P \hat{\bX}_s \hat{\bX}^{\top}_s  P^{\top} \|_F = \arg\max_{P}  \sum_i \lambda^2_i\left( \hat{\bX}^{\top}_t P \hat{\bX}_s \right) .
% $$
% Then we consider an ``matching" using estimated latent positions: 
% %$$
% %\frac{1}{\sqrt{n}}\min_{W \in \mathcal{O}^{d \times d}} \| \hat{\bX}_t - P_1\hat{\bX}_sW\|_2,\text{ where }P_1 = \arg \min_{P} \| \hat{\bX}_t - P\hat{\bX}_s\|_F= \arg\max_{P} \sum_i \lambda_i\left( \hat{\bX}^{\top}_t P \hat{\bX}_s \right).$$
If we minimize this quantity more directly and use the Frobenius norm, via 
$$
\min_{P \in S_n}\frac{1}{\sqrt{n}}\| \hat{\bX}_t - P\hat{\bX}_s\|_F,$$
then we recover
$$\min_{\sigma \in S_n} \left(\frac{1}{n} \sum^n_{i=1} \| (\hat{\bX}_t)_i - (\hat{\bX}_s)_{\sigma(i)}\|^2\right)^{\frac{1}{2}} = W_2(\mu_{\hat{\bX}_t}, \mu_{\hat{\bX}_s} ).
$$
As such, Wasserstein-type metrics can be viewed as a more direct comparison between latent position distributions compared to applying the $d_{MV}$ distance to the graph-matched adjacency matrices. The following two theorems show that Wasserstein-type distances between latent position distributions can be well-approximated by Wasserstein-type distances on the estimated latent positions.

% The above heuristic suggests connections between an aligned $\hat{d}_{MV}$ distance and the Wasserstein distance between adjacency spectral embeddings.
% Thus we consider Wasserstein-type distances as an alternative metric for pairs of random variables in a LPP to extract the dynamics of the TSG.
% Formally, we define:
% \begin{definition}[min-$W_p$ distance]
% For two random variables $X \in \RR^d$ and $Y\in \RR^d$:
% $$
% \text{min-}W_p (X,Y) : = \min_{W \in \mathcal{O}^{d\times d} } W_p(X,WY).
% $$
% And a corresponding estimator is
% for two estimated latent position matrices at $\bX \in \RR^{n \times d}$ and $\bY\in \RR^{n \times d}$:
% $$
% \text{min-}\hat{W}_p (\bX,\bY) : = \min_{W \in \mathcal{O}^{d\times d} } W_p(\bX,\bY W) = \min_{W \in \mathcal{O}^{d\times d}}\min_{\sigma \in S_n} \left( \frac{1}{n} \sum^n_{i=1}\| (\bX)_i -  (\bY W)_{\sigma(i)} \|^p \right)^{\frac{1}{p}}  .
% $$
% when $p = 2$:
% $$
% \text{min-}\hat{W}_2 (\bX,\bY) =  \frac{1}{\sqrt{n}} \min_{W \in \mathcal{O}^{d\times d}}\min_{\sigma \in S_n} \left(\sum^n_{i=1}\| (\bX)_i -  (\bY W)_{\sigma(i)} \|^2\right)^{\frac{1}{2}} = \frac{1}{\sqrt{n}} \min_{W \in \mathcal{O}^{d\times d}}\min_{P \in S_n} \|\bX-P\bY W\|_F  .
% $$
% \end{definition}




% Next we will show that the Wasserstein-type distance between between ASEs of two graphs actually converges to the Wasserstein distance on $X_t$ and $X_s$. 

% Optimal transport distance, also called Wasserstein-1 distance, between the measures $\mu$ and $\nu$ on the metric space $(\mathcal{M},d)$ as $$ \inf_{\pi \in \Gamma(\mu,\nu)} \EE_{(x,y)\sim \pi} d(x,y). $$ Here $\Gamma(\mu,\nu)$ is the set of all couplings between the measures $\mu,\nu$, which is to say the set of measures on $\mathcal{M}\times \mathcal{M}$ with marginal distributions $\mu$ and $\nu$. The Kantorovich dual formulation is given by 
% $$ \max_{f\text{ 1-Lip}} \EE_{x\sim \mu} f(x) - \EE_{y\sim \nu} f(y),$$ where the maximization is over all 1-Lipschitz functions $f:\mathcal{M}\rightarrow\mathcal{M}$ (with respect to the distance $d$).

% In the case that $\mu,\nu$ are discrete uniform distributions, i.e., $\mu = \frac{1}{n}\sum_{i=1}^n \delta_{x_i},\quad \nu = \frac{1}{n}\sum_{j=1}^n \delta_{y_j},$ where $\{x_i\}, \{y_j\}$ are samples from $\mathcal{M}$ (that may not all be distinct), then 
% $$ W^1(\mu,\nu) = \min_\pi \sum_{i,j} \pi(x_i,y_j)d(x_i,y_j). $$
% Since $\sum_i \pi(x_i,y_j) = \nu(y_j)=1/n$ and $\sum_j \pi(x_i,y_j)=\mu(x_i)=1/n$ for any $i,j$, 
% %\textbf{(TC:if $x_i= x_j$, then $\mu(x_i)= \mu(x_j) = \frac{2}{n}$, right? if as you said $x_i$ and $y_i$ may not all be distinct. I think i know what you're saying, please revise a bit to humor me. :-) )} 
% we see that $\pi$ is nothing more than a doubly-stochastic matrix scaled by $1/n$. Since the objective function is linear, there must always be a solution at a (scaled) permutation matrix, which amounts to a matching between $\{x_i\}$ and $\{y_j\}$. 
% That is, for empirical measures $\hat{\mu}_{\hat{\bX}_t}, \hat{\mu}_{\hat{\bX}_s} $:
% $$
% W_1(\hat{\mu}_{\hat{\bX}_t}, \hat{\mu}_{\hat{\bX}_s} ) = \min_{\sigma \in S_n} \frac{1}{n} \sum^n_{i=1} \| (\hat{\bX}_t)_i - (\hat{\bX}_s)_{\sigma(i)}\|.
% $$
% For finite graphs, optimal transport on the latent positions amounts to a version of graph matching that only considers the underlying latent positions, and ignores the actual appearance or non-appearance of edges in the two graphs. 
% This notion of distance extends to the case of continuous distributions, so we can make sense of distributional limits as $n\rightarrow\infty$, even though that becomes complicated for the case of graph matching. When the edges in the two graphs are conditionally independent of one another, this is a perfectly reasonable way to compare the two graph distributions. 
% %(TC:Just to remind even when graphs are conditional independent, this still lose information as W-1 distance between latent positions only consider marginal measure of latent positions while ignore the joint measure) 
% In the case that there is additional information about the relationship between the two graphs encoded in the edges (say, when edge formation is correlated between the two graphs), ignoring the actual edge appearances will, of course, lose information.

% In order to render this practical for the case of two observed graphs, we cannot just compare the true samples of the latent positions, but we have to estimate them in some way. \textcolor{red}{Add words around these two paragraphs and a closing paragraph.}


\begin{thm}\label{thm:W1_continuous}
Let $\bX,\bY\in \RR^{n\times d},$ and suppose that $\hat{\bX}, \hat{\bY}\in \RR^{n\times d}$ satisfy $$ \min_{W\in \mathcal{O}^d}\|\hat{\bX}-\bX W\|_{2\rightarrow\infty}\leq \epsilon,\quad \min_{W\in\mathcal{O}^d}\|\hat{\bY}-\bY W\|_{2\rightarrow\infty}\leq \epsilon.$$ Then 
$$ |\procwphat[1](\hat{\bX},\hat{\bY})-\procwphat[1](\bX,\bY)|\leq 2\epsilon. $$
\end{thm}

\begin{remark}
    Note that Theorem~\ref{thm:W1_continuous} is a deterministic statement, and we make no assumption about the joint distribution of the two samples comprising the rows of $\bX$ and $\bY$ or the relationship between these matrices and their estimates $\hat{\bX},\hat{\bY}$.
\end{remark}


\begin{thm}\label{thm:W1_as_converge}
Let $\mu_X,\mu_Y$ be distributions on $B_d^+$, the subset of the unit ball in $\RR^d$ with all nonnegative entries. Let $\bX,\bY\in\RR^{n\times d}$ be matrices whose rows come from iid samples from $\mu_X,\mu_Y$, respectively. Suppose that $\EE_{\mu_X}[xx^\top], \EE_{\mu_Y}[yy^\top]$ are positive definite. There is a constant $C_d$ such that for sufficiently large $n$,
$$\PP[|\procwphat[1](\bX,\bY)-\procwp[1](\mu,\nu)|>C_d n^{-1/(d+2)}] \leq \exp(-n^{-d/(d+2)}).$$
Furthermore, let $A_X\sim\mathrm{RDPG}(\bX), A_Y\sim\mathrm{RDPG}(\bY)$. Let $\hat{\bX}, \hat{\bY}$ be the ASEs for these $A_X,A_Y$, respectively. Then for any $a>0$, there is a constant $C_{a,d}>0$ such that for sufficiently large $n$, with probability at least $1-2n^{-a}$,
$$|\procwphat[1](\hat{\bX},\hat{\bY})-\procwp[1](\mu,\nu)| \leq C_{a,d} n^{-1/(d+2)}.$$
\end{thm}

This theorem tells us that we can consistently recover $\procwp[1]$ distances between the distributions at various times induced by an LPP from observations of networks at the corresponding times. The rate of this convergence appears to scale badly with the dimension of the latent positions, however this can be improved when the distributions admit a density or additional derivatives: see, e.g., \cite{chewi2025statistical}.

\begin{comment}
What's more we can calculate the $W_1$ distance for both models. 

\begin{thm}\label{thm:A-L-W1}
We consider the properties of $\psi_{W_1}$ for London model and Atlanta model.
\begin{itemize}
    \item
For London model $\{X^L_{t_i} ; i \in[m] \}$, assume $\delta_m = \frac{1}{m}$ then for any $i,j \in [m]$,
$$
\left(\mathcal{D^L}_{W_1}\right)_{i,j} = W_1(X^L_{t_i},X^L_{t_j}) = |\EE [X^L_{t_i}] - \EE[X^L_{t_j}]| = 
\begin{cases}
p|\frac{i}{m}-\frac{j}{m}| \quad i,j \leq t_m^*, \\
p ( \frac{t_m^*}{m}- \frac{i}{m})+q(\frac{j}{m}-\frac{t_m^*}{m})   \quad i < t_m^* < j, \\
q |\frac{i}{m}-\frac{j}{m}| \quad t_m^* < i,j . 
\end{cases}
$$
Thus $\mathcal{D^L}_{W_1}$ is exactly Euclidean-1 realizable, further assume $\frac{t_m^*}{m} = t^*$,
$$
\psi_{W_1}(t_i) = \psi_Z(t_i) ~~ \forall i \in [m]. 
$$  

\item 
For Atlanta model $\{X^A_{t_i} ; i \in[m] \}$, 
$$
\left(\mathcal{D^A}_{W_1}\right)_{i,j} = W_1(X^A_{t_i},X^A_{t_j}) = 0 ~~ \forall ~ i, j \in [m].
$$
\end{itemize}

    
\end{thm}

\begin{figure}[htbp]
    \centering
        \includegraphics[width=.45\linewidth]{figures/Illustration/W1-both.pdf}
    \caption{ The left figure is applying $W_1$ distance mirror to one realization of time series of graphs generated from London model with $n = 100, p=0.3, q=0.9, m = 30, t^* = 0.5, c_L = 0.1, \delta_m = 0.9/30$. The black dots are the first dimension of CMDS result on true $W_1$ distance matrix $\mathcal{D}_{W_1}$, the blue are the first dimension of CMDS result on estimated $W_1$ distance matrix $\hat{\mathcal{D}}_{W_1}$, which is obtained by applying $W_1$ distance on all pairs of ASE results of TSG $\hat{\mu}_{\hat{\bX}_{t_i}}$. Red line is the $0.9\psi_Z$ function. Black dots lie exactly on the red line and blue dots falls closely to the black dots. Compare to left two figures in \ref{fig:London-mirrors}, $W_1$ has smaller variance. 
    The right figure is the similar procedure on Atlanta model with $n=1000, p=0.05, q = 0.45, c_A = 0.8,m=30, t^* = 0.5, N = 50$. We see clearly $W_1$ induced mirror is not informative for $t^*$.}
    \label{fig:W1-mirrors-both}
\end{figure}

\end{comment}
%In this section, we see that applying Euclidean mirror with $d_{MV}$ graph matched TSG is closely related to apply Wasserstein type distance on the ASE of TSG.

% To summarize, for a positive scalar LPP, as the number of vertices $n$ grows, we have two key asymptotic results for any two time points $t$ and $s$:
% \begin{enumerate}
%     \item The $\hat{d}_{MV}$ mirror applied on an $\alpha$-shuffled-TSG converges to the $\alpha$-shuffled-$d_{MV}$ distance on the LPP.
%     \item The $\hat{d}_{MV}$ mirror applied on a graph-matched TSG is conceptually linked to the $W_1$ distance between the empirical distributions of the ASEs, and this distance converges to the $W_1$ on the LPP.
% \end{enumerate}
% To investigate how different strategies for handling vertex correspondence affect the resulting mirror and any subsequent statistical inference, we compare the CMDS embeddings obtained from three distinct dissimilarities matrices, namely, $d_{MV}$ mirror, $\alpha$-shuffled-$d_{MV}$ mirror and $W_1$ mirror. 
%\begin{itemize}
%    \item The \textbf{$d_{MV}$ mirror}, derived from the $d_{MV}$ distance matrix, which applies the mirror to the TSG assuming true vertex alignment.
%    \item The \textbf{$\alpha$-shuffled-$d_{MV}$ mirror}, from the $\alpha$-shuffled-$d_{MV}$ distance matrix, representing the approach for a TSG that has been $\alpha$-shuffled.
%    \item The \textbf{$W_1$ mirror}, from the Wasserstein-1 ($W_1$) distance matrix, which ignores vertex alignment and is related to the ``denoise'' graph matching approach.
%\end{itemize}

% To illustrate these differences, we construct two positive scalar LPPs that serve as edge cases: the London model, $\{X^L_t\}$, and the Atlanta model, $\{X^A_t\}$.
% For these models, the $d_{MV}$, $\alpha$-shuffled $d_{MV}$, and $W_1$ distances, along with their resulting CMDS embeddings, can be analyzed directly.
% We chose these models specifically because the different mirror exhibit dramatically different behaviors for each.
% The theoretical properties of these mirrors will be analyzed in Section \ref{sec:T2M}.
% Subsequently, in Section \ref{sec:Simulations}, we will present simulation results for realized TSGs with a finite $n$.
% It is important to clarify that while applying the mirror to a graph-matched TSG is conceptually related to the Wasserstein distance, there is no guarantee yet that they will yield similar results or performance.
% We use the Wasserstein distance as a conceptual heuristic for understanding the graph matching case.
% Therefore, Section \ref{sec:Simulations} will empirically compare the performance of the $W_1$ representation with that of $\hat{d}_{MV}$ applied to graph-matched TSGs.

%have dramatically different behavior with those four mirrors. 

%And this section shows the different behavior of using $W_1$ metric to induce mirror for changepoint localization.
%And heuristically if we consider graph matching as noisy $W_1$ for both models, this section gives insights about how graph matching will perform for subsequent inference problem, not the problem of whether it can find true alignment. 

%For London model graph matching again as W1 can actually help with subsequent inference, since $W_1$ distance matrix reveal piecewise linear more clearly than dMV because it is exactly while dMV is asymptotic and we will see sometimes it outperforms the true alignment in later simulation section. Thus for London graph matching can't find the true alignment, but it actually benefits the subsequent inference problem. 

%However for Atlanta graph matching not only can't help subsequent inference problem, but also it will returns a flat line with all information washed out. Thus for Atlanta, graph matching can't find the true alignment and can't benefit subsequent inference problem at all. 


%And essentially graph matching as W1 distance means graph matching will ignore the correlation or dependence in the original graphs whether the correspondence is true or not. and it only considers the marginals. If the information is just in marginals, then ignore correspondence might be better as we see in London.  

%Further neither means we shouldn't use graph matching in practice, all models are wrong we will how graph matching wins in a simulated swarm data case.

% practitioners, the main take away is if correspondence is completely lost, try W1 distance induce mirror for subsequent dynamics exploration. If we are lucky to be close to London model setting, it can actually win. 


%{\bf AA}: ZL and I discussed the following. Because the optimal transport problem uses the regularized low-rank $\hat{X}$ to compute things, it's less noisy than graph-matching, which uses the unprocessed adjacencies directly. Therefore, we'd like to compare solving a variant of the graph matching problem with the $\hat{P}_t$ and $\hat{P}_{t'}$ matrices instead, to see how that looks. It will tells us a little bit about regularization (optimal transport) vs retaining some information in the joint distribution (graph matching).



%\section{Comparison between $d_{MV}$, shuffled-$d_{MV}$ and Wasserstein distance}

%In this section we will compare the mirror if we use different metrics: $d_{MV}$, $\text{shuffled-}d_{MV}$ and Wasserstein-1 distance($
%W_1(\mu, \nu)=\inf _{\gamma \in \Gamma(\mu, \nu)}\mathbf{E}_{(x, y) \sim \gamma} |x-y|
%$) on the same LPP, specifically Atlanta model and London model.



%\input{compare table of London and Atlanta}

%Note the W-1 distance between two random variables $X$ and $Y$ on $\RR$ equals to $\EE X-Y$ if $X$ stochastically dominate $Y$ \cite{de20211}. Also note $X$ stochastically dominate $Y$ is equivalent to there exists two identical distributed $X'$ and $Y'$ random variables such that $X'=Y'+\delta$ where $\delta \geq 0$. From this equivalence statement we know for any pair of $X_{t_i}$ and $X_{t_j}$ from London model if $i<j$, then $X_{t_j}$ stochastically dominate $X_{t_i}$ since their difference are always sum of Bernoulli random variables, which are positive. Thus W-1 distance between $X_{t_i}$ and $X_{t_j}$ are simply $\EE X_{t_j} - X_{t_i} $ for $ i < j$. 





\section{Two models: London and Atlanta}\label{sec:T2M}

%\textcolor{red}{intro both, give Theorem 6 and a paragraph out;linign the effect on the two models then a new section for shuffled dmv then sections with more details on the models}



%In this section we introduce two LPP models, as two edge cases, to illustrate how vertex 1-1 correspondence impact the subsequent localization inference problem. Both LPPs have a first-order changepoint. Specifically, London model, originally introduced in \cite{chen2024euclidean}, the changepoint is robust to lost of the true 1-1 correspondence, while in Atlanta model, true 1-1 correspondence is essential to the changepoint localization. 


\begin{figure}
    \centering
    \includegraphics[width=\textwidth]{figures/atlanta-london.pdf}
    \caption{Estimated mirrors for London and Atlanta models after vertex shuffling. Both models exhibit a clear, estimable changepoint in the setting without shuffling. After shuffling the vertices, the London model shows little change in the estimated mirror, while the Atlanta model loses all relevant signal.}
    \label{fig:londonvsatlanta}
\end{figure}

To investigate the effect of vertex misalignment on the recovery of the network dynamics, we introduce two positive scalar LPPs with related but distinct behavior: the \emph{London model}, previously studied in \cite{chen2023discovering}, and the new \emph{Atlanta model}. Associated to these latent position processes is the corresponding London or Atlanta LPP TSG, whose time-varying adjacency matrices can be observed.

Both latent position processes, and hence both the London and Atlanta time series of graphs, have a first-order changepoint $t^*$ which we wish to estimate from observations of the network adjacencies. 
%However, in the London model, signal is contained in the marginal distributions of the latent positions at each time, while in the Atlanta model, the joint distributions between times contain crucial signal for changepoint localization. 
To localize the common changepoint $t^*$ in both models, we consider different functions of network observations, ranging from simple statistics such as average degree to more sophisticated localization techniques based on Euclidean mirrors. We examine properties of Euclidean mirrors for each model, as well as the behavior of the average degree and the behavior of corresponding mirrors under different levels of vertex shuffling. 

Both models have one-dimensional latent position processes that are discrete-time, finite-state space random walks taking values in the unit interval. In the London model, the latent position process begins at a fixed point, and at each time, the latent position $X_t$ can either stay fixed or move one step to the right with probability $p_t$. Moreover, $p_t$ is fixed at $p \in (0,1)$ for all times $t<t^{*}$, and thereafter $p_t$ is fixed at a different value $q$. Thus $t^*$ is a first-order changepoint since the distribution of the LPP increments is different before and after $t^*$. The Atlanta model is similar, but the initial distribution of latent positions is uniform over the state space, and at each time there is an equal probability of moving to the state to the right or to the left with probability $p_t$ (only allowing one direction at the boundaries), which changes similarly to the London model. One of the key differences between these two models is that the marginal distribution of the latent positions is different from one time to the next in the London model, whereas the marginal distributions of the latent positions is fixed for all times in the Atlanta model. This means that while the marginal distributions are informative for the changepoint in the London model (in fact, they are sufficient, as we show in Theorem~\ref{thm:London_MLE}), in the Atlanta model, the marginal distributions contain no information about the changepoint. This results in crucial differences in the estimability of the changepoint after vertex shuffling. 

The London model has an approximating piecewise linear Euclidean mirror that exhibits a slope change at $t^*$, and this structure persists even when vertices undergo shuffling over time. The Atlanta model has an approximating Euclidean mirror that exhibits a similar piecewise linear structure and slope change at $t^*$ {\em when the vertices are correctly aligned}, but this structure and slope change are very rapidly corrupted as the fraction of shuffled vertices increases. This is summarized in Figure \ref{fig:londonvsatlanta}.



%We also consider the utility of average degree as the shuffled mirror, the $\alpha$-shuffled mirror, $W_1$ mirror and an additional simple representation, average degree of each graph, for both models. 

%or London the shuffled mirror is almost identical to the mirror and both are informative of $t^*$; while for Atlanta the shuffled mirror is non informative of $t^*$ at all while mirror is informative $t^*$. We will also analyze the property of $\alpha$-shuffled mirror and property of average degree of the TSG, which is invariant to any permutations on TSG.

% Before proceeding, we introduce the notation for mirrors induced by $d_{MV}$ and shuffled-$d_{MV}$ for London and Atlanta.
% Following the notation in Definition \ref{def:CMDS}, let $\mathbf{\Psi}^{1:c} \in \RR^{m \times c}$ denote the CMDS embedding into $c$ dimensions based on a distance matrix, and let $\psi^1 \in \RR^m$ be its first column.
% Additionally, define $\tilde{\psi}^1:[0,1]\to\RR^1$ as the linear interpolation of the points $\{(t_i,\psi^1(t_i))\}_{i=1}^m$.
% To distinguish mirrors induced by different metrics, we add subscripts. For instance, $\mathcal{D}_{d_{MV}}$ represents the distance matrix for an LPP using the $d_{MV}$ distance, and $\psi^1_{d_{MV}}$ is the first dimension of the CMDS embedding applied to $\mathcal{D}_{d_{MV}}$.
% This notation similarly extends to $d^2_{MV}$, shuffled-$d_{MV}$, or $\alpha$-shuffled-$d_{MV}$, $W_1$ distance.
% To differentiate between the two models, we use the superscripts $\mathcal{L}$ and $\mathcal{A}$. For example, $\mathcal{D^A}_{d^2_{MV}}$ refers to the $d^2_{MV}$ distance matrix for the Atlanta model $\{X^A_{t_i}; i\in[m]\}$.
% We primarily focus on CMDS embeddings into $c=1$; in these cases, we omit the superscript and, for example, write $\psi^L_{d_{MV}}$ for the first CMDS dimension applied to $\mathcal{D^L}_{d_{MV}}$, with $\tilde{\psi}^L_{d_{MV}}$ denoting its linear interpolation.
% When considering embeddings for $c>1$, we explicitly use $\mathbf{\Psi}^{1:c}$ to denote the matrix.
% Finally, for the graph at time $t \in [0,1]$ with $n$ nodes and adjacency matrix $\bA_t$, the average degree is defined as $\frac{1}{n}\sum_{k=1}^n \sum_{h\neq k} (\bA_t)_{k,h}$; we denote this map from time to a real number by $\psi_{\text{avg-degree}}:[0,1]\to\RR^1$.



% %The end goal of the mirror method is to provide a map from a graph to a vector at each time points and so that it preserved the dynamic of the graphs.
% %Mirror method is just finding such map that preserves the pairwise distance we defined on the graphs, furthermore on the random variables that generates the graphs in the LPP model context, thus different metric yields different maps, for either $d_{MV}$ or shuffled-$d_{MV}$ or Wasserstein type distance we will use later. Such map can be constructed using much simpler idea, for example a simple summary statistics on a graph, like average degree of a graph. 
% The most natural Euclidean representation of a stochastic process with a first-order changepoint is a piecewise linear function defined on $[0,1]$ with a slope change at $t^*$.
% Let's denote such function with slope change from $p$ to $q$ at $t^*$ as $\psi_Z$:
% $$\psi_Z^c (t) = \begin{cases}
%     pt +c \quad \text{if } t \leq t^*, \\
%     pt^*+(t-t^*)q +c \quad \text{if } t > t^*. 
% \end{cases} 
% $$ 
% Then we consider sampled $m$ time points $\{t_i : = \frac{i}{m}; i \in[m]\} \subseteq [0,1]$, its corresponding pairwise Euclidean distance matrix is denoted as $\mathcal{D}_Z$:
% $$
% \left(\mathcal{D}_Z\right)_{i,j}:=\psi^c_Z\left(\frac{i}{m}\right)-\psi^c_Z\left(\frac{j}{m}\right), ~~~ \forall i, j \in [m].
% %=
% %\begin{cases}
% %\frac{p}{m}|i-j| \quad \frac{i}{m},\frac{j}{m} \le t^*, \\
% %|\frac{p}{m}(mt^*-i)+\frac{q}{m}(j-mt^*)| \quad \frac{i}{m} \le t^*<\frac{j}{m}, \\
% %\frac{q}{m}|i-j|\quad t^*<\frac{i}{m},\frac{j}{m}.
% %\end{cases}
% $$
% The result of applying CMDS to the $m \times m$ distance matrix $\mathcal{D}_Z$ is a one-dimensional vector in $\mathbb{R}^m$ and it is always centered by construction.
% Thus the specific embedding $\psi^{c_m}_Z$ technically varies with $m$. 
% %This means that for any $m$, the coordinates of the resulting vector are given by $\psi^{c_m}_Z(t_i) = \psi^0_Z(t_i) - c_m$, where $c_m = \frac{1}{m}\sum_{i=1}^m \psi^0_Z(t_i)$.
% %Note that the centering constant $c_m$ depends on $m$, which implies that the specific embedding $\psi^{c_m}_Z$ technically varies with $m$.
% However, this variation is just a constant vertical shift and preserves the pairwise Euclidean distances between all points. 
% %Since a constant shift preserves the pairwise Euclidean distances between all points, this dependency on $m$ is not a central concern for our analysis.
% For simplicity, we will henceforth use the notation $\psi_Z$ to represent the centered embedding obtained from CMDS for any given $m$, i.e., $\sum_{i=1}^m \psi_Z(t_i) = 0$ for any given $m$.
%We show both London and Atlanta has the property that asymptotically their $\tilde{\psi}^L_{d_{MV}}$ or $d^2_{MV}$ will converge to $\psi_Z$ up to a translation determined by $m$. 
\subsection{London model}\label{sec:London}
The London model is formally defined as follows.
%g with an illustration in Figure \ref{fig:L-illustration}.
\begin{model}[London Model LPP and TSG]\label{model:London} Let $m\geq 2$ be an integer. Let $c_L\geq0, \delta_m>0$ be constants satisfying $c+\delta_m m\leq 1$. Let $t^*\in(0,1)$. Denote $t^*_m = \lfloor t^*m\rfloor $, and put $t_i = i/m$ for $i\in\{1,\ldots,m\}$ and define the London model LPP as
\begin{align*}
X^{L}_0 & = c_L \quad \text{with probability 1}, \notag \\
X^{L}_{t_i}&=\begin{cases}
X^{(m)}_{t_{i-1}}+ \delta_m \quad \text{with probability $p_{t_i}$} \\
X^{(m)}_{t_{i-1}} \quad \text{with probability $1-p_{t_i}$},
\end{cases}
\end{align*}
where $p_{t_i} = p$ for $i < t_m^*$ and $p_{t_i} = q$ for $i \geq t_m^*$.

\end{model}

In both this section and the next, we will consider a true 1-dimensional piecewise-linear mirror that will be approximated by various methods, capturing the changepoint at $t^*$ as well as the different slopes $p$ and $q$ before and after the changepoint. We denote this as:

$$
\psi_Z^c (t) = \begin{cases}
     pt +c &\quad \text{if } t < t^*, \\
     pt^*+(t-t^*)q +c &\quad \text{if } t \geq t^*. 
 \end{cases} 
$$ 
When $c=0,$ we will drop the superscript.

%\begin{figure}[htp]
%  \centering
%  \includegraphics[height=4.5cm]{figures/Illustration/100london_sample_paths_p=0.3_q=0.9_m=30_c=0.1_tstar=0.5.pdf}%
%  \hspace{-7mm} 
  %\includegraphics[height=4.8cm]{figures/Illustration/adj100london_graphs_p=0.3_q=0.9_m=30_c=0.1_tstar=0.5.pdf}
%  \caption{The sample paths and realized TSG on London model with $p = 0.3$, $q = 0.9$, $m = 30$ and $t^* = 0.5$, $c_L = 0.1$, $\delta_m = 0.9/30$. The left figure shows $100$ sample paths where each state is represented by a dot and states are connected for each path; The right figure shows the evolution of adjacency matrices of realized time series of graph generated with number of vertices $n = 100$. These figures show that after the changepoint at $t^{*}=0.5$, $X^L_t$ and the density of connections in the adjacency matrices both increase more rapidly, indicating a structural shift both in sample paths and network connectivity.}
%  \label{fig:L-illustration}
% \end{figure}
%\cite{chen2024euclidean} has studied the $d_{MV}$ distance's properties of this model, here we summarize it and introduce more properties when consider shuffled-$d_{MV}$.
%We will show that the mirror induced by either distance will be informative for $t^*$. 
% In the following theorem, we present the properties of $\tilde{\psi}^{L}_{d_{MV}}$, $\tilde{\psi}^L_{\text{shuffled-}d_{MV}}$,
% $\tilde{\psi}^L_{\alpha\text{-shuffled-}d_{MV}}$,
% $\psi^L_{W_1}$,
% namely, CMDS result on $\mathcal{D^L}_{d_{MV}}$, $\mathcal{D^L}_{\text{shuffled}-d_{MV}}$, $\mathcal{D^L}_{\alpha-\text{shuffled}-d_{MV}}$ and $\mathcal{D^L}_{W_1}$.
% Additionally, we examine the average degree representation of the TSG, denoted as $\psi^L_{\text{avg-degree}}$.\\
% %[TOdo, ADD there exists $w_m = \pm1$ so that ..]
% \textcolor{red}{NEED PADDING and context for this, as well as transition.}

Now let us describe mirror estimation for the London model under various metrics: The $\alpha-d_{MV}$ distance for $\alpha\in[0,1]$, which includes the unshuffled case all the way up to the completely shuffled case; the Wasserstein distance $W_1$; and the differences in expected average degree $d_{\mathrm{deg}}(X_t,X_s)=(n-1)|\EE[X_t]^2-\EE[X_s]^2|.$

\begin{thm} \label{thm:London_main}
Consider a London model LPP with $c_L=0, \delta_m=1/m$, and $t^*=t_m^*/m$. Let $\psi_{d}(t)$ denote the zero-skeleton mirror using the distance $d$, and let $\tilde{\psi}_d(t)$ be its linear interpolation between consecutive points. Then the Euclidean mirrors for this model satisfy:
\begin{itemize}
% \item for $d_{MV}$ distance: 
% %\begin{itemize}
% %\item $$
% %\left(\mathcal{D^L}^{(2)}_{d_{MV}}\right)_{i,j} = 
% %d^2_{MV}(X^L_{t_i},X^L_{t_j})=
% %\begin{cases}
% %p^2\left(\frac{i}{m}-\frac{j}{m}\right)^2+\frac{p-p^2}{m}|\frac{i}{m}-\frac{j}{m}| \quad i,j \leq t_m^*, \\
% %\left( p (\frac{t_m^*}{m}- \frac{i}{m})+q(\frac{j}{m}-\frac{t_m^*}{m}) \right)^2 +\frac{p-p^2}{m} | \frac{t_m^*}{m}- \frac{i}{m}|+\frac{q-q^2}{m}|\frac{j}{m}-\frac{t_m^*}{m}| \quad i < t_m^* < j, \\
% %q^2\left(\frac{i}{m}-\frac{j}{m}\right)^2+\frac{q-q^2}{m}|\frac{i}{m}-\frac{j}{m}| \quad t_m^* < i,j . 
% %\end{cases}
% %$$
% %\item
% as $m \to \infty$, $\tilde{\psi}^L_{d_{MV}}$ is asymptotically Euclidean 1-realizable with asymptotic mirror $\psi_Z$, that is, there exists $w \in \pm1$ so that
% $$
% \left(\mathcal{D^L}_{d_{MV}}\right)_{i,j} \to \left(\mathcal{D}_Z\right)_{i,j}~,\quad \sup_{t \in [0,1] } \bigg|\tilde{\psi}^{L}_{d_{MV}}(t)-w\psi_Z(t)\bigg| \to 0, \quad \text{for all }  i,j \in \{1,2,\cdots,m\};
% $$ 
%     %\end{itemize}

% \item for shuffled-$d_{MV}$ distance:
% %\begin{itemize}
%     %\item 
% %\begin{align*}
% %\left(\mathcal{D^L}^{(2)}_{\text{shuffled-}d_{MV}}\right)_{i,j} = 
% %  \text{shuffle-}  d^2_{MV}(X^L_{t_i},X^L_{t_j})=
% %  \begin{cases}
% %  p^2\left(\frac{i}{m}-\frac{j}{m}\right)^2+\frac{p-p^2}{m}\left(\frac{i}{m} {+} \frac{j}{m}\right) \quad i,j \leq t_m^*, \\
% %  \left( p ( \frac{t^*_m}{m}- \frac{i}{m})+q(\frac{j}{m}-\frac{t^*_m}{m}) \right)^2 +\frac{p-p^2}{m} | \frac{t^*_m}{m} {+} \frac{i}{m}|+\frac{q-q^2}{m}|\frac{j}{m} - \frac{t^*_m}{m}| \quad i<t_m^*<j, \\
% %  q^2\left(\frac{i}{m}-\frac{j}{m}\right)^2+\frac{q-q^2}{m}|\frac{i}{m} {+} \frac{j}{m}| + {\frac{2t^*_m}{m^2} \left( p-p^2 -q + q^2 \right)} \quad t_m^*<i,j, \\
% %  0 \quad i=j.
% %  \end{cases}
% %\end{align*}
% %\item 
% as $m \to \infty$, $\tilde{\psi}^L_{\text{shuffled}-d_{MV}}$ is asymptotically Euclidean 1-realizable with asymptotic mirror $\psi_Z$, that is, there exists $w \in \pm1$ so that:
% $$
% \left(\mathcal{D^L}_{\text{shuffled-}d_{MV}}\right)_{i,j}\to\left(\mathcal{D}_Z\right)_{i,j}~~,  \quad
% \sup_{t \in [0,1] } \bigg|\tilde{\psi}^L_{\text{shuffled}-d_{MV}}(t)-w\psi_Z(t)\bigg| \to 0,
% \quad  \text{for all }  i,j \in \{1,2,\cdots,m\};
% $$
% %\end{itemize}

\item For $\alpha-d_{MV}$ distance with $0\le\alpha\leq 1$: As $m \to \infty$, the mirror is asymptotically Euclidean 1-realizable with mirror $\psi_Z$, namely, there is a $w\in\{\pm1\}$ such that:
$$
\sup_{t \in [0,1] } \bigg|\tilde{\psi}_{\alpha-d_{MV}}(t)-w\psi_Z(t)\bigg| \to 0;
$$

\item For $W_1$ distance: The $W_1$ distance is exactly Euclidean-1 realizable, and
$$
\psi_{W_1}(t_i) =\EE[X_{t_i}] =  \psi_Z(t_i)\quad  \forall\; i \in [m]; 
$$  

\item For expected average degree distance:
$$
\sqrt{ \frac {\psi_{\text{deg} }\left(t_i\right) }{n-1}} = \psi_Z(t_i)\quad \forall\;i\in[m].
$$

%The points $\{\left( t_i,  \psi_{\text{avg-degree}}\left(t_i\right) \right)^n_{i=1}\}$ lie on the $\sqrt{n-1}\psi_Z + c$.     
\end{itemize}
\end{thm}

Theorem \ref{thm:London_main} shows as $m \to \infty$, the first dimension of both $\psi_{d_{MV}}$, $\psi_{\alpha-d_{MV}}$ and $\psi_{\text{ind-}d_{MV}}$ converge to $\psi_Z$ uniformly, indicating that shuffling does not affect the mirror's ability to localize $t^*$. 
Further, $\psi_{W_1}$ and the square root of $\psi_{\text{deg}}$ lie on $\psi_Z$ exactly when $t_m^* = mt^*$, so vertex alignments are not necessary for changepoint localization in the London model.
%The reason is that the information in the marginal distribution of the London model is enough for localizing $t^*$.
%Specifically, observe that $\EE[X^L_{t_i}] = \psi_Z(t_i)$ for all $i \in [m]$, that is the evolution of expectation in London model is piecewise linear with slope change at $t^*$, and thus informative of $t^*$.
%What's more, as we will see in the simulation section, using square root of $\psi_{\text{avg-degree}}$ for localization performs the best among all representations for London model because with RDPG assumption, average degree is an unbiased estimator for $\EE[X^L_{t_i}]^2$ with rather small variance. 

In Figure \ref{fig:London-mirrors}, we compare $\psi_{d_{MV}}$,  $\psi_{0.2-d_{MV}}$, $\psi_{\text{ind-}d_{MV}}$, $\psi_{W_1}$ and $\sqrt{\psi_{\text{deg}}/(n-1)}$ (black dots) with their estimates (blue dots), derived from a realized TSG from London model.
The function $0.9\psi_Z$ is also plotted as a red line for reference.
The estimates are the first dimension of CMDS applied to pairwise distance matrices of the ASEs.
% For the TSG with true alignment, the estimated mirror $\hat{\psi}^L_{d_{MV}}$ is obtained from the matrix $\hat{\mathcal{D^L}}_{\hat{d}_{MV}}$. 
% We then shuffle the realized TSG with 20\% and 100\% of vertices, obtaining 20\%-shuffled-TSG and shuffled-TSG respectively.
% For these, we compute the distance matrices $\hat{\mathcal{D^L}}_{\hat{d}_{MV}-20\%\text{-shuffled-TSG}}$ and $\hat{\mathcal{D^L}}_{\hat{d}_{MV}\text{-shuffled-TSG}}$, respectively.
% The first CMDS dimension of these matrices yields the mirrors $\hat{\psi}^L_{\hat{d}_{MV}\text{-}20\%\text{-shuffled-TSG}}$ and $\hat{\psi}^L_{\hat{d}_{MV}\text{-shuffled-TSG}}$.
% These estimates corresponds to their theoretical counterparts $\psi^L_{20\%\text{-shuffled-}d_{MV}}$ and $\psi^L_{\text{shuffled-}d_{MV}}$.
% Finally, we consider two methods that are invariant to vertex alignment.
% The estimated $\hat{\psi}_{W_1}$ is the first dimension of CMDS results on $\hat{\mathcal{D^L}}_{W_1}$, which is obtained by applying $W_1$ distance on all pairs of empirical distributions formed by ASE, estimating $\psi^L_{W_1}$.
% The value $\hat{\psi}_{\text{avg-degree}}$ is the square root of each graph's average degree, which estimates the $\sqrt{\EE\psi^L_{\text{avg-degree}}/n-1}$.
With $n=100$, all estimated mirrors lie close to their theoretical values, and these theoretical values themselves lie close to the asymptotic mirror $\psi_Z$, shown in red. In particular, they all show approximately piecewise linear structure and a slope change at $t^*$. Among these mirrors, the average degree appears to have the best performance.


\begin{figure}[htbp]
    \centering
        \includegraphics[width=.9\linewidth]{figures/Illustration/London-5-metrics.pdf}
    \caption{Comparison of $\psi_{d_{MV}}$,  $\psi_{0.2-d_{MV}}$, $\psi^L_{\text{ind-}d_{MV}}$, $\psi^L_{W_1}$ and $\sqrt{\psi_{\text{deg}}/(n-1)}$ in black and their estimates in blue based on one realized TSG generated from London model with $n = 100$, $p = 0.3$, $q = 0.9$, $m = 30$ and $t^* = 0.5$, $c_L = 0.1$, $\delta_m = 0.9/30$.  
    All mirrors show clear piecewise linear structure with slope change at $t^* = 0.5$.
    Among them, the average degree mirror lies closest to its expected value.}
    \label{fig:London-mirrors}
\end{figure}

One of the key properties of the London model is that no information is lost from shuffling. We can make this statement precise by showing that the marginal counts for each time are sufficient for the LPP, so that no information about $p,q,$ and $t^*$ is lost when the vertex correspondence is removed. This also leads to a maximum likelihood estimator for the changepoint.

\begin{thm}\label{thm:London_MLE}
Let $\{X_i\}$ be an iid sample from the London model, and for each $t\in \{1,\ldots, m\}$, let 
$$
c_t(k):= \frac{\#\{i: X_i(t)=c+\delta k\}}{n},
$$
$0\leq k\leq m$. 
Then there exists a function $h$ such that the likelihood $\mathrm{lik}(\{X_i(t)\}_{i,t}|p,q,t^*)=h(\{c_t(k)\}_{t,k}|p,q,t^*)$, so the marginal counts $c_t(k)$ are sufficient statistics for the London Model. For each $t\in\{1,\ldots, m\}$, the MLE estimates for $\hat{p}$ and $\hat{q}$ assuming that the changepoint occurs at time $t$ are given by
$$
\hat{p}(t):=\frac{1}{t-1}\sum_{j=0}^{t-1} j \frac{c_{t-1}(j)}{n},\quad \hat{q}(t):= \frac{1}{m-t+1}\sum_{k=0}^m k\frac{c_m(k)-c_{t-1}(k)}{n}.
$$
\end{thm}

We can interpret the estimator $\hat{p}(t)$ as the proportion of all opportunities for the $i$th latent position to increase up to time $t-1$ at which it actually increased, averaged over all $n$ latent positions. The estimator $\hat{q}(t)$ is similar, but only counts increases after time $t-1$. In the next section, we study a model where shuffling has much more serious inferential consequences.

\subsection{Atlanta model}\label{sec:Atlanta}

The Atlanta model is similar to the London model, but allows for movement of the latent positions in both directions. We now give the precise definition.

\begin{model}[Atlanta Model]\label{mod:Atlanta} 
Let $N \geq 2$ be a fixed number of states, and let $\delta_N = \frac{c_A}{N-1}$, where $c_A \in (0,1)$ is a constant. Let $m \geq 2$ be the number of time points, and set $t_i = i/m$ for $i \in [m]$. Let $t^* \in (0,1)$, and let $t^*_m = \lfloor t^*m \rfloor$. The Atlanta model LPP is given by:
\begin{align*}
X_{t_0} &\sim \mathrm{Uniform}(\{0,\delta_N,\ldots,(N-1)\delta_N=c_A\}), \\
\text{Middle: if $X_{t_{i-1}} \in (0, c_A)$, } X_{t_i} &=\begin{cases}
X_{t_{i-1}} + \delta_N & \text{with probability $p_{t_i}$}, \\
X_{t_{i-1}} - \delta_N & \text{with probability $p_{t_i}$}, \\
X_{t_{i-1}} & \text{with probability $1 - 2p_{t_i}$},
\end{cases} \\
\text{Left boundary: if $X_{t_{i-1}} = 0$, } X_{t_i} &=\begin{cases}
X_{t_{i-1}} + \delta_N & \text{with probability $p_{t_i}$}, \\
X_{t_{i-1}} & \text{with probability $1 - p_{t_i}$},
\end{cases} \\
\text{Right boundary: if $X_{t_{i-1}} = c_A$, } X_{t_i} &=\begin{cases}
X_{t_{i-1}} - \delta_N & \text{with probability $p_{t_i}$}, \\
X_{t_{i-1}} & \text{with probability $1 - p_{t_i}$},
\end{cases}
\end{align*}
where $p_{t_i} = p$ for $i < t^*_m$, and $p_{t_i} = q$ for $i \geq t^*_m$, with $p \neq q$.
\end{model}

The Atlanta model is parameterized by $N$, $m$, $p$, $q$, $t^*$, and $c_A$. 
It is a Markov process with the uniform distribution as a stationary distribution. 
%with transition matrix $T_p$ before the changepoint $t^*_m$ and $T_q$ after the changepoint, where 
%$$
%T_p :=
%\begin{bmatrix}
%1-p & p   & 0   & 0   & \cdots & 0 \\
%p   & 1-2p & p   & 0   & \cdots & 0 \\
%0   & p   & 1-2p & p   & \cdots & 0 \\
%\vdots & \vdots & \vdots & \vdots & \vdots & \vdots \\
%0   & 0   & \cdots & p   & 1-2p & p \\
%0   & 0   & \cdots & 0   & p   & 1-p
%\end{bmatrix}_{N \times N}.
%$$
%Note that the row vector $[1/N, \ldots, 1/N]$ is a left eigenvector for both $T_p$ and $T_q$. 
Since the process starts in this distribution, the marginal distribution remains unchanged for all times. Namely, $X_{t_i} \overset{\mathcal{\text{dist}}}{=} X_{t_0}$ for any $i \in [m]$.
This already suggests that any method that utilizes only marginal information will not be informative for $t^*$.

% In the following theorem, we present the properties of $\psi^{A}_{d^2_{MV}}$, $\psi^A_{\text{shuffled-}d_{MV}}$,
% $\mathbf{\Psi}^A_{\alpha\text{-shuffled-}d_{MV}}$,
% $\psi^A_{W_1}$,
% namely, CMDS result on $\mathcal{D^A}_{d^2_{MV}}$, $\mathcal{D^A}_{\text{shuffled-}d_{MV}}$, $\mathcal{D^A}_{\alpha\text{-shuffled-}d_{MV}}$ and $\mathcal{D^A}_{W_1}$.
% Additionally, we examine the average degree representation of the TSG, denoted as $\psi^L_{\text{avg-degree}}$..
% Additionally, we examine the average degree representation, denoted as $\psi^A_{\text{avg-degree}}$.

The following theorem provides an analogue of Theorem~\ref{thm:London_main} for the Atlanta model.

\begin{thm}\label{thm:Atlanta-main}
Consider an Atlanta model LPP for some $N,m,c_A,$ and $t^*=t_m^*/m$. The distances and Euclidean mirrors for this model satisfy:
%We consider the properties of maps $\psi^A_{d^2_{MV}}$, $\psi^A_{\text{shuffled}-d_{MV}}$, and $\psi^A_{\text{avg-degree}}$ for Atlanta model. For the following statements we will abuse notation and drop the superscript $A$ for simplicity.  
\begin{itemize}
\item For the $d_{MV}$ distance:
%\begin{itemize}
%\item
\begin{comment}
   $$
\left(\mathcal{D^A}^{(2)}_{d_{MV}} \right)_{i,j}= 
d^2_{MV}(X^A_{t_i},X^A_{t_j})=
\begin{cases} 
\frac{c^2}{(N-1)^2N} \tr\left(T_p^{|i-j|} M \right) \quad i,j<t_m^*, \\

\frac{c^2}{(N-1)^2N} \tr\left(T_p^{|i-t_m^*|}T_q^{|t_m^*-j|} M \right)  \quad i<t_m^*<j, \\

\frac{c^2}{(N-1)^2N} \tr \left(T_q^{|i-j|} M \right) \quad t_m^*<i,j, 
\end{cases}
$$
where $(M)_{i,j} = (i-j)^2$.
Note for non-negative integer $2 \leq k < N$, we have 
\begin{equation} \label{equ:m1}
\mathrm{tr }\left(T_p^k\,M\right) =  2(N-1)kp -4\sum_{t=2}^k\frac{(-1)^t}{t-1}\binom{k}{t} \binom{2t-4}{t-2}p^t,
\end{equation}
\begin{equation} \label{equ:m2}
\mathrm{tr }\left(T_p^kT_q^l\,M\right) =2(N-1)(kp+lq)  -\sum_{d=2}^{k+l}\sum_{t=\max(0,d-l)}^{\min(d,k)}(-1)^{d}\frac{4}{d-1} \binom{2d-4}{d-2}\binom{k}{t}\binom{l}{d-t}p^tq^{d-t}.
\end{equation}
 
\end{comment}
With fixed $m$, there exists a $w \in \{\pm1\}$ such that as $N \to \infty$,
$$
\max_{i \in [m]}\bigg| \frac{N(N-1)}{2c_A^2m}\psi_{d^2_{MV}}(t_i) -  w\psi_Z(t_i) \bigg| \to 0.
$$

%\end{itemize}    

\item For the ind-$d_{MV}$ distance:
$$ \text{ind-}d^2_{MV}(X_{t_i}, X_{t_j}) = 
   \begin{cases} 
     2\,\mathrm{Var}(u) = \frac{c_A^2}{6}\frac{N+1}{N-1} , & i \neq j,\\
     0 & i = j,
   \end{cases} 
$$
and the first $m-1$ dimensions of CMDS applied to these distances may be written as
$$
\mathbf{\Psi}_{\text{ind}-d_{MV}} = \sqrt{\frac{c_A(N+1)}{6(N-1)}} \bU_{m-1},
$$
where $\bU_{m-1}\in \RR^{m\times m-1}$ is any matrix satisfying $\bU^{\top}_{m-1}\bU_{m-1} = I_{m-1}$ and $e^\top \bU_{m-1}=0$, for $e\in\RR^m$ the matrix of all ones.

\item For $\alpha$-$d_{MV}$ distance, $\alpha\in[0,1)$: the first dimension of CMDS satisfies
$$
\psi_{\alpha -d_{MV}} = \sqrt{1-\alpha + \frac{\alpha c_A(N+1)}{12(N-1)\lambda_1}   }      \psi_{d_{MV}},
$$ 
where $\lambda_1$ is the largest eigenvalue of the doubly-centered matrix of squared $d_{MV}$ distances, $-\frac{1}{2}H \mathcal{D}_{d_{MV}}^{(2)}H$, $H=I-ee^\top/m$.

\item For $W_1$ distance: The distance matrix is exactly 1-Euclidean realizable with
$$
\psi_{W_1}(t_i) = 0 \quad\forall\; i \in [m].
$$

\item For expected average degree distance: The distance matrix is exactly 1-Euclidean realizable with
$$
\psi_{\text{deg}}(t_i) = 0 \quad \forall\;i \in [m].
$$
\end{itemize}
\end{thm}


%\begin{figure}[htp]
%  \centering
%  \includegraphics[height=4.5cm]{figures/Illustration/100atlanta_sample_paths_p=0.05_q=0.45_m=30_c=0.8_N=50_tstar=0.5.pdf}
%  \hspace{-7mm}
  %\includegraphics[height=4.8cm]{figures/Illustration/adj100atlanta_graphs_p=0.05_q=0.45_m=30_c=0.8_N=50_tstar=0.5.pdf}
%  \caption{The sample paths and realized TSG based on Atlanta model with $p = 0.05$, $q = 0.45$, $m = 30$ and $t^* = 0.5$, $c_A = 0.8$, $N = 50$, $\delta_N=\frac{0.8}{49}$. The left figure shows $100$ sample paths where each state is represented by a dot and states are connected for each path; The right figure shows the evolution of adjacency matrices of realized from sample paths on the left. This figure shows that after the changepoint at $t^{*}=0.5$, $X^A_t$, the sample paths undergoes increased oscillation because we increase jump probability from $0.05$ to $0.45$ and %the network exhibits more dynamic and variable connectivity patterns while 
% the density of graphs remains relatively consistent over time. }
 %\label{fig:A-illustration}
%\end{figure}

Note that in the Atlanta model, the number of time points $m$ and the number of states $N$ are two independent parameters.
Unlike Theorem \ref{thm:London_main}, the convergence of $\psi^A_{d^2_{MV}}$ occurs as $N$ grows, while $m$ can remain fixed.
However for both models, the partitioning in time and in latent position space becomes increasingly fine, and in-fills the respective space. In the Atlanta model, $\psi^A_{d^2_{MV}}$ is asymptotically piecewise linear, while $\psi^A_{d_{MV}}$ is not, as shown in Figure \ref{fig:Atlanta-mirrors} (second and third panels), though both exhibit an abrupt change at $t^*$.
In the simulation section, we demonstrate that the MSE is comparable for both $\hat{\psi}^A_{d^2_{MV}}$ and $\hat{\psi}^A_{d_{MV}}$, with $\hat{\psi}^A_{d^2_{MV}}$ performing slightly better.

Under shuffling, the first $m-1$ dimensions of CMDS are given by any orthonormal basis of the subspace of $\RR^{m}$ orthogonal to the vector of all ones, so this matrix is completely uninformative about $t^*$. Similarly, $\psi^A_{W_1}$ and $\psi^A_{\text{avg-degree}}$ are completely flat, since they only depend on the marginal distribution at each time, which stays fixed for this model. 

The behavior of the mirror under partial shuffling merits additional attention. When $\alpha\in(0,1)$, it is simply a rescaling of $\psi_{d_{MV}}^A$, but when $\alpha=1$, we obtain the uninformative, fully-shuffled mirror. While theoretically this mirror contains full information about $t^*$ even with $\alpha$ just slightly less than 1, in the case of estimating the mirror from an observed matrix, the noise level increases, as seen in panels 2 and 3 of Figure~\ref{fig:Atlanta-mirrors}.

% Furthermore, behavior of $\psi^A_{\alpha\text{-shuffled-}d_{MV}}$ is interesting.
% It is a rescaling of $\psi^A_{d_{MV}}$ as long as $0<\alpha < 1$, but when $\alpha=1$, its first $m-1$ dimension is any orthonormal basis of the $\RR^{m-1}$ space. 
% That is, partial shuffling with $\alpha < 1$ preserves all the signal up to rescaling, but when $\alpha = 1$, the signal at $t^*$ is lost entirely.
% To analyze the reason, we recall that CMDS on a distance matrix $\mathcal{D}$ involves double-centering the entrywise squared distance matrix to obtain the Gram matrix $- \frac{1}{2}H\mathcal{D}^{(2)}H$, see Definition \ref{def:CMDS}.
% The eigenvectors of the Gram matrix are then scaled by the square root of their corresponding eigenvalues, which is the CMDS result.
% The behavior arises because $\mathcal{D^A}_{\alpha \text{-shuffled-}d_{MV}}^{(2)}$ is a convex combination of $\mathcal{D^A}_{\text{shuffled-}d_{MV}}^{(2)}$ and $\mathcal{D^A}_{d_{MV}}^{(2)}$.
% Here, $\mathcal{D^A}_{\text{shuffled-}d_{MV}}^{(2)}$ is constant in the off-diagonal entries and zero on the diagonal.
% After double-centering, $\mathcal{D^A}_{\alpha \text{-shuffled-}d_{MV}}^{(2)}$ adds to the Gram matrix $-\frac{1}{2}H\mathcal{D^A}_{d_{MV}}^{(2)}H$ a term proportional to $H$.
% Since the Gram matrix $-\frac{1}{2}H\mathcal{D^A}_{d_{MV}}^{(2)}H$ and $H$ share the same eigenvectors, the new Gram matrix for $\mathcal{D^A}_{\alpha \text{-shuffled-}d_{MV}}^{(2)}$ also has the same eigenvectors.
% This modification only affects the eigenvalues, resulting in the scaling of the original CMDS outcome.

% Although shuffling with $\alpha < 1$ does not affect the signal in $\psi^A_{\alpha\text{-shuffled-}d_{MV}}$ without noise, it has a significant impact when noise is present, i.e., when the number of vertices is finite.
% As shown in Figure \ref{fig:Atlanta-mirrors} third panel, when consider applying mirror to $20\%$-shuffled TSG with $n=1000$, the blue dots are farther from the black dots compared to the blue dots in the second panel, indicating higher noise in the estimation.
% To understand this behavior, we decompose $\hat{\mathcal{D^A}}_{\hat{d}_{MV}-\alpha\text{-shuffled-TSG}}^{(2)}$ similarly as in equation \ref{equ:decompose_partial_shuffling}:
% $$
% \hat{\mathcal{D^A}}_{\hat{d}_{MV}-\alpha\text{-shuffled-TSG}}^{(2)}
% = \frac{(n-\alpha_n)}{n} \hat{\mathcal{D^A}}_{\hat{d}_{MV}}^{(2)} + \frac{\alpha_n}{n} \hat{\mathcal{D^A}}_{\hat{d}_{MV}\text{-shuffled-TSG}}^{(2)}.
% $$
% %where $\hat{\mathcal{D}}_{\hat{d}_{MV}-100\%-\text{shuffled-TSG}}^{(2)}$ contains entries computed using $\hat{d}_{MV}$ applied to the last $\alpha_n$ rows of $\{\hat{\bX}_t ;~ t \in [m]\}$, which are fully shuffled.
% %Similarly, $\hat{\mathcal{D}}_{\hat{d}_{MV}}^{(2)}$ contains entries computed using $\hat{d}_{MV}$ applied to the first $n-\alpha_n$ rows of $\{\hat{\bX}_t ;~ t \in [m]\}$, which are not shuffled.
% And note CMDS result of 
% $
% \frac{\alpha_n}{n} \hat{\mathcal{D^A}}_{\hat{d}_{MV}\text{-}100\%\text{-shuffled-TSG}} +
% \frac{(n-\alpha_n)}{n} \hat{\mathcal{D^A}}_{\hat{d}_{MV}}^{(2)}
% $
% is not simply a rescaling of $\hat{\mathcal{D^A}}_{\hat{d}_{MV}}^{(2)}$.
% Although each entry of $\hat{\mathcal{D^A}}_{\hat{d}_{MV}\text{-shuffled-TSG}}^{(2)}$ converges to $\mathcal{D^A}_{\text{shuffled-}d_{MV}}^{(2)}$ as $n$ grows large (by Theorem \ref{thm:shuffleddMV_convergence}), the matrix is not fully constant and does not necessarily share the same eigenvectors as the double-centered matrix of $\hat{\mathcal{D^A}}_{d_{MV}}^{(2)}$.
% Thus, as shown in Figure \ref{fig:Atlanta-mirrors}, the blue dots in the third panel are noisier compared to the blue dots in the second panel, though they still retain some signal.

\begin{figure}[htbp]
    \centering
        \includegraphics[width=\linewidth]{figures/Illustration/Atlanta-6-metrics.pdf}
    \caption{ Comparison of $\psi_{d^2_{MV}}$,  $\psi_{d_{MV}}$, $\psi_{0.2-d_{MV}}$, $\psi_{\text{ind-}d_{MV}}$, $\psi_{W_1}$ and $\psi_{\text{deg}}$ in black and their estimates in blue based on one realized TSG generated from Atlanta model with $n = 1000$, $p = 0.05$, $q = 0.45$, $m = 30$ and $t^* = 0.5$, $c_A = 0.8$, $N = 50$. 
    %The first figure shows the mirror induced by $d^2_{MV}$ mirror: black dots are the first dimension of CMDS of $\mathcal{D^A}_{d^2_{MV}}$ scaled by $\frac{N(N-1)}{2c_Am}$ from Theorem \ref{thm:Atlanta-main}; the blue dots are the first dimension of CMDS on $\hat{\mathcal{D}}_{d^2_{MV}}$, the red line is the $\psi_Z$ function.  The second figure shows the $d_{MV}$ mirror result on the same TSG.
    %Black dots are the first dimension of CMDS on $\mathcal{D^A}_{d_{MV}}$ while blue dots are first dimension of CMDS on $\hat{\mathcal{D}}_{d_{MV}}$.
    %The third figure shows $d_{MV}$ mirror result on the same TSG but with last 20\% of nodes shuffled, black dots are the first dimension of CMDS on true $\mathcal{D}_{\text{20\%-shuffled-}d_{MV}}$ while blue dots are first dimension of CMDS on $\hat{\mathcal{D}}_{ \hat{d}_{MV}\text{-}20\%\text{-shuffled-TSG}}$.
    %The fourth figure shows $d_{MV}$ mirror result on the same TSG but with 100\% of nodes shuffled: black dots are the first dimension of CMDS on $\mathcal{D}_{\text{shuffled-}d_{MV}}$ , while the blue dots are the first dimension of CMDS of $\hat{\mathcal{D}}_{ \hat{d}_{MV}\text{-}100\%\text{-shuffled-TSG}}$. Finally the last figure is the square root of the average degree of time series of graphs, black dots are the expectation $\EE\psi_{\text{avg-degree}}$ while blue dots are the realization: $\psi_{\text{avg-degree}}$. We see the left two figures are informative of $t^*$ while $d^2_{MV}$ lies closely to $\psi_Z$ and show piecewise linearity, $d_{MV}$ show more curvature at the boundary, but still show abrupt change at $t^*$. What's more $d_{MV}$ can tolerate small percent of misalignment as shown in the third figure: blue dots have more variance around black dots but still show change around $t^*$.
    Only the left three panels show an abrupt change at $t^* = 0.5$, indicating only mirrors using the true vertex alignment are informative about $t^*$.
    Since the ind-$d_{MV}$, $W_1$, and average degree mirrors either lose the true alignment or ignore it, they have little information about $t^*$.
    %Thus in Atlanta model, the representation that doesn't consider alignment won't work while $d_{MV}$ with not too much shuffling will still work.
    }
    \label{fig:Atlanta-mirrors}
\end{figure}

In this section we introduced two models, London and Atlanta, and thoroughly analyze their mirrors' properties.
While information about $t^*$ can be obtained from the observed adjacency matrices when the vertex correspondence is known, the situation is very different once the vertex correspondence is lost. The fundamental reason for this is that in London model the marginal distributions $\EE[X^L_t]$ are sufficient for $t^*$, while in the Atlanta model, the marginal is constant across time. As such, true alignment is essential for estimation of $t^*$ in the Atlanta model. We will investigate this in more detail through simulations in the next section.

%Another layer is, the true alignment serves only as a nuisance parameter, and the final task is the changepoint localization for time series of graphs, which aim to find the dissimilarity among time series of graphs, the graph matched time series of graphs will inevitably ``smoothing" the dynamics and weaken the signal. 

\section{Simulations}\label{sec:Simulations}
%We will use the same localization estimator thus the performance of $\hat{t}$ only depends on the embeddings. 
%The $\hat{t}$ is estimated with TSG with true alignment and $\hat{t}_\text{shuffled}$ is the estimated result with TSG with shuffling.
%For both model we will consider mirror method with true alignment and with shuffled alignment, which gives two embeddings and then perform localization on the embedding and get $\hat{t}$ and $\hat{t}_\text{shuffled}$.
%Thus the impact of vertex misalignment can be quantified by the difference between mean square error(MSE) of $\hat{t}$ and MSE $\hat{t}_\text{shuffled}$.
%We will show later in the simulation section while $ \text{MSE}(\hat{t}^L) \approx \text{MSE}(\hat{t}^L_\text{shuffled})$ in London, meaning $\hat{t}$ is robust to vertex misalignment for London; $ \text{MSE}(\hat{t}_\text{shuffled}) \gg \text{MSE}(\hat{t})$ in Atlanta, saying true alignment is essential for Atlanta. 
\subsection{Preliminaries and construction}\label{sec:Simulation_Preliminaries}

We aim to localize $t^*$ in an observed TSG generated from the London or Atlanta model.
We set the number of vertices, $n$, for the graphs and generate the time series of networks comprising $m$ graphs according to London model \ref{model:London} and Atlanta model \ref{mod:Atlanta}. \footnote{All of our codes are available at \url{https://github.com/Graph-matched-mirrors/Vertex-alignment-and-changepoint-localization-in-network-time-series}.} To prevent the graph from being too sparse or too dense, for the London model, we choose the starting point as $c_L=0.1$, resulting in $\delta_m=\frac{0.9}{m}$; for the Atlanta model, we set $u_N$ as a uniform distribution supported on discrete set $\{0.1, 0.1+\delta_N,...,0.9\}$, with $\delta_N = \frac{.9-.1}{N-1}$.
Throughout, we set $t^*$ to be $\frac{1}{2}$, select an even number of time points $m$, then choose $p$ and $q$.
In the Atlanta model, we always pick $p,q$ such that $0 \le p ,q \le 0.5$. 

%For fixed values of $n$, $m$, $p$, $q$, and $t_m^*=\frac{m}{2}$, we generate a time series of graphs from London model as follows.
%We fix a vector $\bX_{t_0} \in \R^{n\times 1}$, with all entries set to $0.1$, representing the starting point for the $n$ latent positions. For subsequent time points, we independently draw $n$ samples $\{u_1,\cdots,u_n\}$ from $\mathrm{Bernoulli}(p)$. The $i$th entry of $(\bX_{t_k})_i=(\bX_{t_{k-1}})_i$ if $u_i=0$, or $(\bX_{t_k})_i=(\bX_{t_{k-1}})_i+\frac{0.9}{m}$ if $u_i=1$.
%This process repeats until $\bX_{\frac{m}{2}}$ is reached.
%For $\bX_{\frac{m}{2}+1}$, we switch the jump probability to $q=0.3$. Eventually, we obtain $\{\bX_{t_1},\cdots,\bX_{t_m}\}$ as the realized latent position matrices for $n$ vertices. For each time, we generate an RDPG using the latent positions in $\bX_t$, obtaining a TSG as defined in \ref{def:tsg}.
%We denote the adjacency matrices of the aligned TSG as $\{\bA_{t_1},\cdots,\bA_{t_m}\}$. 

%As for Atlanta model, we fix $n$, $m$, $p$, $q$, $t_m^*=\frac{m}{2}$ and additional parameter number of states $N$. Now for $\bX_{t_1}$ we draw $n$ i.i.d sample from $u_N$, representing the latent positions for the fist time point. 
%For the following time points $t_k$, we will consider $i \in [n]$ one by one: if $(\bX_{t_{k-1}})_i$ equals to $0.1$ then $(\bX_{t_k})_i = (\bX_{t_k})_i + \delta_N$ with probability $p$ and  $(\bX_{t_k})_i = (\bX_{t_k})_i$ with probability $1-p$; if $(\bX_{t_{k-1}})_i$ equals to $0.9$, then $(\bX_{t_k})_i = (\bX_{t_k})_i - \delta_N$ with probability $p$ and  $(\bX_{t_k})_i = (\bX_{t_k})_i$ with probability $1-p$, and if $0.1< (\bX_{t_{k-1}})_i < 0.9$, then $(\bX_{t_k})_i = (\bX_{t_k})_i + \delta_N$ with probability $p$, $(\bX_{t_k})_i = (\bX_{t_k})_i - \delta_N$ with probability $p$ and $(\bX_{t_k})_i = (\bX_{t_k})_i$ with probability $1-2p$.
%We repeat this until $\bX_{t_{\frac{m}{2}}}$. For $\bX_{t_{\frac{m}{2}+1}}$ ,we switch the probability from $p$ to $q$. Thus we obtain $\{\bX_{t_1},\cdots,\bX_{t_m}\}$ as the realized latent position matrices for $n$ vertices.
%Then as above we generate a TSG. 

In the London model, for fixed values of $n$, $m$, $p$, $q$, and $t_m^*=\frac{m}{2}$, we generate a time series of latent position matrices according to Model~\ref{model:London}.
We fix a vector $\bX_{t_0} \in \R^{n\times 1}$, with all entries set to $0.1$.
Until a changepoint at $t_m^*=m/2$, each entry independently increases by $\delta_m$ with probability $p$ at each time step, and stays constant with probability $1-p$.
After the changepoint, this probability switches to $q$.
We obtain $\{\bX^L_{t_1},\cdots,\bX^L_{t_m}\}$ as the realized latent position matrices for London model. 

In the Atlanta model, we fix $n$, $m$, $p$, $q$, $t_m^*=\frac{m}{2}$ and number of states $N$.
For $\bX_{t_1}$ we draw $n$ i.i.d sample from $u_N$.
Before the changepoint at $t_m^*=m/2$, each entry can move up or down by $\delta_N$ with probability $p$ or remain unchanged with probability $1-2p$. 
If the entry is at either boundary $0.1$ or $0.9$, it will remain unchanged with probability $1-p$ and jump away from the boundary with probability $p$.  
After the changepoint, the probability $p$ is replaced by $q$.
We then obtain $\{\bX^A_{t_1},\cdots,\bX^A_{t_m}\}$ as the realized latent position matrices for Atlanta model.  For both models, the final time series of graphs, with adjacency matrices $\{\mathbf{A}_{t_1},\cdots,\mathbf{A}_{t_m}\}$, is generated by creating a Random Dot Product Graph (RDPG) from the latent positions $\mathbf{X}_t$ at each time point.


For a realized aligned TSG $\bA_{t_1},\cdots,\bA_{t_m}$ with number of vertices $n$, we denote the ASE results in 1 dimension as $\hat{\bX}_{t_1},\cdots,\hat{\bX}_{t_m}$.
First we consider $\hat{d}_{MV}$ on the aligned TSG: we construct 
$$
(\hat{\mathcal{D}}_{\hat{d}_{MV}})_{k,s} = \frac{1}{\sqrt{n}}\|\hat{\bX}_{t_k} - \hat{\bX}_{t_s}\|_F, ~~ \forall k,s \in [m],$$ then apply CMDS to this matrix into $c=1$ dimension, obtaining $\hat{\psi}_{\hat{d}_{MV}}$.

To study the effects of losing vertex alignment, we fix a proportion of the vertices to be shuffled, $\alpha$, then obtain an $\alpha$-shuffled TSG  $\{P_{n,\alpha}^1\bA_{t_1}(P_{n,\alpha}^1)^{\top},\cdots,P_{n,\alpha}^m\bA_{t_m}(P_{n,\alpha}^m)^{\top}\}$ as in Definition~\ref{def:alpha-shuffled-tsg}, but proceed as though the TSG were aligned. This gives a distance matrix
$$
\left(\hat{\mathcal{D}}_{\hat{d}_{MV}-\alpha\text{-shuffled-TSG}}\right)_{k,s} = 
\frac{1}{\sqrt{n}} \|P^k_{n, \alpha} \hat{\bX}_{t_k} - P^s_{n, \alpha}\hat{\bX}_{t_s} \|_F, ~~~~ \forall k,s \in [m] ,
$$ with corresponding mirror $\hat{\psi}_{\hat{d}_{MV}\text{-}\alpha\text{-shuffled-TSG}}$.
We also consider applying graph matching to the shuffled TSG before computing $\hat{d}_{MV}$: this is described further in Section \ref{sec:Simulations_MSE_Metrics}.  

For the first of our measures which does not make use of the vertex alignment, we use the $W_1$ distance on the empirical measures formed by ASEs of the TSG with $d_{\text{ASE}}=1$. This gives a distance matrix 
$$
\left(\hat{\mathcal{D}}_{W_1}\right)_{k,s} 
 {=}  \sum^n_{i=1}  \biggl|(\hat{\bX}_{t_k})_{(i)} - (\hat{\bX}_{t_s} )_{(i)}\biggr|, ~~~ \forall k,s\in[m],
$$
where $(\hat{\bX})_{(i)}$ denotes the $i$-th order statistic of $\hat{\bX}$, which is possible since this is a 1-dimensional vector. The corresponding mirror $\hat{\psi}_{W_1}$ is the first dimension of CMDS applied to this matrix. Lastly we consider the average degree of each graph as a trivial representation of the graph dynamics:
$$
\hat{\psi}_{\text{avg-degree}}(t_k)= \frac{1}{n} \sum^n_{i =1} \sum^n_{j \neq i} \left(\bA_{t_k}\right)_{i,j} ~~ \forall k \in[m]. 
$$

As shown in Theorem \ref{thm:London_main} and Theorem \ref{thm:Atlanta-main}, $t^*$ in both the London and Atlanta models can be revealed as the point of slope change in a piecewise linear mirror. 
This piecewise linearity is achieved either asymptotically or by applying specific transformations, depending on the notion of distance. 
To enable a principled comparison between representations, we apply transformations to make all mirrors approximately piecewise linear.
For the London model, we use the $\hat{d}_{MV}$ distance and choose a large value for $m$, since $\psi^L_{d_{MV}}$ becomes asymptotically piecewise linear as $m$ increases.
For the average degree in the London model, we use the square root of the average degree to achieve piecewise linearity.
%We will abuse the notation and refer to both of these as the average degree.
For the Atlanta model, we use $\hat{d}^2_{MV}$ instead of $\hat{d}_{MV}$, and select a large value for $N$ for the same reason.

We localize $t^*$ by estimating the slope change in our one-dimensional mirror estimate.
To that end, we design an $l_2$ localization estimator described in Algorithm \ref{alg:2}, which is revised from the $l_{\infty}$ localizer in \cite{chen2024euclidean}.
\begin{algorithm}
\caption{$l_2$ localization for piecewise linear data}
\label{alg:2}
\begin{algorithmic}[1]
\State Input: $m$ pairs of observations $\{(t_1,y_1),(t_2,y_2),\cdots,(t_m,y_m)\}$ with each $t_i\in \R$,$y_i\in \R$.
\State \textbf{for k from $2$ to $m-1$:}
\State \quad Define $$S_k:=\min_{\alpha, \beta_L, \beta_R \in \R} 
\sum^m_{ i = 1 } \left( \alpha + \beta_L(t_i-t_k)+(\beta_R-\beta_L)(t_i-t_k)I_{\{t_i > t_k\}} - y_i \right)^2.$$
\State Find the smallest $k_0$ such that $S_{k_0}= \min\{S_2,S_3,\cdots,S_{m-1}\}$ and set $\hat{t}=t_{k_0}$.
\State Output: $\hat{t}$. 
\end{algorithmic}
\end{algorithm}

Note that the minimization problem defining $S_k$ can be solved analytically by ordinary least squares for each fixed value of $k$. For each Monte Carlo realization, we obtain an estimated mirror, then apply this $l_2$ localization method to get an estimated changepoint $\hat{t}_{mc}$. We compute the squared error $(\hat{t}_{mc} - t^*)^2$, then run the Monte Carlo simulation $nmc$ times, recording the mean of these as the mean square error (MSE), as well as the sample standard deviation over the realizations: 
$$
\text{MSE} := \frac{1}{nmc} \sum_{mc=1}^{nmc} \left(\hat{t}_{mc} - t^*\right)^2, \quad 
\text{std} := \sqrt{\frac{1}{nmc-1} \sum_{mc=1}^{nmc} \left( \left(\hat{t}_{mc} - t^*\right)^2 - \text{MSE} \right)^2}.
$$
Because $\hat{t}_{mc}$ are i.i.d with expectation $\EE[(\hat{t} - t^*)^2]$, using a normal approximation, we can construct a confidence interval for the $\EE[(\hat{t} - t^*)^2]$ as:$$
\text{error bars} = \text{MSE} \pm \frac{1.96}{\sqrt{nmc}} \times \text{std},
$$ which we use to plot error bars for the MSE. 

To understand the scale of errors we should expect for MSE in this setting, we note that throughout all simulations, we choose $t^*= \frac{1}{2}$ and $m$ even, thus $\hat{t}_{mc} \in \{\frac{2}{m},\frac{3}{m}, \cdots,\frac{m-1}{m}\}$ (note in Algorithm \ref{alg:2} that the estimated changepoint cannot be a boundary point). As such, an estimator $u_m$ that randomly guesses among $\{\frac{2}{m},\frac{3}{m}, \cdots,\frac{m-1}{m}\}$ should have a chance MSE given by $\EE[ (u_m - \frac{1}{2})^2 ]$, which will be indicated as a vertical dashed red line in all simulations. Error rates above this level suggest a lack of consistency for the corresponding estimator, at least with the given sample size.


\subsection{MSE with different shuffling ratio}\label{Simulation_MSE}
First we study how the proportion of shuffled vertices $\alpha$ and signal strength $|p-q|$ affect the changepoint estimation in both models. 
We apply the $\hat{d}_{MV}$ mirror for the London model and the $\hat{d}^2_{MV}$ mirror for the Atlanta model to the $\alpha$-shuffled TSGs, treating them as if they were aligned. We consider 100 Monte Carlo runs ($nmc=100$) for each pair of parameters, and vary $\alpha$ and $|p-q|$ in each model to observe the effect, shown in Figure:~\ref{fig:MSEvsShuffleRatio}. 

For the Atlanta model in Figure \ref{fig:Atlanta-MSEvsShuffleRatio}, MSE increases dramatically with increasing $\alpha$. For $p=0.4$, $q = 0.1$ or $0.2$, when $\alpha < 15
\%$, the MSE for either model is better than chance, indicating that the $d^2_{MV}$ mirror is robust to a small amount of shuffling. This indicates that there is still a fair amount of signal in the aligned vertices, allowing for recovery in the case of the Atlanta model. When this alignment is not used, which is the case when we use $W_1$ or average degree, the performance is much worse than chance (see Table \ref{table:A}).

However, as soon as $\alpha\geq 15\%$, the signal disappears in the Atlanta model, and all methods perform worse than chance. The London model in Figure \ref{fig:London-MSEvsShufflingRatio} is robust to misalignment and the MSE is better than chance across proportions of shuffled vertices and values of $q\neq p$.

% It is worth noting that when $\alpha = 0$, for $p=0.4$ and $q = 0.1$ or $q = 0.2$, Atlanta model MSE is worse than chance level while London model's MSE is much smaller than chance level. 
% This is because for Atlanta model, $\psi^A_{d^2_{MV}} \approx \frac{m}{N^2}\psi_Z$ while for London model $\psi^L_{d_{MV}} \approx \psi_Z$, thus for same variance level(same $n$), $\psi^A_{d^2_{MV}}$ has smaller signal strength and is harder to be estimated, causing the slope change harder to localize.  
% Also note for Atlanta model, when shuffling ratio is 0 but $q = 0.35$ or $0.4$, MSE are worse than chance level, indicating that $\hat{t}$ more frequently selects points farther from the middle point $t^* = 0.5$, which we call as boundary effect.
% This occurs because, when the slope difference $p - q$ is small, the change near $t^*$ is subtle and difficult to localize.
% In such cases, the $l_2$ estimator tends to select points near the boundaries (e.g., $t = 1$ or $t = m = 30$), where the estimated mirror exhibits artificial slope changes.
% This may be due to higher variance near the boundaries and also because with finite $N$, even when $n = \infty$, the $\psi^A_{d^2_{MV}}$ is not exactly piecewise linear.
% %Similarly in London, with finite $m$, $\psi^L_{d^2_{MV}}$ is not exactly piecewise linear. 
% %Thus it is possible the localizer will tend to choose boundary as changepoint when the signal around $t^*$ is not strong enough. 
% Meanwhile we note the boundary effect persists as shuffling ratio increases. 
% Also we note $l_{\infty}$ localizer suffers less from the boundary effect as shown in Appendix Figure \ref{fig:atlanta_mse_vs_shuffle_appendix} the MSE are closer to chance level using $l_{\infty}$ especially when shuffling ratio increases. 

\begin{figure}[htbp]
    \centering
    \begin{subfigure}[b]{0.48\linewidth}
        \centering
        \includegraphics[width=\linewidth]{figures/Simulation_results/Atlanta_D2_l2.pdf}
        \caption{Atlanta model with $p=0.4$ varying $q$ and $\alpha$.}
        \label{fig:Atlanta-MSEvsShuffleRatio}
    \end{subfigure}
    \hfill
    \begin{subfigure}[b]{0.48\linewidth}
        \centering
        \includegraphics[width=\linewidth]{figures/Simulation_results/London_l2.pdf}
        \caption{London model with $p=0.4$ varying $q$ and $\alpha$.}
        \label{fig:London-MSEvsShufflingRatio}
    \end{subfigure}\caption{
    MSE versus shuffle ratio \(\alpha\) for two models: (a) Atlanta and (b) London.
    For both models, we fix \(p = 0.4\) with \(n = 300\), \(m = 40\), and \(t^* = 0.5\); in London, $c_L = 0.1, \delta_L = 0.9/40$, in Atlanta, we set \(N = 50\) and $\delta_A = 0.8$.
    Then \(\alpha\) ranges from 0 to 1 simulating increasing extent of vertex misalignment.
    Colored curves correspond to \(q = 0.1, 0.2, 0.35, 0.4, 0.5\), with each point averaged over $ nmc = 100$ and the error bars show 1.96 standard deviation.
    The red dotted horizontal line marks the chance level of $0.075$ for $m = 40$.
    %We apply the \(d_{MV}\) mirror localization procedure blindly as if the partially shuffled TSG were aligned, and report average MSE of the estimated location \(\hat{t}\).
    Atlanta shows high sensitivity to vertex misalignment, with MSE rapidly rising above chance once \(\alpha \ge 0.15\).
    In contrast, the London model exhibits robustness to misalignment, with MSE largely unaffected across all shuffle ratios and non-zero signal strengths.
    }
    \label{fig:MSEvsShuffleRatio}
\end{figure}



\subsection{MSE with different distances}\label{sec:Simulations_MSE_Metrics}
In this section we compare the performance of various distances with respect to the MSE for estimating the changepoint in the London and Atlanta models. We consider fully shuffled TSGs, and in addition to the four distances considered above ($\hat{d}_{MV}, \hat{d}_{MV-\alpha\text{-shuffled-TSG}}, W_1,$ and $d_{\text{avg-degree}}$),
we also consider applying graph matching to the shuffled TSG and then apply $\hat{d}_{MV}$ on the graph matched TSG. Let us denote the adjacency matrices of the shuffled TSG as 
$\bA'_{t_1},...,\bA'_{t_m}$, with ASEs 
$\{\hat{\bX}_{t_1}':=P_1\hat{\bX}_{t_1}, %\hat{\bX}_{t_2}':=P_{2}\hat{\bX}_{t_2} ,
... , \hat{\bX}_{t_m}':=P_{m}\hat{\bX}_{t_m}\}$.

A naive approach to graph matching simply aligns each pair of adjacency matrices, giving us
$$
\left(\hat{\mathcal{D}}_{\hat{d}_{MV}\text{-pairwise-GM-TSG}}\right)_{k,s} 
= 
\frac{1}{\sqrt{n}} \| \hat{\bX}_{t_k}' - P_{s \to k}\hat{\bX}_{t_s}'\|_F,~~~~ P_{s \to k} = \arg \min_{P} \| \bA'_{t_k} - P \bA'_{t_s} P^{\top} \|_F, ~~~~ \forall k,s \in [m].
$$ 
This requires $O(m^2)$ many graph matchings, which is computationally expensive, so we will use two alternative approaches to reduce this number. 

First we consider an ``all to one" graph matching in which $\bA'_{t_1}$ remains unchanged, but each $\bA'_{t_k}$ is matched to $\bA'_{t_1}$ to obtain $P_{k\to1}$. This gives
\begin{equation}\label{eq:alltoonegm}
\left(\hat{\mathcal{D}}_{\hat{d}_{MV}\text{-all to one-GM-TSG}}\right)_{k,s} 
=  \frac{1}{\sqrt{n}} \| P_{k\to1} \hat{\bX}'_{t_k} - P_{s\to1}\hat{\bX}'_{t_s}\|_F
\end{equation}
where $P_{k\to1} = \arg \min_{P} \| P \bA'_{t_k} P^{\top} -  \bA'_{t_1}\|_F $ for all $k \in [m]$. Note that $P_{1 \to 1} = I$. 

We also consider ``consecutive pair" graph matching, where we match $\bA'_{t_2}$ to the previous adjacency matrix $\bA^{\ast}_{t_1} = \bA'_{t_1}$ with solution $P_{2\to1^{\ast}}$, then match $\bA_{t_3}$ to the matched $\bA^{\ast}_{t_2} := P_{2\to1^{\ast}} \bA'_{t_2} P^{\top}_{2\to1^{\ast}}$ and find the solution $P_{3 \to 2^{\ast}}$, and so on, yielding
$$
\left(\hat{\mathcal{D}}_{\hat{d}_{MV}\text{-consecutive-GM-TSG}}\right)_{k,s}
=  \frac{1}{\sqrt{n}} \| P_{k\to (k-1)^{\ast}} \hat{\bX}'_{t_k} - P_{s\to(s-1)^{\ast}}\hat{\bX}'_{t_s} \|_F.
%\overset{\text{London and Atlanta}}{=}   \frac{1}{\sqrt{n}} \| \hat{\bX}_t - P_4\hat{\bX}_s\|_2  \overset{\text{1d}}{ = } \frac{1}{\sqrt{n}} \| \hat{\bX}_t - P_4\hat{\bX}_s\|_F,
$$
where $P_{k\to(k-1)^{\ast}} = \arg \min_{P} \| P \bA'_{t_k} P^{\top} -  \bA^{\ast}_{t_{k-1}} \|_F, ~~ \bA^{\ast}_{t_{k-1}} = P_{(k-1) \to (k-2)^{\ast}}\bA'_{t_{k-1}}P^{\top}_{(k-1) \to (k-2)^{\ast}} $ for all $k \in \{2,3,...,m\}$, define $P_{1 \to 0^\ast} := I$. 

\begin{table}[htbp]
% \centering
% \begin{minipage}{0.48\textwidth}
% \textbf{London}
\centering
\begin{tabular}{lccc||ccc}
&&\textbf{London}&&&\textbf{London}&\\
\hline
\textbf{Representation} & \textbf{Lower} & \textbf{Mean} & \textbf{Upper} &\textbf{Lower} & \textbf{Mean} & \textbf{Upper} \\
\hline
$d_{MV}$ with true alignment  & 0.0042 & 0.0051 & 0.0061 & 0.0242 & 0.0289 & 0.0335 \\
$d_{MV}$ with 100\% shuffling    & 0.0106 & 0.0125 & 0.0144 & 0.0178 & 0.0214 & 0.0251 \\
$d_{MV}$ with GM all to one       & 0.0091 & 0.0108 & 0.0125 & 0.0187 & 0.0225 & 0.0263 \\
$d_{MV}$ with GM consecutive      & 0.0229 & 0.0265 & 0.0300 & 0.0102 & 0.0129 & 0.0156 \\
$W_1$ 
& 0.0064 & 0.0076 & 0.0088 & 0.0155 & 0.0189 & 0.0223 \\
Average degree
& 0.0032 & 0.0040 & 0.0048 & 0.0028 & 0.0035 & 0.0041 \\
\hline
&\multicolumn{3}{c}{$p=0.4,q=0.3$}&\multicolumn{3}{c}{$p=0.3,q=0.4$}
\end{tabular}
% \caption*{\small $p = 0.4$, $q = 0.3$}
% \end{minipage}
% \hfill
% \begin{minipage}{0.48\textwidth}
% \textbf{London}
% \centering
% \begin{tabular}{lccc}
% \hline
% \textbf{representations} & \textbf{Lower} & \textbf{Mean} & \textbf{Upper} \\
% \hline
% $d_{MV}$ true alignment  & 0.0242 & 0.0289 & 0.0335 \\
% $d_{MV}$ with 100\% shuffling        & 0.0178 & 0.0214 & 0.0251 \\
% $d_{MV}$ with GM all to one       & 0.0187 & 0.0225 & 0.0263 \\
% $d_{MV}$ with GM consecutive      & 0.0102 & 0.0129 & 0.0156 \\
% $W_1$                & 0.0155 & 0.0189 & 0.0223 \\
% Average degree          & 0.0028 & 0.0035 & 0.0041 \\
% \hline
% \end{tabular}
% \caption*{\small $p = 0.3$, $q = 0.4$}
% \end{minipage}
\caption{Comparison of MSEs for different distances using the $l_2$ localizer for $n = 200$, $m = 20$, $nmc = 500$ with different $(p, q)$ settings for the London model. Chance level for $m=20$ is $6.8\times 10^{-2}$. Thus all representations have MSE significantly better than chance, meaning they capture information about $t^*$.}
\label{tab:L}
\end{table}

For numerical solution of graph matching,
we use the ``indefinite" method provided by the ``gm" function in the R package igraphmatch \cite{qiao2021igraphmatch}, which relaxes the problem from the set of permutation matrices to its convex hull, then applies the Frank-Wolfe algorithm to find a local optimum for the relaxed problem. We set the max number of iterations = 100 for all simulations and set no seed vertices. This runs the alignment without any known correspondences, starting from a random-chosen doubly stochastic matrix.

These results are summarized in Table \ref{tab:L} for the London model and Table \ref{table:A} for the Atlanta model. 
For the London model, the MSE is significantly smaller than chance, across all choices of distance. The best performing estimator is given by average degree, which is comparable to $d_{MV}$ with true alignment when $p = 0.4, q = 0.3$, but is significantly better when $p = 0.3, q = 0.4$. This appears to be a result of smaller variance for the average degree compared to other choices of distance, as we saw in Figure \ref{fig:London-mirrors}. In this setting, $d_{MV}$ on consecutive pair graph matching and $W_1$ both result in much smaller MSE compared to even the true alignment. In all three cases, utilizing marginal information brings the estimator closer to a function of the sufficient statistics described in Theorem~\ref{thm:London_MLE}, which should result in decreased variance in light of the Rao-Blackwell theorem. In particular, for the London model, ignoring vertex alignment does not hinder inference about $t^*$, and could actually result in improved estimation. 
%While when $p = 0.4, q = 0.3$ it is much worse than the true alignment. 
%Consecutive could make the signal smaller 
%and in former section we claim but that is to $W_2$ which we didn't include because it doesn't have piecewise linear structure, thus not comparable, thus in table $W_1$ and GM performance are different, in 0.4 0.3 case $W_1$ is much better while in 0.3 0.4 they have overlapped error bar. 

On the contrary, for the Atlanta model in either of the $n=500$ or $n=800$ cases, only $\hat{d}^2_{MV}$ on the true alignment has MSE significantly better than chance. Importantly, the MSE for all to one graph matching, consecutive graph matching, or ind-$d_{MV}$ are all comparable, indicating that for the Atlanta model, once the alignment is lost, the information contained within that alignment cannot be recovered. All three methods have MSE worse than chance level, as seen in Figure \ref{fig:Atlanta-MSEvsShuffleRatio}.


\begin{table}[htbp]
\centering
%\begin{minipage}{0.48\textwidth}
\centering
%\textbf{Atlanta}
\begin{tabular}{lccc||ccc}
&&\textbf{Atlanta}&&&\textbf{Atlanta}&\\
\hline
\textbf{Representation} & \textbf{Lower} & \textbf{Mean} & \textbf{Upper} &\textbf{Lower} & \textbf{Mean} & \textbf{Upper} \\
\hline
$d^2_{MV}$ with true alignment  & 0.0163 & 0.0195 & 0.0228 &0.0095 & 0.0112 & 0.0129 \\
$d^2_{MV}$ with 100\% shuffling         & 0.1007 & 0.1082 & 0.1157 &0.0959 & 0.1040 & 0.1120 \\
$d^2_{MV}$ with GM all to one        & 0.0875 & 0.096 & 0.1044 &0.0849 & 0.0929 & 0.1010 \\
$d^2_{MV}$ with GM consecutive       & 0.0904 & 0.0952 & 0.1000 &0.0896 & 0.0941 & 0.0985 \\
$W_1$                                 & 0.0769 & 0.0845 & 0.0920 &0.0533 & 0.0603 & 0.0674 \\
Average degree                          &  0.0690 & 0.0765 & 0.0841 &0.0626 & 0.0697 & 0.0767 \\
\hline
&&$n=500$&&&$n=800$&
\end{tabular}
% \caption*{\(n = 500\)}
% \end{minipage}
% \hfill
% \begin{minipage}{0.48\textwidth}
% \centering
% \textbf{Atlanta}\\
% \begin{tabular}{ccc}
% \hline
% \textbf{Lower} & \textbf{Mean} & \textbf{Upper} \\
% \hline
% 0.0095 & 0.0112 & 0.0129 \\
% 0.0959 & 0.1040 & 0.1120 \\
% 0.0849 & 0.0929 & 0.1010 \\
% 0.0896 & 0.0941 & 0.0985 \\
% 0.0533 & 0.0603 & 0.0674 \\
% 0.0626 & 0.0697 & 0.0767 \\
% \hline
% \end{tabular}
% \caption*{\(n = 800\)}
% \end{minipage}
\caption{Comparison of MSEs for different distances using the $l_2$ localizer for $p = 0.4$, $q = 0.2$, $m = 20$, $N = 50$, and $nmc = 300$ with different $n$ settings for Atlanta model. 
Only the $d^2_{MV}$ dissimilarity with true vertex alignment provides the information to localize the changepoint $t^*$, yielding an MSE substantially below the chance level of $0.0679$. 
All other representations are uninformative and perform no better than chance.}
\label{table:A}
\end{table}

\subsection{Simulated swarm data}\label{sec:Simulations_Swarm}

We consider a time series of graphs derived from the behavior of interacting agents in a swarm, simulated in \cite{hindes2024swarming}. This data is a 1000-seconds-long video consisting of 10,001 sequential frames (1 frame per 0.1 second), representing the movements of 100 agents in a 2D plane with known agent correspondences across all frames.
For simplicity, we focus on a specific subinterval of time, spanning from 780 seconds to 820 seconds, resulting in \( m = 401 \) frames with \( t_1 = 780 \, s \) and \( t_{401} = 820 \, s \).
The video for this time period can be found here \footnote{\url{https://www.cis.jhu.edu/~parky//SofA/NRL-swarm-movie.mp4}}. 

Our goal is to identify the structural changepoint times within this video.
Upon visually reviewing the video, we identify two distinct structural changes: at 800 seconds, the agents form two major groups that gradually merge back into a single larger group; and around 810 seconds, they begin to split into two unequal groups.
To apply our mirror method to detect and localize the changes, we first convert the video into a time series of graphs.

For each frame \( t \), we observe the locations in \( \RR^2 \) of \( n = 100 \) agents.
We use the locations as the latent positions for each node, represented as a matrix \( \mathbf{X}_t \in \RR^{100 \times 2} \) by stacking the locations row-wise.
Next, we generate a RDPG using \( \mathbf{X}_t \) as latent position matrix, thresholding the \( \mathbf{P}_t = \mathbf{X}_t \mathbf{X}^{\top}_t \) matrix such that entries below 0 are set to 0 and entries above 1 are set to 1.
This process yields a time series of undirected graphs without self-loops, characterized by \( n = 100 \) nodes and \( m = 401 \) time points.
We can then implement our \( d_{MV} \)-based mirror method.
In practice, when the mirror exhibits a manifold structure, we can further reduce the dimension of the mirror using Isometric Mapping (Isomap) \cite{isomap_science} on the results obtained from CMDS, which results in more straightforward changepoint estimation.
This results in an \emph{iso-mirror}, the algorithm for which is summarized in Algorithm~\ref{alg:iso-mirror}. We note that in step 3 of the algorithm, to approximate the minimizing rotation matrix $W\in \mathcal{O}^{d_{\text{ASE}}\times d_{\text{ASE}}}$ for $\|\hat{\bX}_t-\hat{\bX}_s W\|_2$, we instead use the closed-form solution to the corresponding problem with the Frobenius norm instead of the spectral norm. This approach can be justified by the bounds in \cite{cape2019two} or computational comparisons provided in \cite{jasa2025procrustes}.
%When applying Isomap, we consider $k$-nearest-neighbor graph, where $k$ is the smallest value such that the graph is connected.
%.
\begin{algorithm}
\caption{Iso-mirror estimation}
\label{alg:iso-mirror}
\begin{algorithmic}[1]
\State Input: TSG $\{\bA_1,\bA_2,\cdots,\bA_m\}$, ASE dimension $d_{\text{ASE}}$, CMDS dimension $d$.
\State Compute $\hat{X}_t=$ASE$(A_t)\in \RR^{n\times d_{\text{ASE}}}, 1\leq t\leq m$.
\State {Construct the distance matrix $\hat{\mathcal{D}}\in\R^{m\times m}$ where $\hat{\mathcal{D}}_{t,s}=\frac{1}{\sqrt{n}}\|\hat{\bX}_{t}-\hat{\bX}_{s}W_F\|_2$,} where $W_F=\arg \min_{W\in\mathcal{O}^{d_{\text{ASE}}\times d_{\text{ASE}}  }} \|\hat{\bX}_{t}-\hat{\bX}_{s}W\|_F$.
\State Compute CMDS$(\hat{\mathcal{D}})=\{\hat{\psi}(1),\hat{\psi}(2),\cdots,\hat{\psi}(m)\}$, with each $\hat{\psi}(t) \in \R^d$. 
\State Apply Isomap on $\{\hat{\psi}(1),\hat{\psi}(2),\cdots,\hat{\psi}(m)\}$ using $k$ nearest neighbors, where $k$ is the smallest value such that the $k$-nearest-neighbor graph is connected, obtaining $\{\hat{\psi}_{d\to1}(1),\hat{\psi}_{d\to1}(2),\cdots,\hat{\psi}_{d\to1}(m)\}$, with each $\hat{\psi}_{d\to 1}(t)\in \RR$.
\State Output: $\{\hat{\psi}_{d\to1}(1),\hat{\psi}_{d\to1}(2),\cdots,\hat{\psi}_{d\to1}(m)\}$.
\end{algorithmic}
\end{algorithm}

For this video-induced TSG, first we apply the iso-mirror Algorithm \ref{alg:iso-mirror} on the true aligned TSG with $d_{\text{ASE}} =2$ and CMDS embedding dimension $d=2$. 
%Note in step 3 we need to solve a 2 dimensional orthogonal Procrustes problem with respect to spectral norm, which doesn't have analytical solution.
%We will use Frobenius orthogonal Procrustes problem solution as an approximate, which by our empirical experiments shows similarly mirror results. 
Second, to investigate the impact of not knowing the true alignment, we randomly shuffle 100\% of the vertices across all time points and then apply the iso-mirror algorithm to the shuffled TSG with $d_{ASE} =2$ and $d =2$.
Finally, we perform the all to one graph matching with max number of iterations $200$ as described in Equation~\ref{eq:alltoonegm}, followed by applying the iso-mirror algorithm to the graph-matched TSG with $d_{ASE} =2$ and $d =2$.
We also consider some metrics that do not make use of alignment information, including the Wasserstein  distances $W_1$ and $W_2$, and the average degree distance.
For the Wasserstein distances, we chose $d_{ASE} = 2$. 
%and we replace the third step in Algorithm \ref{alg:iso-mirror}. 
As in Algorithm~\ref{alg:iso-mirror}, for any pair $s,t\in[m]$, we first Procrustes align $\hat{\bX}_s$ to $\hat{\bX}_t$ with respect to Frobenius norm, then calculate $W_p\left(\hat{\mu}_{\hat{\bX}_{t}} , \hat{\mu}_{\hat{\bX}_{s}}\right)$ on the appropriately-rotated matrices of latent positions to obtain the $s,t$ entry of the distance matrix.
We choose CMDS dimension $d=1$ for both of the Wasserstein distance matrices.

To allow for the simultaneous detection and localization of first-order changepoints on the final 1-dimensional iso-mirrors above, we utilize the function ``selgmented" in R package Segmented \cite{muggeo2017interval,segmented} to perform segmentation regression.
This approach fits a piecewise linear function onto the  representations while accommodating an unknown number of slope changes.
We set the max number of slope changes to be $20$ and select the optimal number of changepoints based on the Bayesian Information Criterion (BIC).
\begin{figure}
    \centering
    \includegraphics[width=0.9\linewidth]{figures/NRL_2.pdf}
    \caption{Gray dots represent the iso-mirror embeddings for the time series of graphs derived from simulated swarm data between 780 seconds and 820 seconds. Black dashed lines indicate two structural changes detected by eye at 800s and 810s. The red lines are the segmentation regression fits for the estimated iso-mirrors. The panels in the left column, arranged from top to bottom, display the iso-mirror results for the true aligned TSG, the TSG with 100\% vertex shuffling, and the all to one graph-matched TSG with max number of iterations $200$. The right column shows the corresponding results for the $W_1,$ $W_2$, and average degree distances. While shuffling completely disrupts the original dynamics, graph matching partially recovers them. }
    \label{fig:NRL}
\end{figure}

These results are shown in Figure \ref{fig:NRL}. We see that the iso-mirror under the true alignment effectively captures the dynamics of the swarm movement. Segment regression detects four slope changes: $787(1.1), 798(0.1),801(0.1), 811(0.3)$, with 3 occurring around time points $800$ and 810—closely aligning with the changepoints identified by visual inspection. 
In contrast, for the shuffled TSG, the iso-mirror becomes zigzagged and uninformative, and no slope changes are detected: the segment regression yields a flat line.
This indicates that the data is sensitive to lost of true correspondence.
After applying all to one graph matching, segmentation regression once again detects two slope changes $798(0.2)$ and $801(0.2)$, but it fails to recover the change around 810. 
When ignoring alignment completely, we see using the mirror induced by $W_1$ or $W_2$ detects more slope changes (6 and 10 respectively) than $d_{MV}$ and its variants, and both detect changes around 800 and 810.
This indicates that the structural changes happening around times 800 and 810 are encoded in the marginal distributions of the TSG, in a way that the $W_p$ distances are more sensitive to than the $d_{MV}$ distance, even with vertex alignment. It is unclear whether these additional detected changepoints are spurious, or just subtler than the ones detected by eye. Finally, using average degree as a simple summary statistic detects 14 slope changes, the most among all methods. This quantity seems to be smoothly varying over time, rather than exhibiting piecewise-linear structure, so the current changepoint detection method may not be well-suited for this input.




\section{Conclusion and discussion}\label{sec:Conclusion}

This paper investigates the impact of vertex misalignment for changepoint detection in time series of networks. We study the effect of this shuffling on the joint distributions of the latent positions between times, as well as methods for comparing the latent position distributions at different times using only marginal information. We propose two tractable models for time series of networks that represent the extreme ends of a spectrum of sensitivity to the vertex alignment: in the London model, marginal information is sufficient, so losing the vertex correspondence does not degrade inferential power. On the other hand, in the Atlanta model, once the vertex alignment is lost, all information about the changepoint is gone, and cannot be recovered through graph matching or marginal-based distances. As we demonstrate in Section~\ref{sec:Simulations_Swarm}, more complicated data-generating processes can lead to phenomena somewhere between these two extremes, where certain changepoints may be detectable through marginal information, or can be recovered through graph matching, while others may be lost without known vertex correspondence. 

Besides serving as examples of the most extreme phenomena for recovery of changepoints with and without vertex alignment, the tractability of the London and Atlanta models lead to new insights about distances between networks and their use with the Euclidean mirror methodology, since these mirrors are explicitly computable. For example, we find that in the London model, $d_{MV}$-based distances look approximately Euclidean 1-realizable regardless of the amount of vertex shuffling, but the Wasserstein metric $W_1$ and the average degree metric are both exactly 1-realizable. Additional phenomena appear for the Atlanta model, where CMDS applied to the matrix of \emph{squared} $d_{MV}$ distances leads to approximate 1-realizability for a large number of states $N$. In the fully shuffled case, the mirror in the Atlanta model becomes completely degenerate. These phenomena were all previously unobserved for Euclidean mirrors generated from LPPTSGs before the present work, and provide a more complete picture of what is captured by different notions of distance on latent position random variables. 


Our takeaway for practitioners is as follows.
If the signal of interest can be found in the marginal distributions, as in the London model, it may be worth considering Wasserstein-based metrics for inference, and the vertex alignment may be non-essential. When the signal is believed to be in the joint distribution, as in the Atlanta model, the alignment can be crucial for accurate recovery. 
The standard $d_{MV}$ mirror is a robust option if a known alignment is available with only a small amount of error.
However, when the vertex alignment is completely unknown, practitioners should use $d_{MV}$ on graph-matched data with caution, since this approach may only capture marginal information (akin to a $W_1$ mirror), while appearing to capture information about the joint distribution. Because it is often unclear in practice whether a signal lies in the marginal or joint, we recommend that practitioners explore a variety of dissimilarity-induced mirrors to ensure a comprehensive analysis. It may also be worth exploring the effect of shuffling a small portion of the vertices to see how sensitive the results are to an errorful vertex alignment.

Some limitations of the current paper include Theorem~\ref{thm:independentdMV_convergence}, which imposes the restrictive assumption of a positive scalar LPP. Second, a principled method for determining whether signal resides in the joint or marginal distributions would make these results more practical for real data. 

As an initial step, we propose computing the ratio $\frac{d_{MV}}{\text{ind-}d_{MV}}$ for all pairs of time points. When this is consistently close to 1, it suggests that shuffling has a minimal impact on the mirror. Larger deviations may indicate sensitivity to shuffling.

Third, the current London and Atlanta models represent two edge cases: one contains signals only in the marginal distributions, while the other has signals only in the joint distribution.
This raises the question of whether we can construct intermediate cases, where the signal is contained in both marginal and joint information, and finding appropriate distances for measuring the dynamics in such cases. For example, neither of the current models can reproduce the phenomena observed in the simulated swarm data example in Section~\ref{sec:Simulations_Swarm}, where the signal is lost after shuffling but is partially recovered after graph matching.
This amounts to finding a LPP whose $W_1$ mirror is partially informative while its ind-$d_{MV}$ is uninformative.

Finally, the use of the ind-$d_{MV}$ and Wasserstein distances motivates us to consider alternative dissimilarities for inducing mirrors.
This raises a key question: how does one choose a dissimilarity in a principled way to best extract information about the dynamics within a class of LPPs? Formally defining this question and solving it represents an interesting and important topic for future study.


\section*{Acknowledgements}

Tianyi Chen was supported by The Acheson J. Duncan Fund for the Advancement of Research in Statistics. Tianyi Chen and Carey E Priebe were supported by Office of Naval Research (ONR) Science of Autonomy award number N00014-24-1-2278. Carey E Priebe was supported by the Office of Naval Research (ONR) Applied Physics Lab (APL) Grant 151530. Vince Lyzinski was supported by Air Force Office of Scientific Research (AFOSR) Complex Networks award number FA9550-25-1-0128 and The Johns Hopkins University HLT COE.

\section{Appendix}

\begin{figure}[htp]
  \centering
  \includegraphics[height=4.5cm]{figures/Illustration/100london_sample_paths_p=0.3_q=0.9_m=30_c=0.1_tstar=0.5.pdf}%
  \hspace{-7mm} 
 % \includegraphics[height=4.8cm]{figures/Illustration/adj100london_graphs_p=0.3_q=0.9_m=30_c=0.1_tstar=0.5.pdf}
  \caption{Sample paths from the London model with $p = 0.3$, $q = 0.9$, $m = 30$ and $t^* = 0.5$, $c_L = 0.1$, $\delta_m = 0.9/30$. Each state is represented by a dot and states are connected for each path. After the changepoint at $t^{*}=0.5$, the density of connections in the adjacency matrices increase more rapidly, indicating a structural shift both in sample paths and network connectivity.}
  \label{fig:L-illustration}
\end{figure}


\begin{figure}[htp]
  \centering
  \includegraphics[height=4.5cm]{figures/Illustration/100atlanta_sample_paths_p=0.05_q=0.45_m=30_c=0.8_N=50_tstar=0.5.pdf}
  \hspace{-7mm}
  %\includegraphics[height=4.8cm]{figures/Illustration/adj100atlanta_graphs_p=0.05_q=0.45_m=30_c=0.8_N=50_tstar=0.5.pdf}
  \caption{Sample paths for the Atlanta model with $p = 0.05$, $q = 0.45$, $m = 30$ and $t^* = 0.5$, $c_A = 0.8$, $N = 50$, $\delta_N=\frac{0.8}{49}$. Each state is represented by a dot and states are connected for each path. After the changepoint at $t^{*}=0.5$, the sample paths show increased oscillation because the jump probability increases from $0.05$ to $0.45$. However, the density of connections in the corresponding adjacency matrices remains stable over time.}
 \label{fig:A-illustration}
\end{figure}




\subsection{Additional simulation results}

First we compare 
$\psi^A_{\alpha-d_{MV}}$ and 
$\psi^A_{\alpha-d^2_{MV}}$.
Note as shown in Theorem \ref{thm:Atlanta-main},
$\psi^A_{\alpha-d_{MV}}$ is a scaling of $ \psi^A_{d_{MV}}$. 
But it is no longer the case for $\psi^A_{\alpha-d^2_{MV}}$:
$$
\mathcal{D^A}_{\alpha-d^2_{MV}} 
=\mathcal{D^A}_{\alpha-d_{MV}}^{(2)} 
=   \alpha \mathcal{D^A}_{\text{ind-}d_{MV}}^{(2)} +(1-\alpha) \mathcal{D^A}_{d_{MV}}^{(2)}.
$$
Recall CMDS on $\mathcal{D^A}_{\alpha-d^2_{MV}}$  amounts to considering 
\begin{align*}
\mathcal{D^A}^{(2)}_{\alpha-d^2_{MV}} 
&=  
\left(\alpha \mathcal{D^A}_{\text{ind-}d_{MV}}^{(2)} +(1-\alpha) \mathcal{D^A}_{d_{MV}}^{(2)}
\right)^2\\
&= \alpha \mathcal{D^A}_{\text{ind-}d^2_{MV}}^{(2)} + (1-\alpha) \mathcal{D^A}_{d^2_{MV}}^{(2)} + 2\alpha(1-\alpha) \mathcal{D^A}_{\text{ind-}d_{MV}}^{(2)} \odot \mathcal{D^A}_{d_{MV}}^{(2)}, 
\end{align*}
where $\odot$ denote the entrywise product of two matrices. 
Because of the third term $\mathcal{D^A}_{\text{ind-}d_{MV}}^{(2)} \odot \mathcal{D^A}_{d_{MV}}^{(2)}$, $\psi^A_{\alpha-d^2_{MV}}$ is no longer a scaling of $\psi^A_{d^2_{MV}}$, $\psi^A_{d_{MV}}$ or $\psi^A_{\alpha-d_{MV}}$. 
However, as shown in Figure \ref{fig:atlanta-iso-mirror_0.2}, they are similar to each other.

Then we compare the iso-mirror result as in Algorithm \ref{alg:iso-mirror} with $d_{MV}$, not $d^2_{MV}$. 
We choose CMDS embedding dimension $d$ to be 3, 5, 8, and 10. 
Interestingly, we can see that as $d$ gets larger, the iso-mirror appears more piecewise linear, but when $d$ is too high, the mirror becomes highly variable and exhibits false changepoints. 
Meanwhile, if we focus on the estimated iso mirror (the blue dots), we see using $d=3$ or $d = 5$ adds more noise.
Higher shuffling percentage adds more variance to the blue dots while the true mirror remains unchanged. 
To summarize, the iso-mirror successfully represents the dynamics when there is no noise, but when $n$ is finite, there is a bias-variance tradeoff, and the optimal choice of CMDS dimension prior to applying ISOMAP may be greater than 1. 

% Note we decide in the main text for simulation of MSE we use $d^2_{MV}$ only for simplicity.
% Here we also include the same simulation setting and use iso mirror with $d_{MV}$ for Atlanta here.
% As in Figure \ref{fig:atlanta_mse_vs_shuffle_appendix}(a)(c)(e) is the MSE versus shuffle ratio for Atlanta model using $d_{MV}$ or $d^2_{MV}$ or iso mirror with $d=10$ and the localizer is $l_2$.
% And summarized in Table \ref{tab:A-appendix} is the MSE for Atlanta model with using $d_{MV}$ or $d^2_{MV}$ or iso mirror with $d=4$ or $d=8$ on shuffled TSG or graph matched TSG.
% We also include the result with $l_{\infty}$ localizer as in Figure \ref{fig:atlanta_mse_vs_shuffle_appendix}(b)(d)(f).
% We also show the MSE versus shuffle ratio for London model using $d_{MV}$ with both localizer \ref{fig:london_mse_vs_shuffle_appendix}. 

%However behavior is more complicated when we consider the $\alpha-$shuffled-$d_{MV}$ distance matrix.
%First note that since it is a convex combination of $\mathcal{D}_{d_{MV}}$ and $\mathcal{D}_{\text{shuffled-}d_{MV}}$ thus it doesn't have the property of $\mathcal{D}_{d_{MV}}$ , that is entrywise square won't make the CMDS result linear.
%However from empirical results we show that ISOMAP with higher dimension of CMDS does help to make it looks linear. But you can’t use a too high dimension in the CMDS embedding! As the figures show, if you use too high dimension, it will look weird and zigzagging! But if you look at the blue dots the $n = 1000$ case, ISOMAP add more variance compare to the linearity it adds. Thus use isomap with high CMDS embedding is not a good idea when n is small!
%Also note how shuffling percent increase(5\% to 20\%) adds more variance to the blue dots while black dots remain unchanged. To summarize, ISOMAP does extract the linearity and slope changepoint in the dynamics successfully and serves as(at least one of) the BEST representation when there is no noise.
%When $n$ is finite, it is another bias variance trade off story. Thus we decide in the main text for simulation of MSE we use $d^2_{MV}$ only for simplicity. But we did also run isomap result below. 


\begin{figure}
    \centering
    \includegraphics[width=0.8\linewidth]{figures/Illustration/n=1000m=30N=50_shuffling_0.2.pdf}
    \caption{The mirror and iso-mirror(with CMDS embedding dimension $d$ being 3 ,5, 8, and 10) with different CMDS embedding dimension
    based on the same realized 20\% shuffled-TSG from Atlanta model in Figure \ref{fig:Atlanta-mirrors} with $n = 1000$, $p = 0.05$, $q = 0.45$, $m = 30$ and $t^* = 0.5$, $c_A = 0.8$, $N = 50$.
    }
    \label{fig:atlanta-iso-mirror_0.2}
\end{figure}

\begin{figure}
    \centering
    \includegraphics[width=0.8\linewidth]{figures/Illustration/n=1000m=30N=50_shuffling_0.05.pdf}
    \caption{he mirror and iso-mirror with different CMDS embedding dimension based on the same realized 5\%-shuffled-TSG from Atlanta model in Figure \ref{fig:Atlanta-mirrors} with $n = 1000$, $p = 0.05$, $q = 0.45$, $m = 30$ and $t^* = 0.5$, $c_A = 0.8$, $N = 50$. }
    \label{fig:atlanta-iso-mirror_0.05}
\end{figure}



%\begin{figure}
%    \centering
%    \includegraphics[width=0.5\linewidth]{figures/Simulation results/London-both-localizer n = 200 nmc = 500 p = 0.4 q = 0.3.pdf}
%    \caption{This figure is London model. With $n=200$, $p=0.3$ and $q = 0.4$ This figure shows that $l_{\infty}$ performs worse than $l_2$ localizer in almost all cases. Meanwhile $W_2$ performs worst, which is expected since because $\psi_{W2}$ is not piecewise linear so not really comparable with others.}
%    \label{fig:sim_London_app1}
%\end{figure}


%For Atlanta model, there is another way to obtain approximate piecewise linear other than $d^2_{MV}$ as in the paper.
%We can CMDS $d_{MV}$ distance matrix into high dimension $c = 4, 8$ and then apply ISOMAP to the CMDS results to embed into 1 dimension, and our empirical results shows that the ISOMAP to 1d result is approximate piecewise linear again, thus here we show the MSE of $\hat{t}$ on using this iso-mirror.  
\begin{figure}[htbp]
  \centering

  % Row 1
  \begin{subfigure}[b]{0.48\textwidth}
    \includegraphics[width=\textwidth]{figures/Simulation_results/Atlanta_D_l2.pdf}
    \caption{$l_2$ localizer on $d_{MV}$ mirror}
  \end{subfigure}
  \hfill\begin{subfigure}[b]{0.48\textwidth}
    \includegraphics[width=\textwidth]{figures/Simulation_results/Atlanta_d_linf.pdf}
    \caption{$l_\infty$ localizer on $d_{MV}$ mirror}
  \end{subfigure}
  
  \vspace{1em}

  % Row 3
  \begin{subfigure}[b]{0.48\textwidth}
    \includegraphics[width=\textwidth]{figures/Simulation_results/Atlanta_D2_l2.pdf}
    \caption{$l_2$ localizer on $d_{MV}^2$ mirror}
  \end{subfigure}
  \hfill
  \begin{subfigure}[b]{0.48\textwidth}
    \includegraphics[width=\textwidth]{figures/Simulation_results/Atlanta_D2_linf.pdf}
    \caption{$l_\infty$ localizer on $d_{MV}^2$ mirror}
  \end{subfigure}

  

  \vspace{1em}
  % Row 2
  \begin{subfigure}[b]{0.48\textwidth}
    \includegraphics[width=\textwidth]{figures/Simulation_results/Atlanta_iso_l2.pdf}
    \caption{$l_2$ localizer on iso mirror with $d = 10$}
  \end{subfigure}
  \hfill
  \begin{subfigure}[b]{0.48\textwidth}
    \includegraphics[width=\textwidth]{figures/Simulation_results/Atlanta_iso_linf.pdf}    
    \caption{$l_\infty$ localizer on iso mirror with $d = 10$}
  \end{subfigure}

  \caption{MSE versus shuffle ratio $\alpha$ under the Atlanta model for various combinations of mirror distances and localization methods. Each panel shows results using either the $d_{MV}$, $d_{MV}^2$, or iso mirror (with CMDS embedding dimension $d = 10$), paired with an $l_2$ or $l_\infty$ localizer. For fixed  $p = 0.4$ and each $q \in \{0.1, 0.2, 0.35, 0.4, 0.5\}$, results are averaged over $nmc = 100$ Monte Carlo replicates. Vertical bars denote 1.96 standard deviations, and the red dotted line marks the chance MSE level of 0.075 for $m = 40$. Across all configurations, none of the methods successfully detect the true changepoint location under increasing vertex misalignment. MSE under the $l_2$ localizer often lies above the chance level, primarily due to bias toward boundary estimates, as the procedure tends to overfit noise near the endpoints of the time series.
  }
  \label{fig:atlanta_mse_vs_shuffle_appendix}
\end{figure}

\begin{figure}[htbp]
  \centering

  % Row 1
  \begin{subfigure}[b]{0.48\textwidth}
    \includegraphics[width=\textwidth]{figures/Simulation_results/London_l2.pdf}
    \caption{$l_2$ localizer on $d_{MV}$ mirror}
  \end{subfigure}
  \hfill\begin{subfigure}[b]{0.48\textwidth}
    \includegraphics[width=\textwidth]{figures/Simulation_results/London_Linf.pdf}
    \caption{$l_\infty$ localizer on $d_{MV}$ mirror}
  \end{subfigure}

  \caption{
    MSE versus proportion of shuffled vertices $\alpha$ under the London model for two variants of the mirror localization procedure. The plots show performance using the $d_{MV}$ mirror combined with either an $l_2$ or $l_\infty$ localizer. For $p = 0.4$ and each $q \in \{0.1, 0.2, 0.35, 0.4, 0.5\}$, results are averaged over $nmc = 100$ Monte Carlo replicates. Error bars represent 1.96 standard deviations, and the red dotted line denotes the chance MSE level of 0.075 for $m = 40$. Across all shuffle ratios, the London model demonstrates robust changepoint recovery, with MSE consistently below chance and stable across increasing levels of vertex misalignment.
    Both localization methods exhibit strong performance, and no significant difference is observed between the $l_2$ and $l_\infty$ localizers.}
  \label{fig:london_mse_vs_shuffle_appendix}
\end{figure}



\begin{sidewaystable}[htbp]
\centering
\begin{tabular}{cc}
\begin{subtable}[t]{0.48\textwidth}
\centering
\begin{tabular}{lccc}
\hline
\textbf{Metric} & \textbf{Lower} & \textbf{Mean} & \textbf{Upper} \\
\hline
$d_{MV}$ with true alignment & 0.0333 & 0.0402 & 0.0471 \\
$d_{MV}$ with shuffled        & 0.1051 & 0.1128 & 0.1206 \\
$d_{MV}$ with GM all to one       & 0.0857 & 0.0940 & 0.1024 \\
$d_{MV}$ with GM consecutive      & 0.1447 & 0.1513 & 0.1578 \\
\hline
\end{tabular} 
\caption{c = 4, $n = 500$}
\end{subtable}
&
\begin{subtable}[t]{0.48\textwidth}
\centering
\begin{tabular}{lccc}
\hline
\textbf{Metric} & \textbf{Lower} & \textbf{Mean} & \textbf{Upper} \\
\hline
$d_{MV}$ with true alignment & 0.0087 & 0.0123 & 0.0159 \\
$d_{MV}$ with shuffled        & 0.1043 & 0.1121 & 0.1200 \\
$d_{MV}$ with GM all to one       & 0.0941 & 0.1021 & 0.1101 \\
$d_{MV}$ with GM consecutive      & 0.1558 & 0.1615 & 0.1672 \\
\hline
\end{tabular}
\caption{c = 4, $n = 800$}
\end{subtable}
\\[12ex]
\begin{subtable}[t]{0.48\textwidth}
\centering
\begin{tabular}{lccc}
\hline
\textbf{Metric} & \textbf{Lower} & \textbf{Mean} & \textbf{Upper} \\
\hline
$d_{MV}$ with true alignment & 0.0326 & 0.0389 & 0.0452 \\
$d_{MV}$ with shuffled        & 0.0999 & 0.1074 & 0.1150 \\
$d_{MV}$ with GM all to one       & 0.0790 & 0.0873 & 0.0955 \\
$d_{MV}$ with GM consecutive      & 0.0647 & 0.0725 & 0.0803 \\
\hline
\end{tabular}
\caption{c = 8, $n = 500$}
\end{subtable}
&
\begin{subtable}[t]{0.48\textwidth}
\centering
\begin{tabular}{lccc}
\hline
\textbf{Metric} & \textbf{Lower} & \textbf{Mean} & \textbf{Upper} \\
\hline
$d_{MV}$ with true alignment & 0.0285 & 0.0346 & 0.0407 \\
$d_{MV}$ with shuffled        & 0.0984 & 0.1063 & 0.1142 \\
$d_{MV}$ with GM all to one       & 0.0824 & 0.0906 & 0.0988 \\
$d_{MV}$ with GM consecutive      & 0.0675 & 0.0757 & 0.0838 \\
\hline
\end{tabular}
\caption{c = 8, $n = 800$}
\end{subtable}
\end{tabular}

\vspace{3mm}
\caption{
Comparison of MSEs across metrics using $l_2$ localizer for $p = 0.4$, $q = 0.2$, $m = 20$, $N = 50$, and $nmc = 500$ with different $n$ for the Atlanta model.
Here, we apply ISOMAP to CMDS $d_{MV}$ distance matrices with embedding dimension $c = 4, 8$, and use the resulting 1 dimensional embedding for localization.
As with $d_{MV}^2$, localization using $d_{MV}$ with true alignment remains the only method consistently below the chance level $0.068$. All shuffled methods, including those with graph matching, fail to recover $t^*$. 
Both choices of higher CMDS dimension, $c=4$ and $c=8$, fail to perform better than the $d^2_{MV}$ mirror. 
%confirming that in the Atlanta model, once alignment is shuffled, the dynamic structure provides no usable information for time localization.
}
\label{tab:A-appendix}
\end{sidewaystable}






% \newpage


% \subsection{Proof of Lemma \ref{thm:independentdMV_convergence}}

% \begin{proof}

% \item[1.] let $C_i$ be the event that $\sigma$ fixes the index $i$. 
% We note that $\PP(C_{i})=
% \frac{(n-1)!}{n!}=\frac{1}{n}.$
% Therefore, since $N^{\sigma}_{fix}=\sum_{i=1}^n I_{C_i}$, $\EE[N^{\sigma}_{\text{fix}}] =\sum_{i=1}^n \PP(C_i) =1$.  


% \item[2.] Let \(\sigma\) be a shuffling. We want to show that \(\sigma^{-1}\) is also a shuffling. For any permutation \(\pi \in S_n\),
% \[
%     \PP\bigl(\sigma^{-1} = \pi\bigr) = \PP\bigl(\sigma = \pi^{-1}\bigr) = \frac{1}{n!}\,.
% \]
% This shows that \(\sigma^{-1}\) is uniformly distributed over \(S_n\).

% \item[3.] For any $\pi \in S_n$,
% $$
% \PP(\sigma_0 \circ \sigma = \pi ) = \PP( \sigma = \sigma_0^{-1} \circ \pi ) = \frac{1}{n!};\quad 
% \PP(\sigma \circ \sigma_0 = \pi ) = \PP( \sigma = \pi \circ \sigma_0^{-1} ) = \frac{1}{n!}
% $$
% Thus $\sigma_0 \circ \sigma$ is uniformly distributed over $S_n$ and $\sigma \circ \sigma_0$ is uniformly distributed over $S_n$ . 

% \item[4.] For two independent shufflings \(\sigma_1\) and \(\sigma_2\), we compute the probability of their composition:
% \begin{align*}
%     \PP\bigl(\sigma_1 \circ \sigma_2 = \pi\bigr) &= \sum_{\tau \in S_n} \PP(\sigma_2 = \tau) \PP\bigl(\sigma_1 \circ \sigma_2 = \pi\mid \sigma_2 = \tau \bigr) \quad (\text{law of total probability}) \\
%     &= \sum_{\tau \in S_n} \PP(\sigma_2 = \tau) \PP\bigl(\sigma_1 = \pi \circ \tau^{-1}\bigr) \quad (\sigma_1 \ind \sigma_2) \\
%     &= \sum_{\tau \in S_n} \frac{1}{n!} \cdot \frac{1}{n!} \quad (\text{both } \sigma_1 \text{ and } \sigma_2 \text{ are uniformly distributed}) \\
%     &= \frac{1}{n!}.
% \end{align*}
% This shows that \(\sigma_1 \circ \sigma_2\) is also uniformly distributed over \(S_n\).
% \item[5.] Assume $\sigma_1$ and $\sigma_2$ are two independent $\alpha-$shufflings, denote their corresponding permutation matrix as $P_1$ and $P_2$. Then $P_1 = I_{n-\alpha_n} \oplus P^1_{\alpha_n}$, $P_2 = I_{n-\alpha_n} \oplus P^2_{\alpha_n}$. Note $P_1 P_2 = I_{n-\alpha_n} \oplus P^1_{\alpha_n} P^2_{\alpha_n}$. Also note $P^1_{\alpha_n} P^2_{\alpha_n}$ corresponds to the composite of two independently shufflings on $S_{\alpha_n}$ and from above we know  $P^1_{\alpha_n} P^2_{\alpha_n}$ is also a shuffling on $S_{\alpha}$. Thus $\sigma_1 \circ \sigma_2$ is also a $\alpha-$shuffling.


% \end{proof}


% \newpage

\subsection{Proof of Theorem \ref{thm:independentdMV_convergence}}
We begin with an elementary bound that addresses the impact of derangement; we include it for reader reference.

\begin{proposition}\label{prop:derangement-2d}
For any $n>2$, denote $\sigma$ as a fixed derangement in $S_n$, that is $\forall i \in [n], \sigma(i) \neq i$. Consider two random variables in $X,Y\in\RR^1$ with a joint measure $\mu_{X,Y}$ on $\RR^2$. For a function $g: \RR^d \to \RR$, denote $\EE_{\mu}[g] : =\EE_{(x_1,x_2,\cdots,x_d) \sim \mu}[g(x_1,x_2,\cdots,x_d)]$.
Denote $\{(x_i,y_i)^n_{i=1}\}$ as $n$ i.i.d samples from $\mu_{X,Y}$ defined on the probability space $((\RR^2)^{\times n}, (\mathcal{B}(\RR^2))^{\times n}) = : (\Omega, \mathcal{B} ).$ The measure on $\Omega$ is denoted as $\PP_{\Omega}$. Denote $
\{ \left(x_i, y _{\sigma(i)} \right); i \in [n]\}$ as the sample after the derangement. Then we have for any function bounded on support of $\mathcal{S}: = \text{support of } \mu_X \otimes \mu_Y $, $f: \RR^2 \to \RR$, denote 
$
B_{f,\mathcal{S}} = \sup_{(x,y) \in \mathcal{S}} |f(x,y)|,
$
and for any $\epsilon > 0$, 
%there exists two constant $c'_{f,\mathcal{S}}, c''_f$ that only depends on $f$ and $\mathcal{S}$ such that
$$
\PP_{\Omega} \left( \left| \frac{1}{n} \sum^n_{i=1} f\left(x_i, y_{\sigma(i)}\right)  -  \EE_{ \mu_X \otimes \mu_Y}[f]   \right| \ge \epsilon \right) \le \frac{ 3 B_{f,\mathcal{S}}^2 }{n \epsilon ^2}.
$$
\end{proposition}

% We will use the following lemma:

% \begin{lemma}\label{lem:ind}
% Assume $X \in \RR^d$ and $Y \in \RR^{d'}$ are two random variables, and $f: \RR^d \to \RR^c$, $g : \RR^{d'} \to \RR^{c'}$ are two measurable functions, then $X \ind Y$ implies that $f(X) \ind g(Y)$. 
% \end{lemma}
% % \begin{proof}
% % Assume $B_1 \subset \RR^c$ and $B_2 \subset \RR^{c'}$ are any two Borel sets, denote $f^{-1}(C)$ as the preimage of $C$ under $f$ for any $f$ and any $C$. 
% % \begin{align*}
% % \PP\left( f(X) \in B_1 , g(Y) \in B_2   \right) &= \PP\left( X \in f^{-1}(B_1) , Y \in g^{-1} (B_2)   \right) \\
% % & = \PP\left( X \in f^{-1}(B_1)\right) \PP \left( Y \in g^{-1} ( B_2)   \right) \\
% % & =  \PP\left( f(X) \in B_1 \right) \PP \left( g(Y) \in  B_2   \right).
% % \end{align*}
% % Thus $f(X) \ind g(Y)$. 
% % \end{proof}
% Now we prove the theorem.
\begin{proof}
First we note for any $i$ recall $\sigma(i) \neq i$, then $(x_i, y_i) \ind (x_{\sigma(i)}, y_{\sigma(i)}) $. Now consider two projections $\Pi_x(x,y) = x$ and $\Pi_y(x, y)= y$. By independence, $\Pi_x(x_i, y_i) \ind \Pi_y(x_{\sigma(i)}, y_{\sigma(i)})$, that is, $x_i \ind y_{\sigma(i)}$. Further $x_i \sim \mu_X$, $y_{\sigma(i)} \sim \mu_Y$, thus for any $i$, $(x_i, y_{\sigma(i)}) \sim \mu_X \otimes \mu_Y $. Then 
$$
\EE \biggl[ \frac{1}{n} \sum^n_{i=1} f\left(x_i, y_{\sigma(i)}\right) \biggr] = \frac{1}{n} \sum^n_{i=1} \EE \biggl[ f\left(x_i, y_{\sigma(i)}\right) \biggr] =  \EE_{ \mu_X \otimes \mu_Y}[f].
$$
Next we show the variance is shrinking with $n$. Note
$$
\Var\biggl[ \frac{1}{n} \sum^n_{i=1} f\left(x_i, y_{\sigma(i)}\right) \biggr] = \frac{1}{n^2} \left( \sum^n_{i=1} \Var \biggr[ f\left(x_i, y_{\sigma(i)}\right) \biggr]  + \sum^n_{k=1} \sum^n_{t \neq k } \Cov \biggl[f\left(x_k, y_{\sigma(k)}\right) , f\left(x_t, y_{\sigma(t)}\right)   \biggr] \right),
$$
Note $f$ is bounded on $\mathcal{S}$, i.e., the support of $\mu_X \otimes \mu_Y$, so there exists a $c'_{f,\mathcal{S}}$ such that for all $k \in [n]$ we have
$$\Var [f\left(x_k, y_{\sigma(k)}\right)] = \Var_{(x,y)\sim\mu_X \otimes \mu_Y} [f(x,y)] \leq \EE_{(x,y)\sim\mu_X \otimes \mu_Y} [f(x,y)^2] \leq B_{f,\mathcal{S}}^2.$$ 
Thus the first term is bounded:
$$
\sum^n_{i=1} \Var \biggr[ f\left(x_i, y_{\sigma(i)}\right) \biggr] \leq n B_{f,\mathcal{S}}^2.
$$
Now consider the second term. For a given $k \in [n]$, there are at most two terms in the summand that are nonzero, which are both bounded by a constant. Given the derangement $\sigma$ and a fixed $k$, denote $t_1 :=\sigma(k)$; $t_2:=\sigma^{-1}(k)$. As long as $t \neq k, t_1, t_2$, then $\{k,\sigma(k)\}\cap\{t,\sigma(t)\}=\varnothing$ and 
$$
f\left(x_k, y_{\sigma(k)}\right) \ind f\left(x_t, y_{\sigma(t)}\right) \Rightarrow \Cov \biggl[f\left(x_k, y_{\sigma(k)}\right) , f\left(x_t, y_{\sigma(t)}\right)\biggr] = 0.
$$
For the at most two nonzero terms in the sum, Cauchy-Schwarz yields
$$\mathrm{Cov}\left[f(x_k,y_{\sigma(k)}),f(x_t,y_{\sigma(t)})\right]\leq \sqrt{\mathrm{Var}\left[f(x_k,y_{\sigma(k)})\right]\mathrm{Var}\left[f(x_t,y_{\sigma(t)})\right]}\leq B_{f,\mathcal{S}}^2,$$
so
$$
\sum^n_{k=1} \sum_{t \neq k } \Cov \biggl[f\left(x_k, y_{\sigma(k)}\right) , f\left(x_t, y_{\sigma(t)}\right) \biggr]\leq 2nB_{f,\mathcal{S}}^2.
$$
Thus
$$
\Var \biggl[ \frac{1}{n} \sum^n_{i=1} f\left( x_{i} - y_{\sigma(i)} \right) \biggl] \leq \frac{3B_{f,\mathcal{S}}^2}{n}.
$$
and the result now follows from Chebyshev's inequality.

\end{proof}

We now prove Theorem~\ref{thm:independentdMV_convergence}. 
\begin{proof}
% The proof idea is a classic good set bad set argument. We will show that if we randomly draw a permutation $\sigma$, as long as $N_{\text{fix}}$ is small, $A_F$ is small and $\frac{N_{\text{fix}}}{n}\EE_{ \mu_X \otimes \mu_Y}$ is small, thus for event $A$ to happen $A_D$ must be big, using Theorem \ref{thm:derangement-2d} we already know probability of $|A_D| > \epsilon ~|~ \sigma$ is small, thus $\PP(A)$ is controllable. Meanwhile we know the probability of a random permutation's $N_{\text{fix}}$ to be big is small. 


Denote $\PP_{S_n}$ as the uniform distribution on $S_n$.
Because $\sigma$ is independent of the samples, the probability measure on $\Omega \times S_n$ is the product measure $\PP_{\Omega} \otimes \PP_{S_n}$, denoted as $\PP_{\Omega\times S_n}$.
Let $A$ denote the event
$$
A  :=\left\{\omega \times \sigma \in \Omega \times S_n: \left| \frac{1}{n} \sum^n_{i=1} f\left(x_i, y_{\sigma(i)}\right)  -  \EE_{ \mu_X \otimes \mu_Y}[f]   \right|> \epsilon\right\}. 
$$ Note $A\subseteq \Omega \times S_n $. We control $\PP_{\Omega\times S_n}(A)$. 
For any permutation $\sigma$, we divide $[n]$ into two disjoint sets based on the points fixed by $\sigma$, namely $\{i : \sigma(i) = i\}$ and its complement $\{i : \sigma(i) \neq i\}$. Denote the cardinality of the first set by $N^{\sigma}_{\text{fix}}$. We deduce that 
\begin{align*}
\frac{1}{n} &\sum^n_{i=1} f\left(x_i, y_{\sigma(i)}\right)  -  \EE_{ \mu_X \otimes \mu_Y}[f]\\
&= 
\underbrace{\frac{1}{n} \sum_{i:i=\sigma(i)}  f\left(x_i, y_{\sigma(i)}\right)}_{=:Y_n^{\sigma}} - \frac{N^{\sigma}_{\text{fix}}}{n}\EE_{ \mu_X \otimes \mu_Y}\\
&\quad +  \underbrace{\frac{n-N^{\sigma}_{\text{fix}}}{n} \left(\frac{1}{n-N^{\sigma}_{\text{fix}}} \sum_{i:i \neq \sigma(i)}  f\left(x_i, y_{\sigma(i)}\right) - \EE_{ \mu_X \otimes \mu_Y}  \right) }_{=:W_n^{\sigma}}.
\end{align*}
where $Y_n^{\sigma}$ and $W_n^{\sigma}$ are random variables defined above. 

Assume $f$ is upper bounded by $B_{f,\mathcal{S}}>\frac{\epsilon}{4}$ on $\mathcal{S}$. Denote $\alpha :  = \frac{\epsilon}{4B_{f,\mathcal{S}}}$, $ 0< \alpha <1 $. 
Denote by $E$ the event
$$
E : = \Omega \times \{\sigma \in S_n:N^{\sigma}_{\text{fix}} \leq \alpha n \}. $$
%Intuition of this $\alpha$ is we can tolerate $\alpha n$ number of fixed points and $N_{\text{fix}} \leq \alpha n$ is what we meant as small.  
Since $|\sup_{(x,y) \in \mathcal{S}} f(x,y)| \leq B_{f,\mathcal{S}}$, we know that
$$
|Y_n^{\sigma}| \leq \frac{N^{\sigma}_{\text{fix}}B_{f,\mathcal{S}}}{n},
$$
and hence on $E$, we have that 
$|Y_n^{\sigma}| \leq \alpha B_{f,\mathcal{S}}=\frac{\epsilon}{4}$ and
$$\bigg|\frac{N^{\sigma}_{\text{fix}}}{n} \EE_{ \mu_X \otimes \mu_Y}[f]\bigg| \leq \frac{N^{\sigma}_{\text{fix}}}{n} \EE_{ \mu_X \otimes \mu_Y}[|f|] \leq  \frac{B_{f,\mathcal{S}} N^{\sigma}_{\text{fix}}}{n} \leq \alpha B_{f,\mathcal{S}}=\frac{\epsilon}{4}.
$$
Now, on $A$, we see that   
$|Y_n^{\sigma}+W_n^{\sigma} - \frac{N^{\sigma}_{\text{fix}}}{n} \EE_{ \mu_X \otimes \mu_Y}[f] | > \epsilon
$, so 
$$
|W_n^{\sigma} | > \epsilon - |Y_n^{\sigma}| - \bigg|\frac{N^{\sigma}_{\text{fix}}}{n} \EE_{ \mu_X \otimes \mu_Y}[f]\bigg|,
$$ 
Thus, on $E$, we see that $|Y_n^{\sigma}| \leq \frac{\epsilon}{4}$ and $\bigg|\frac{N^{\sigma}_{\text{fix}}}{n} \EE_{ \mu_X \otimes \mu_Y}[f]\bigg| \leq \frac{\epsilon}{4} $, and therefore on $A \cap E$, we have $|W_n^{\sigma}| > \frac{\epsilon}{2}$.
If we define $A'$ by
$$
A' := \left\{ \omega \times \sigma \in \Omega \times S_n:  |W_n^{\sigma}| > \frac{\epsilon}{2} \right\}
$$
then we conclude from the above that
$A \cap E \subseteq A' \cap E.$
We bound $\PP_{\Omega \times S_n}(A)$ by noting that 
\begin{align*}
\PP_{\Omega \times S_n}(A) = & \PP_{\Omega \times S_n}(A\cap E) + \PP_{\Omega \times S_n}(A \cap E^c) \\
\le & \PP_{\Omega \times S_n}(A' \cap E) + \PP_{\Omega \times S_n}(E^c)
\end{align*}
Consider the first term:
\[
\begin{aligned}
&\PP_{\Omega \times S_n}(A' \cap E) \\
& =\sum_{\sigma_0: N^{\sigma_0}_{\text{fix}} \leq \alpha n }  
\PP_{\Omega}\!\left(|W_n^{\sigma_0}|>\frac{\epsilon}{2}\right) \, \PP_{S_n}(\sigma_0) \quad %(\text{property of product measure} )
\\
& = \sum_{\sigma_0: N^{\sigma_0}_{\text{fix}} \leq \alpha n }  \PP_{\Omega}\left( \bigg|\frac{n-N^{\sigma_0}_{\text{fix}}}{n} \left(\frac{1}{n-N^{\sigma_0}_{\text{fix}}} \sum_{i:i \neq \sigma_0(i)}  f\left(x_i, y_{\sigma_0(i)}\right) - \EE_{ \mu_X \otimes \mu_Y}[f]  \right) \bigg|> \frac{\epsilon}{2}  ~  \right)\PP_{S_n}(\sigma_0) \\
&=   \sum_{\sigma_0: N^{\sigma_0}_{\text{fix}} \leq \alpha n }  \PP_{\Omega}\left( \bigg| \frac{1}{n-N^{\sigma_0}_{\text{fix}}} \sum_{i:i \neq \sigma_0(i)}  f\left(x_i, y_{\sigma_0(i)}\right) - \EE_{ \mu_X \otimes \mu_Y}[f] \bigg|>\frac{n}{n-N^{\sigma_0}_{\text{fix}}} \frac{\epsilon}{2}  \right) \PP_{S_n}(\sigma_0) \\
& \leq \sum_{\sigma_0: N^{\sigma_0}_{\text{fix}} \leq \alpha n } 
\frac{ (12B_{f,\mathcal{S}}^2)(n-N^{\sigma_0}_{\text{fix}})^2 }{(n - N^{\sigma_0}_{\text{fix}} ) n^2 \epsilon ^2}  \PP_{S_n}(\sigma_0) \quad (\text{Proposition }\ref{prop:derangement-2d}) \\
& = \sum_{\sigma_0: N^{\sigma_0}_{\text{fix}} \leq \alpha n }  \frac{ ( 12B_{f,\mathcal{S}}^2)(n-N^{\sigma_0}_{\text{fix}}) }{n^2 \epsilon ^2} \PP_{S_n}(\sigma_0) \\
& \leq \sum_{\sigma_0: N^{\sigma_0}_{\text{fix}} \leq \alpha n } \frac{  12B_{f,\mathcal{S}}^2}{n \epsilon ^2} \PP_{S_n}(\sigma_0)\\
& = \frac{  12B_{f,\mathcal{S}}^2}{n \epsilon ^2} \PP_{S_n}\left( N^{\sigma_0}_{\text{fix}} \leq \alpha n \right)  \\
& \le \frac{  12B_{f,\mathcal{S}}^2}{n \epsilon ^2}.
\end{aligned}
\]


An easy computation shows $\EE_{S_n}[N^{\sigma}_{\text{fix}}] = 1$, so we bound $\PP(E^c)$ by Markov's inequality:
$$
\PP_{\Omega \times S_n}(E^c ) = \PP_{S_n}(N^{\sigma}_{\text{fix}} > \alpha n) \leq \frac{1}{\alpha n}.
$$
Combining these completes the proof.
\end{proof}





\subsection{Proof of Corollary \ref{Cor:partial shuffling}}
Let $A \subseteq \Omega \times S_{\alpha_n}$ denote the event
$$
\resizebox{\textwidth}{!}{$\displaystyle A : = \{ \omega \times \sigma_{\alpha_n} \in \Omega \times S_{\alpha_n} , \left| \hat{d}^2_{MV}(\mathbf{X}^{t_i}, P_{\alpha_n} \mathbf{X}^{t_j})-  
\left( \left(1-\frac{\alpha_n}{n}\right) d_{MV}(X_{t_i},X_{t_j}) +  \frac{\alpha_n}{n}\, \text{ind-}d^2_{MV}(X_{t_i}, X_{t_j})\right)
\right| \geq \epsilon  \}.$}
$$
We derive that 
\begin{align*}
\hat{d}^2_{MV}(\mathbf{X}_{t}, P_{n,\alpha} \mathbf{X}_{t'})  = & \frac{1}{n}\| \bX_t - P_{\alpha_n}  \bX_{t'}\|^2_F   \\
= & \frac{1}{n}\| \bX_t - (I_{n - \alpha_n} \oplus P_{\alpha_n} )   \bX_{t'}\|^2_F \\
= & \frac{1}{n} \left(  \sum^{n - \alpha_n}_{i=1} \left( \left( \bX_t \right)_i - (\bX_{t'})_i \right)^2  + \sum_{i = n - \alpha_n + 1}^n \left( \left( \bX_t \right)_i - (\bX_{t'})_{\sigma^{\alpha_n}(i)} \right)^2     \right) \\
= & \frac{n - \alpha_n}{n} ~ \frac{1}{n - \alpha_n}  \sum^{n - \alpha_n}_{i=1} \left( \left( \bX_t \right)_i - (\bX_{t'})_i \right)^2  + \frac{\alpha_n}{n} \frac{1}{\alpha_n}\sum_{i = n - \alpha_n + 1}^n \left( \left( \bX_t \right)_i - (\bX_{t'})_{\sigma^{\alpha_n}(i)} \right)^2. 
\end{align*}

Recall that $\{\left(\left( \bX_t \right)_i - (\bX_{t'})_i \right)^2 \} ^{n - \alpha_n}_{i=1}$ are i.i.d samples with mean $\EE[(X_t - X_{t'})^2] = d^2_{MV}(X_{t_i}, X_{t_j})$ and variance $\Var[(X_{t}-X_{t'})^2]$. Further, $X_t,X_{t'}$ are both supported on a subset of $[0,1]$, so $\Var[(X_{t}-X_{t'})^2] \leq 1$. Chebyshev's inequality guarantees that
\begin{align*}
& \PP_{\Omega \times S_{\alpha_n}} \left( \underbrace{  \bigg| \frac{1}{n - \alpha_n}  \sum^{n - \alpha_n}_{i=1} \left( \left( \bX_t \right)_i - (\bX_{t'})_i \right)^2  - d^2_{MV}(X_{t_i}, X_{t_j}) \bigg| \ge \epsilon }_{ = : B} \right) \\
= & \PP_{\Omega} \left(  \bigg| \frac{1}{n - \alpha_n}  \sum^{n - \alpha_n}_{i=1} \left( \left( \bX_t \right)_i - (\bX_{t'})_i \right)^2  - d^2_{MV}(X_{t_i}, X_{t_j}) \bigg| \ge \epsilon  \right)  \\
\le &  \frac{1}{(n-\alpha_n) \epsilon^2}.    
\end{align*}
For the second term we apply Theorem \ref{thm:independentdMV_convergence} with $f(x,y)=(x-y)^2$ and conclude that

$$
\PP_{\Omega \times S_{\alpha_n}} \left(  \underbrace{ \bigg| \frac{1}{ \alpha_n}  \sum^{n }_{i = n - \alpha_n +1 } \left( \left( \bX_t \right)_i - (\bX_{t'})_{\sigma^{\alpha_n}(i)} \right)^2  - \text{ind-}d^2_{MV}(X_{t_i}, X_{t_j}) \bigg| \ge \epsilon }_{ = : C} \right)  
\le   \frac{12}{\alpha_n \epsilon^2} + \frac{4}{\alpha_n \epsilon}  .  
$$
Finally, $A \subseteq B \cup C$, and hence 
$$
\PP_{\Omega \times S_{\alpha_n}} (A) \le \PP_{\Omega \times S_{\alpha_n}}( B \cup C) \leq \PP_{\Omega \times S_{\alpha_n}}(B) + \PP_{\Omega \times S_{\alpha_n}}(C) \leq 
\frac{1}{(n-\alpha_n) \epsilon^2} + \frac{12}{\alpha_n \epsilon^2} + \frac{4}{\alpha_n \epsilon}. 
$$

%\subsection{Proof of Theorem \ref{thm:W1_continuous}}



\subsection{Proof of Theorem \ref{thm:London_main}}

The first bullet point in the unshuffled case is a property of $\tilde{\psi}^1_{d_{MV}}$, and it is proved in \cite{chen2024euclidean} Theorem 2.26. 
For the completely shuffled case $\alpha=1$, we calculate the ind-$d_{MV}$ for $X^L_{t_i}$ and $X^L_{t_j}$ from London Model:

When $t_i,t_j < t^*$, note $X^L_{t_i} \sim \frac{1}{m} \text{Bin}(i,p)$, $X^L_{t_j} \sim \frac{1}{m} \text{Bin}(j,p)$. 
\begin{align*}
\text{ind-}d_{MV}(X^L_{t_i},X^L_{t_j})^2 & =  \EE[(X^L_{t_i}-X^L_{t_j})^2] \quad \text{ ($w=1$ is the minimizer.)}  \\
& =  \Var[X^L_{t_i}-X^L_{t_j}]   + \left(\EE[X^L_{t_i}]-\EE[X^L_{t_j}]\right)^2 \\
& = \Var[X^L_{t_i}]+\Var[X^L_{t_j}] + \left(\EE[X^L_{t_i}]-\EE[X^L_{t_j}]\right)^2  \quad (X^L_{t_i} \ind X^L_{t_j})\\
& = \frac{1}{m^2}  (i+j)(p-p^2) + \frac{1}{m^2}\left(p(i-j)\right)^2. 
\end{align*}
Similarly, when $t_i < t^* < t_j$, $X^L_{t_i} \sim \frac{1}{m} \text{Bin}(i,p)$, $X^L_{t_j} = B_j+B_2$ with $ B_j\sim \frac{1}{m}\text{Bin}(j - t_m^*,q), B_2 \sim \frac{1}{m}\text{Bin}(t_m^*,p), B_j \ind B_2$.
\begin{align*}
\text{ind-}d_{MV}(X^L_{t_i},X^L_{t_j})^2
& = \Var[X^L_{t_i}]+\Var[X^L_{t_j}] + \left(\EE[X^L_{t_i}]-\EE[X^L_{t_j}]\right)^2  \\
& = \Var[X^L_{t_i}]+\Var[B_j] + \Var[B_2]+ \left(\EE[X^L_{t_i}]-\EE[X^L_{t_j}]\right)^2  \\
& = \frac{1}{m^2}\left( \left(p-p^2\right)i + \left(p-p^2\right)t^*_m + (j-t^*_m)(q-q^2) + \left(\left( j-t^*_m   \right)q +\left(t^*_m - i \right)p\right)^2   \right).
\end{align*}
Finally, when $t_i, t_j > t^*$, $X^L_{t_i} = B_i + B_2, X^L_{t_j} = B_j + B_2$, then 
\begin{align*}
\text{ind-}d_{MV}(X^L_{t_i},X^L_{t_j})^2
& = \Var[X^L_{t_i}]+\Var[X^L_{t_j}] + \left(\EE[X^L_{t_i}]-\EE[X^L_{t_j}]\right)^2  \\
& = \Var[B_i]+\Var[B_2]+\Var[B_j] + \Var[B_2]+ \left(\EE[X^L_{t_i}]-\EE[X^L_{t_j}]\right)^2  \\
& = \frac{1}{m^2}\left( \left(q-q^2\right)(i-t^*_m) + 2\left(p-p^2\right)t^*_m + (j-t^*_m)(q-q^2) + q(i-j)^2  \right) \\
& = \frac{1}{m^2}\left(  (q-q^2)(i+j)+2t^*_m(p-p^2-q+q^2)   + q(i-j)^2  \right).
\end{align*}
The above show the form of $\mathcal{D^L}^{(2)}_{\text{ind-}d_{MV}}$. Note that for any $m$, for all $i,j\in[m]$,
$$
(\mathcal{D^L}^{(2)}_{\text{ind-}d_{MV}})_{i,j} - (\mathcal{D}^{(2)}_Z)_{i,j} \sim O\left(\frac{1}{m}\right).
$$
And in the proof of Theorem 2.26 in \cite{chen2024euclidean}, all other properties are consequences of this property. Thus the same proof can go through for $\mathcal{D^L}^{(2)}_{\text{ind-}d_{MV}}$ too. Thus $\tilde{\psi}^L_{\text{ind-}d_{MV}}$ has the same property. 

For $0<\alpha<1$, we combine these two calculations to obtain for all $i,j\in[m]$
$$
(\mathcal{D^L}^{(2)}_{\alpha-d_{MV}})_{i,j} - (\mathcal{D}^{(2)}_Z)_{i,j} = \alpha (\mathcal{D^L}^{(2)}_{\text{ind-}d_{MV}})_{i,j} + (1-\alpha)(\mathcal{D^L}^{(2)}_{d_{MV}})_{i,j} - (\mathcal{D}_Z^{(2)})_{i,j} \sim O\left(\frac{1}{m}\right).
$$
By the same argument as above we have the result. 

For the second bullet point, we observe that whenever $s\leq t$, $X_s$ is stochastically dominated by $X_t$, since $\PP[X_s\leq X_t]=1$, so by \cite{de20211}, 
$$
W_1(X^L_{t_i},X^L_{t_j}) = |\EE [X^L_{t_i}] - \EE[X^L_{t_j}]|.
$$

For the third bullet point, consider a graph $\bA$ with $n$ vertices that is generated from a random dot product model with the distribution of random variable $X$. That is, latent positions $\{x_1, x_2, ... x_n\}$ are drawn i.i.d from $X \in \R^1$, then each edge is a Bernoulli sampled with probability as inner product of latent positions. Then the average degree is 
\begin{align*}
\EE \biggl[ \frac{1}{n} \sum^n_{i =1} \sum^n_{j \neq i} \bA_{i,j} \biggr] & = 
\EE \biggl[  \EE \biggl[ \frac{1}{n} \sum^n_{i =1} \sum^n_{j \neq i} \bA_{i,j} | \{x_1, x_2, ... x_n\}  \biggr] \biggr]  \\
& =
\EE \biggl[   \frac{1}{n} \sum^n_{i =1} \sum^n_{j \neq i} \EE[\bA_{i,j} | \{x_1, x_2, ... x_n\}]   \biggr] \\ 
& = 
\EE \biggl[   \frac{1}{n} \sum^n_{i =1} \sum^n_{j \neq i} x_i x_j  \biggr]  \\
& = 
\frac{1}{n} \sum^n_{i =1} \sum^n_{j \neq i} \EE[x_i] \EE[x_j]  \\
& = (n-1) \left(\EE[X]\right)^2.
\end{align*}
For the London model, assuming $c=0$,$t^*_m/m = t^*$,
$$
\EE[X^L_{t_i}] = \begin{cases}
p t_i & t_i \le t^*, \\
p \frac{t^*_m}{m} + q(t_i - \frac{t^*_m}{m}) = p t^* + q(t_i - t^*) & t_i > t^*.
\end{cases}
$$
\begin{align*}
d_{\mathrm{deg}}(X_t,X_s)&=(n-1)|\EE[X_t]^2-\EE[X_s]^2|\\
    &=(n-1)|\psi_Z(t)^2-\psi_Z(s)^2|\\
    &=|\psi_{\mathrm{deg}}(t)-\psi_{\mathrm{deg}}(s)|
\end{align*}
when we set $\psi_{\mathrm{deg}}(t)=(n-1)\psi_Z(t)^2$.

\subsection{Proof of Theorem~\ref{thm:London_MLE}}

Consider $n$ sample trajectories from the London model, $\{X_i(t):i\in[n], t\in[m]\}$. The likelihood is given by
\begin{align*}
    \mathrm{lik}(\{X_i(t)\}_{i,t}|p,q,t^*) &=\prod_{i=1}^n\prod_{t=1}^{t^*-1}p^{\mathbf{1}_{>0}(X_i(t)-X_i(t-1))}(1-p)^{1-\mathbf{1}_{>0}(X_i(t)-X_i(t-1))}\\
    &\qquad \times\prod_{t=t^*}^m q^{\mathbf{1}_{>0}(X_i(t)-X_i(t-1))}(1-q)^{\mathbf{1}_{>0}(X_i(t)-X_i(t-1))}\\
    &=\prod_{t=1}^{t^*-1}p^{\sum_i \mathbf{1}_{>0}(X_i(t)-X_{i}(t-1))}(1-p)^{n-\sum_i \mathbf{1}_{>0}(X_i(t)-X_i(t-1))}\\
    &\qquad\times\prod_{t=t^*}^m q^{\sum_i \mathbf{1}_{>0}(X_i(t)-X_i(t-1))}(1-q)^{n-\sum_i \mathbf{1}_{>0}(X_i(t)-X_i(t-1))}.
\end{align*}
To complete the proof of sufficiency, we observe that
$$\sum_{i=1}^n \mathbf{1}_{>0}(X_i(t)-X_i(t-1))=\sum_{i=1}^n (X_i(t)-X_i(t-1))/\delta=\sum_{j=0}^m j(c_t(j)-c_{t-1}(j))/\delta.$$

Now we consider maximizing the likelihood. If we fix the supposed changepoint at time $t$, the log-likelihood becomes
\begin{align*}
\ell(p,q)&= \sum_{s=1}^{t-1} \sum_{i=1}^n \left(\mathbf{1}_{>0}(X_i(s)-X_i(s-1))\log\left(\frac{p}{1-p}\right)\right)+n(t-1)\log(1-p)\\
&\qquad +\sum_{s=t}^m \sum_{i=1}^n\left(\mathbf{1}_{>0}(X_i(s)-X_i(s-1))\log\left(\frac{q}{1-q}\right)\right)+n(m-t+1)\log(1-q).
\end{align*}
For a given $i$, we have
\begin{align*}
\sum_{s=1}^{t-1}\mathbf{1}_{>0}(X_i(s)-X_i(s-1))&=\sum_{s=1}^{t-1}(X_i(s)-X_i(s-1))/\delta=(X_i(t-1)-c)/\delta,\\
\sum_{s=t}^m \mathbf{1}_{>0}(X_i(s)-X_i(s-1))&=\sum_{s=t}^m (X_i(s)-X_i(s-1))/\delta=(X_i(m)-X_i(t-1))/\delta.
\end{align*}
So summing over $i$ gives
\begin{align*}
\ell(p,q)&= \sum_{j=0}^{t-1}jc_{t-1}(j)\log\left(\frac{p}{1-p}\right)+n(t-1)\log(1-p)\\
&\qquad +\sum_{j=0}^m j(c_m(j)-c_{t-1}(j))\log\left(\frac{q}{1-q}\right)+n(m-t+1)\log(1-q).
\end{align*}
The remaining optimization now follows the standard argument for the MLE of a Bernoulli distribution.

\subsection{Proof of Theorem \ref{thm:Atlanta-main}}

We want to determine the behavior of $d_{MV}^2(X_s^A,X_t^A)$ when $N$ is large. First we show $d_{MV}^2$ can be written as a trace. Let $T_p$ denote the transition probability matrix for the Markov chain defining one step of the Atlanta model when the transition probability is $p$, and let $M=[(i-j)^2]_{i,j=1}^N$. 
Since all random variables in Atlanta model are positive supported, we know $w=1.$ Assume $i<j<t^*_m$:
\begin{align*}
d^2_{MV}\left(X^A_{t_i} , X^A_{t_j}\right) & = \EE[  \left(X^A_{t_i} - X^A_{t_j}\right)^2 ] \\
&=\sum_{n=1}^N\sum_{l=1}^N\PP\bigl( X_{t_i}^A=S_n, X_{t_j}^A=S_l\bigr)(S_l-S_n)^2\\
& =  \sum^N_{n=1} \frac{1}{N}\sum^N_{l=1} \PP\left(X^A_{t_j} = S_l | X^A_{t_i} = S_n \right)  \left(S_l - S_n \right)^2 (X^A_{t_i}\text{ is uniform}) \\
&= \frac{1}{N} \sum^N_{n=1}\sum^N_{l=1} \left(T^{j-i}_{p}\right)_{n,l}(\delta_N(n-l))^2 \\
& = \frac{\delta^2_N}{N} \tr\left(T_p^{j-i} M \right).
\end{align*}
When $i<t^*_m<j$, 
$$
\PP(X^A_{t_i} = S_l | X^A_{t_j} = S_n ) = \left(T_p^{t_m^*-i}T_q^{j-t_m^*}\right)_{n,l},
$$ and the case with $t_m^*<i<j$ looks like the first case but with $T_q$ replacing $T_p$. Now we simplify the trace expression. 
We know from \cite{kouachi2006eigenvalues} the eigenvalues of $T_p$ are: 
\[
\boxed{
\lambda_k^{(p)}
\;=\;
1 \;-\; 2p
\;+\;
2p\;\cos\Bigl(\frac{(k-1)\,\pi}{N}\Bigr),
\quad
k \;= \;1,\;\dots,\;N.
}
\]
and corresponding eigenvectors are:

$$
\boxed{
(v_k)_j = \cos\left( \frac{ (k-1) \pi}{2 N} (2j-1) \right) , k=1,2,\cdots ,N , \quad j = 1,2, \cdots, N.}
$$

Also note $$
||v_k||^2 = \sum_{j=1}^N \cos ^2\left(\frac{k \pi(2 j-1)}{2 N}\right)=\frac{N}{2}, \quad \text { for } \quad k \neq 0
$$ then we have the orthonormal eigenvectors $u_k = \sqrt{\frac{2}{N}}v_k$ for $k>1$ and $u_1 = \sqrt{\frac{1}{N}}v_1$. Note that $u_k$s are not functions of $p$ and are common eigenvectors of all $T_p$ matrices.
\begin{lemma}
    Sum of cosines and sines of an arithmetic sequence.
    \[
\sum_{k=0}^{n-1} \cos(a + k \cdot d) = \frac{\sin\left(\frac{n \cdot d}{2}\right)}{\sin\left(\frac{d}{2}\right)} \cdot \cos\left(\frac{2a + (n-1) \cdot d}{2}\right)
\]
\[
\sum_{k=0}^{n-1} \sin(a + k \cdot d) = \frac{\sin\left(\frac{n \cdot d}{2}\right)}{\sin\left(\frac{d}{2}\right)} \cdot \sin\left(\frac{2a + (n-1) \cdot d}{2}\right)
\]
\end{lemma}
\begin{proof}
Note that
\[
\sum_{k=0}^{n-1} e^{(a + k \cdot d)i} = e^{ai}\frac{e^{ndi} - 1}{e^{di}-1}= e^{ai}\frac{e^{ndi/2} - e^{-ndi/2}}{e^{di/2}-e^{-di/2}}\frac{e^{ndi/2}}{e^{di/2}} =  \frac{\sin\left(\frac{n \cdot d}{2}\right)}{\sin\left(\frac{d}{2}\right)}e^{ai + ndi/2-di/2}
\]
By taking the real part of the above equation we get
\[
\sum_{k=0}^{n-1} \cos(a + k \cdot d) = \frac{\sin\left(\frac{n \cdot d}{2}\right)}{\sin\left(\frac{d}{2}\right)} \cdot \cos\left(\frac{2a + (n-1) \cdot d}{2}\right),
\]
and by taking the imaginary part we get
\[
\sum_{k=0}^{n-1} \sin(a + k \cdot d) = \frac{\sin\left(\frac{n \cdot d}{2}\right)}{\sin\left(\frac{d}{2}\right)} \cdot \sin\left(\frac{2a + (n-1) \cdot d}{2}\right).
\]
\end{proof}
\begin{lemma}
    Define $\alpha_k = u_k^\top M u_k$, then with $\omega_k=(k-1)\pi/2,$
    $$
\alpha_k = 
\begin{cases}
    \frac{N(N+1)(N-1)}{6} &\quad k = 1\\
    -\frac{\cos^2\left(\omega_k/N\right)}{N\sin^4\left(\omega_k/N\right)}&\quad \text{even } k\\
    0 &\quad \text{odd } k>1\\
\end{cases}.
$$

\end{lemma}
\begin{proof}
    For $k = 1$, 
    \[
    \alpha_1 = u_1^\top M u_1 = \frac{1}{N} \sum_{i=1}^N\sum_{j=1}^N M_{ij} = \frac{1}{N} \sum_{i=1}^N\sum_{j=1}^N (i-j)^2 = \frac{1}{N} \sum_{i=1}^N\sum_{j=1}^N (i^2-2ij+j^2)
    \]
    \[ 
    =\frac{1}{N} \left(\frac{N(N+1)(2N+1)}{6}\cdot N - 2\left(\frac{N(N+1)}{2}\right)^2+N\cdot \frac{N(N+1)(2N+1)}{6}\right)
    \]
    \[ 
    =\frac{N(N+1)(N-1)}{6}
    \]
    If $k>1$, $\left(v_k\right)_i = \cos\left(\frac{\omega_k}{N}(2i-1)\right)$
    \[
    \alpha_k = u_k^\top M u_k = \frac{2}{N} \sum_{i=1}^N\sum_{j=1}^N \left(v_{k}\right)_i\,M_{ij}\, \left(v_{k}\right)_j = \frac{2}{N} \sum_{i=1}^N\sum_{j=1}^N \cos\left( \frac{\omega_k}{N} (2i-1) \right)(i-j)^2\cos\left( \frac{\omega_k}{N} (2j-1) \right)
    \]
    \[= \frac{2}{N} \left(\sum_{i=1}^N i^2 \cos\left( \frac{\omega_k}{N} (2i-1) \right)\right)\left(\sum_{j=1}^N\cos\left( \frac{\omega_k}{N} (2j-1) \right)\right)
    \]
    \[- \frac{2}{N} \cdot 2\left(\sum_{i=1}^N i \cos\left( \frac{\omega_k}{N} (2i-1) \right)\right)\left(\sum_{j=1}^N j\cos\left( \frac{\omega_k}{N} (2j-1) \right)\right)
    \]
    \[+ \frac{2}{N} \left(\sum_{i=1}^N \cos\left( \frac{\omega_k}{N} (2i-1) \right)\right)\left(\sum_{j=1}^N j^2 \cos\left( \frac{\omega_k}{N} (2j-1) \right)\right)
    \]
    Let $f_d(k) = \sum_{i=1}^N i^d \cos\left(\frac{\omega_k}{N}(2i-1)\right)$, then we have \[
    \alpha_k = \frac{2}{N}\left(f_2(k)\cdot f_0(k)-2f_1(k)^2 + f_0(k)f_2(k)\right).
    \]
    Note that $f_0(k) = \sum_{i=1}^N \cos\left(\frac{\omega_k}{N}(2i-1)\right)$, which is a sum of cosines of an arithmetic sequence, so
    \[
    f_0(k) = \frac{\sin\left(\omega_k\right)}{\sin\left(\omega_k/N\right)}\cdot\cos\left(\omega_k\right).
    \]
    Therefore, for integer $k$ we have $f_0(k) = 0$.
    Now $\alpha_k$ simplifies to $-\frac{4}{N}f_1(k)^2$. In order to find $f_1(k)$, first, we calculate sum of sines of the same arithmetic sequence
    \[
    \sum_{i=1}^N \sin\left(\frac{\omega_k}{N}(2i-1)\right) = \frac{\sin\left(\omega_k\right)}{\sin\left(\omega_k/N\right)}\cdot\sin\left(\omega_k\right) = \frac{\sin^2\left(\omega_k\right)}{\sin\left(\omega_k/N\right)}.
    \]
    Now if we take derivative of both sides with respect to $\frac{\omega_k}{N}$, we get
    \[
    \sum_{i=1}^N \left(2i-1\right)\cdot\cos\left(\frac{\omega_k}{N}(2i-1)\right)  = \frac{2N \cos\left(\omega_k\right) \sin\left(\omega_k\right)}{\sin\left(\omega_k/N\right)} - \frac{\cos\left(\omega_k/N\right) \sin^{2}\left(\omega_k\right)}{\sin^{2}\left(\omega_k/N\right)}.
    \]
    Therefore,
    \[
    f_1(k) = \frac{N \cos\left(\omega_k\right) \sin\left(\omega_k\right)}{\sin\left(\omega_k/N\right)} - \frac{\cos\left(\omega_k/N\right) \sin^{2}\left(\omega_k\right)}{2\sin^{2}\left(\omega_k/N\right)}+ \frac{f_0(k)}{2}
    \]
    which for integer $k$ simplifies to
    \[
    f_1(k) = - \frac{\cos\left(\omega_k/N\right) \sin^{2}\left(\omega_k\right)}{2\sin^{2}\left(\omega_k/N\right)}
    \]
    If $k$ is odd, then $\sin^{2}\left(\omega_k\right) = 0$, and if $k$ is even, $\sin^{2}\left(\omega_k\right) = 1$, so we have 
    \[
f_1(k) = 
\begin{cases}
    - \frac{\cos\left(\omega_k/N\right)}{2\sin^{2}\left(\omega_k/N\right)}&\quad \text{even } k\\
    0 &\quad \text{odd } k
\end{cases}
\]
Recall that $\alpha_k = -\frac{4}{N}f_1(k)^2$, therefore, for $k > 1$ we have 
\[
\alpha_k = 
\begin{cases}
    -\frac{\cos^2\left(\omega_k/N\right)}{N\sin^4\left(\omega_k/N\right)}&\quad \text{even } k\\
    0 &\quad \text{odd },
\end{cases}
\]
or in general,
\[
\alpha_k = 
\begin{cases}
    \frac{N(N+1)(N-1)}{6} &\quad k = 1\\
    -\frac{\cos^2\left(\omega_k/N\right)}{N\sin^4\left(\omega_k/N\right)}&\quad \text{even } k\\
    0 &\quad \text{odd } k>1\\
\end{cases}
\]
\end{proof}
\begin{lemma}
    For non-negative integer $k < N$, we have 
    \[
\mathrm{tr }(T_p^k\,M) =  2(N-1)kp -4\sum_{t=2}^k\frac{(-1)^t}{t-1}\binom{k}{t} \binom{2t-4}{t-2}p^t
\] 
\end{lemma}
\begin{proof}
    Let $T_p = U\Sigma_pU^\top$ be the eigendecomposition of $T_p$. Note that since the eigenvectors do not depend on $p$, we do not index the matrix $U$ with $p$. 
    \[
\mathrm{tr }(T_p^k\,M)
= \mathrm{tr} \left(  U \Sigma_p^k U^{\top} M \right)
= \mathrm{tr} \left(   \Sigma_p^k U^{\top} M U \right)
= \sum_{m=1}^{N} \alpha_m\,\bigl(\lambda_m^{(p)}\bigr)^k \text{ with  } \alpha_m = u^{\top}_m M u_m.
\] 
Before simplifying the last summation, note that $\sum_{m = 1}^N \alpha_m = \tr\left(M\right) = 0$, we will use this later. Now plugging in the values for $\alpha_m$ and $\lambda_m^{(p)}$ from the previous lemmas,
\begin{align*}
\mathrm{tr }(T_p^k\,M) &= \sum_{m=1}^{N} \alpha_m\,\bigl(\lambda_m^{(p)}\bigr)^k 
= \alpha_1 +\sum_{\text{even }m} \alpha_m\,\bigl(\lambda_m^{(p)}\bigr)^k\\
&= \alpha_1 -\sum_{\text{even }m} \frac{\cos^2\left(\omega_m/N\right)}{N\sin^4\left(\omega_m/N\right)}\,\left(1-2p+2p\cos\left(\frac{(m-1)\pi}{N}\right)\right)^k\\
&= \alpha_1 -\sum_{\text{even }m} \frac{\cos^2\left(\omega_m/N\right)}{N\sin^4\left(\omega_m/N\right)}\,\left(1-4p\sin^2\left(\frac{\omega_m}{N}\right)\right)^k \\
&=\alpha_1 - \sum_{\text{even }m} \frac{\cos^2\left(\omega_m/N\right)}{N\sin^4\left(\omega_m/N\right)} + \frac{4kp}{N}\sum_{\text{even }m} \frac{\cos^2\left(\omega_m/N\right)}{\sin^2\left(\omega_m/N\right)}\\
&\quad -\sum_{t=2}^k\binom{k}{t}\frac{(-4p)^t}{N}\sum_{\text{even }m}\cos^2\left(\omega_m/N\right)\sin^{2t-4}\left(\omega_m/N\right)
\end{align*}
Note that 
\[
\alpha_1 - \sum_{\text{even }m} \frac{\cos^2\left(\omega_m/N\right)}{N\sin^4\left(\omega_m/N\right)} = \sum_{m=1}^N\alpha_m = 0,
\]
Now we calculate the third term's summation,  we follow the idea of Cauchy's proof of the Basel problem.  

\[
\frac{\cos\left(Nx\right) + i\sin\left(Nx\right)}{\sin^{N}\left(x\right)} = \frac{\left(\cos\left(x\right) + i\sin\left(x\right)\right)^{N}}{\sin^{N}\left(x\right)} =\left( \cot(x) + i\right)^{N}=\sum_{j = 0}^{N}\binom{N}{j}\cot^{N-j}\left(x\right) i^{j}
\]
By taking the real part of the above equations we get
\[
\frac{\cos\left(Nx\right)}{\sin^{N}\left(x\right)} = \sum_{j = 0}^{\lfloor N/2\rfloor}\binom{N}{2j}\cot^{N-2j}\left(x\right) (-1)^{j}
\]
Note that for even $m\leq N$, $\cos\left(N\left(\omega_m/N\right)\right) = \cos\left((m-1)\pi\right/2) = 0$ and $\sin\left(\omega_m/N\right)\neq 0$, so $\omega_m/N$ is a root of the left hand side of the above equation. If $N$ is an even number, then $\cot^2\left(\omega_m/N\right)$ are roots of the polynomial
\[
P(t) = \sum_{j = 0}^{N/2}\binom{N}{2j}(-1)^{j} t ^{N/2 - j}.
\]
All of the numbers $\cot^2\left(\omega_m/N\right)$ for even $0<m\leq N$ are distinct and there are $N/2$ of them. Therefore, they are all of the roots of the above polynomial. From the Vieta formula for the sum of the roots of a polynomial, we get
\[
\sum_{\text{even } m}\cot^2\left(\omega_m/N\right) =  -\frac{-\binom{N}{2}}{\binom{N}{0}}  = \frac{N(N-1)}{2}
\]
For odd $N$, we have 
\[
\frac{\cos\left(Nx\right)}{\cot\left(x\right)\sin^{N}\left(x\right)} = \sum_{j = 0}^{ (N-1)/2}\binom{N}{2j}\cot^{N-1-2j}\left(x\right) (-1)^{j},
\]
and $\omega_m/N$ is again a root of left hand side, therefore, $\cot^2\left(\omega_m/N\right)$ are roots of the polynomial
\[
P(t) = \sum_{j = 0}^{(N-1)/2}\binom{N}{2j}(-1)^{j} t ^{(N-1)/2 - j}.
\]
As $P$ is a $(N-1)/2$ degree polynomial, $\cot^2\left(\omega_m/N\right)$ are all the roots of $P$, and again by the Vieta formula
\[
\sum_{\text{even } m}\cot^2\left(\omega_m/N\right)=  -\frac{-\binom{N}{2}}{\binom{N}{0}}  = \frac{N(N-1)}{2}.
\]
Therefore, in general,
\[
\sum_{\text{even }m} \frac{\cos^2\left(\omega_m/N\right)}{\sin^2\left(\omega_m/N\right)} = \sum_{\text{even } m} \cot^2\left(\omega_m/N\right) = \frac{N(N-1)}{2}
\]

Now we show that for $N > k\geq t$
\[
\sum_{\text{even }m}\cos^2\left(\omega_m/N\right)\sin^{2t}\left(\omega_m/N\right) = c_t N,
\]
where \[
c_t = \frac{1}{2^{2t+2}(t+1)} \binom{2t}{t}
\]
First, note that $\sin^{2t}\left(x\right)$ is an even function with period $2\pi$, so we can write a Fourier series using $\cos\left(rx\right)$ functions for nonnegative integer $r$. 
\[
\sin^{2t}\left(x\right) = \sum_{r= 0 }^\infty A_{r,t} \cos\left(rx\right),
\]
where,
\newcommand{\diff}{\mathop{}\!\mathrm{d}}
\[
A_{0,t} = \frac{1}{2\pi}\int_0^{2\pi} \sin^{2t}\left(x\right)\diff x
\]
and for $r > 0$
\[
A_{r,t} = \frac{1}{\pi}\int_0^{2\pi}\sin^{2t}\left(x\right)\cos\left(rx\right)\diff x.
\]
Let $z = e^{ix}$, then
\[
I_{r,t} = \int_0^{2\pi}\sin^{2t}\left(x\right)\cos\left(rx\right)\diff x = \int_{|z| = 1}\left(\frac{z-z^{-1}}{2i}\right)^{2t}\left(\frac{z^r+z^{-r}}{2}\right)\frac{1}{iz}\diff z 
\]
\[
= \frac{1}{2^{2t+1}i\,(-1)^{t}}\int_{|z| = 1}\frac{\left(z^2-1\right)^{2t}\left(z^{2r}+1\right)}{z^{2t+r+1}}\diff z .
\]
The final integral has a singularity at $z = 0$, therefore, $I_{r,t}$ simplifies to 
\[
I_{r,t} = \frac{1}{2^{2t+1}i\,(-1)^{t}}\,\,2\pi i \,\,\mathop{\mathrm{Res}}_{z = 0}\left(\frac{\left(z^2-1\right)^{2t}\left(z^{2r}+1\right)}{z^{2t+r+1}}\right) =  \frac{\pi}{2^{2t}\,(-1)^t} \,\,\mathop{\mathrm{Res}}_{z = 0}\left(\frac{\left(z^2-1\right)^{2t}\left(z^{2r}+1\right)}{z^{2t+r+1}}\right)
\]
Note that the residue term is equal to the coefficient of $z^{-1}$ in the Laurent series of the integrand around the singularity. Which here is equal to coefficient of $z^{2t+r}$ in expansion of $\left(z^2-1\right)^{2t}\left(z^{2r}+1\right)$. 
\[
\left(z^2-1\right)^{2t}\left(z^{2r}+1\right) = \sum_{j = 0}^{2t}\binom{2t}{j} z^{2j+2r}\left(-1\right)^{2t-j}+\sum_{j = 0}^{2t}\binom{2t}{j} z^{2j}\left(-1\right)^{2t-j}
\]
\[
= \sum_{j = r}^{2t+r}\binom{2t}{j-r} \left(-1\right)^{j-r}z^{2j}+\sum_{j = 0}^{2t}\binom{2t}{j} \left(-1\right)^{j}z^{2j}
\]
Now note that if $r$ is odd, $2t+r$ is odd, but the above series has no non-zero odd power of $z$. Therefore, the residue term would be zero, and $I_{r,t} = 0$ for odd $r$. For even $r$, the terms where $j = t + r/2$ are the desired ones in both summations. If $r > 2t$, we have $t + r/2 > 2t$ and $t + r/2 < r$, so there is not a $z^{2t+r}$ in any of the summations. But if $r \leq 2t$, the coefficient would be
\[
\binom{2t}{t-r/2} \left(-1\right)^{t -r/2}+\binom{2t}{t + r/2} \left(-1\right)^{t + r/2} = 2\binom{2t}{t-r/2} \left(-1\right)^{t +r/2}
\]
Therefore,
\[
I_{r,t} =  \begin{cases}
    \frac{\pi(-1)^{r/2}}{2^{2t-1}}\binom{2t}{t-r/2}\quad &\text{even} \; r \leq 2t\\
    0 \quad &\text{otherwise},
\end{cases}
\]
which results in
\[
A_{r,t} =  \begin{cases}
    \frac{1}{2^{2t}}\binom{2t}{t}\quad &r = 0\\
    \frac{(-1)^{r/2}}{2^{2t-1}}\binom{2t}{t-r/2}\quad &\text{even} \; 0<r \leq 2t\\
    0 \quad &\text{otherwise},
\end{cases}
\]

By using the above Fourier series we have 
\[\sum_{\text{even } m} \sin^{2t}\left(\omega_m/N\right) =\sum_{r= 0 }^t A_{2r,t} \sum_{\text{even } m} \cos\left(2r\omega_m/N\right). \]
When $N>k\geq r$, for $r>0$, $N$ does not divide $r$ and we can use the same lemma on sum of cosines of an arithmetic sequence to obtain
\[
\sum_{\text{even } m} \cos\left(2r\omega_m/N\right) = 
\begin{cases}
    \frac{\sin\left((N-1)r\pi/(2N)\right)}{\sin\left(r\pi/N\right)}\cdot\cos\left((N-1)r\pi/(2N)\right) \quad & \text{odd } N\\
    \frac{\sin\left(r\pi/2\right)}{\sin\left(r\pi/N\right)}\cdot\cos\left(r\pi/2\right) \quad & \text{even } N,
\end{cases}
\]
which simplifies to 
\[
\sum_{\text{even } m} \cos\left(2r\omega_m/N\right) = 
\begin{cases}
    -\frac{(-1)^{r}}{2} \quad & \text{odd } N\\
    0 \quad & \text{even } N
\end{cases}
\]
after using the double-angle and angle addition formulas for $\sin(x)$. For $r = 0$, we have \[
\sum_{\text{even } m} \cos\left(r\omega_m/N\right) = \left\lfloor\frac{N}{2}\right\rfloor
\]
For even $N>2t$,
\[
\sum_{\text{even } m} \sin^{2t}\left(\omega_m/N\right) = A_{0,t}\frac{N}{2} =  \frac{N}{2^{2t+1}}\binom{2t}{t}.
\]
For odd $N>2t$,
\begin{align*}
\sum_{\text{even } m} \sin^{2t}\left(\omega_m/N\right) &= A_{0,t}\frac{N-1}{2} + \sum_{r=1}^{t} A_{2r}\frac{(-1)^{r}}{2} =  \frac{N-1}{2^{2t+1}}\binom{2t}{t} - \sum_{r=1}^{t} \frac{1}{2^{2t}} \binom{2t}{t-r}\\
&= \frac{N-1}{2^{2t+1}}\binom{2t}{t} - \frac{1}{2^{2t}}\sum_{j = 0 }^{t-1}  \binom{2t}{j}\\
&= \frac{N-1}{2^{2t+1}}\binom{2t}{t} - \frac{1}{2^{2t+1}}\left(2^{2t}-\binom{2t}{t}\right)\\
&= \frac{N}{2^{2t+1}}\binom{2t}{t} -\frac{1}{2}
\end{align*}

Now note that the original summation
\(
\sum_{\text{even }m}\cos^2\left(\omega_m/N\right)\sin^{2t}\left(\omega_m/N\right)
\)
can be written as 
\[
\sum_{\text{even }m}\left(1-\sin^2\left(\omega_m/N\right)\right)\sin^{2t}\left(\omega_m/N\right)
= \sum_{\text{even }m}\sin^{2t}\left(\omega_m/N\right) - 
\sum_{\text{even }m}\sin^{2t+2}\left(\omega_m/N\right)
\] 
For even $N$,
\[
\sum_{\text{even }m}\cos^2\left(\omega_m/N\right)\sin^{2t}\left(\omega_m/N\right) = \frac{N}{2^{2t+1}}\binom{2t}{t} - \frac{N}{2^{2t+3}}\binom{2t+2}{t+1} = \frac{N}{2^{2t+2}(t+1)} \binom{2t}{t}
\]
If $N$ is odd,
\[
\sum_{\text{even }m}\cos^2\left(\omega_m/N\right)\sin^{2t}\left(\omega_m/N\right) = \frac{N}{2^{2t+1}}\binom{2t}{t}-\frac{1}{2} - \frac{N}{2^{2t+3}}\binom{2t+2}{t+1}+\frac{1}{2} = \frac{N}{2^{2t+2}(t+1)} \binom{2t}{t}
\]
So we proved that
\[
\sum_{\text{even }m}\cos^2\left(\omega_m/N\right)\sin^{2t}\left(\omega_m/N\right) = c_t N
\]
where
\[
c_t = \frac{1}{2^{2t+2}(t+1)} \binom{2t}{t}
\]

Combining all of these terms, we obtain
\begin{align*}
\mathrm{tr }(T_p^k\,M) &= \alpha_1 - \sum_{\text{even }m} \frac{\cos^2\left(\omega_m/N\right)}{N\sin^4\left(\omega_m/N\right)} + \frac{4kp}{N}\sum_{\text{even }m} \frac{\cos^2\left(\omega_m/N\right)}{\sin^2\left(\omega_m/N\right)}\\
&\quad -\sum_{t=2}^k\binom{k}{t}\frac{(-4p)^t}{N}\sum_{\text{even }m}\cos^2\left(\omega_m/N\right)\sin^{2t-4}\left(\omega_m/N\right)\\
&= 0 + \frac{4kp}{N}\cdot\frac{N(N-1)}{2}-\sum_{t=2}^k\binom{k}{t}\frac{(-4p)^t}{N}\cdot c_{t-2} N\\
&= 2(N-1)kp -4\sum_{t=2}^k\frac{(-1)^t}{t-1}\binom{k}{t} \binom{2t-4}{t-2}p^t 
\end{align*} 


\begin{lemma}
    For non-negative integers $k$, $l$, and $N> k + l$
    \[
\mathrm{tr }(T_p^kT_q^l\,M) =2(N-1)(kp+lq)  -\sum_{d=2}^{k+l}\sum_{t=\max(0,d-l)}^{\min(d,k)}(-1)^{d}\frac{4}{d-1} \binom{2d-4}{d-2}\binom{k}{t}\binom{l}{d-t}p^tq^{d-t}
\]
\end{lemma}
\begin{proof}
Proof of this lemma is similar to the previous one. 
\begin{align*}
\mathrm{tr }(T_p^kT_q^l\,M) &= 
\mathrm{tr} \left(  U \Sigma_p^k U^{\top}U \Sigma_q^l U^{\top} M \right)
= \mathrm{tr} \left(   \Sigma_p^k \Sigma_q^l U^{\top} M U \right)\\
&= \sum_{m=1}^{N} \alpha_m\,\bigl(\lambda^{(p)}_m\bigr)^k\bigl(\lambda^{(q)}_m\bigr)^l\\
&=\alpha_1 +\sum_{\text{even }m} \alpha_m\,\bigl(\lambda^{(p)}_m\bigr)^k\bigl(\lambda^{(q)}_m\bigr)^l\\
&= \alpha_1 -\sum_{\text{even }m} \frac{\cos^2\left(\omega_m/N\right)}{N\sin^4\left(\omega_m/N\right)}\,\left(1-4p\sin^2\left(\frac{\omega_m}{N}\right)\right)^k \left(1-4q\sin^2\left(\frac{\omega_m}{N}\right)\right)^l\\
&= \alpha_1 -\sum_{t=0}^k\sum_{s=0}^l\sum_{\text{even }m} \frac{\cos^2\left(\omega_m/N\right)}{N\sin^4\left(\omega_m/N\right)}\binom{k}{t}\binom{l}{s}\,\left(-4p\sin^2\left(\frac{\omega_m}{N}\right)\right)^t \left(-4q\sin^2\left(\frac{\omega_m}{N}\right)\right)^s\\
&= \alpha_1 -\sum_{t=0}^k\sum_{s=0}^l\frac{(-4p)^t(-4q)^s}{N}\binom{k}{t}\binom{l}{s}\sum_{\text{even }m}\cos^2\left(\frac{\omega_m}{N}\right)\sin^{2(t+s)-4}\left(\frac{\omega_m}{N}\right)\\
&= \alpha_1 -\sum_{d=0}^{k+l}\sum_{t=\max(0,d-l)}^{\min(d,k)}\frac{(-4p)^t(-4q)^{d-t}}{N}\binom{k}{t}\binom{l}{d-t}\sum_{\text{even }m}\cos^2\left(\frac{\omega_m}{N}\right)\sin^{2d-4}\left(\frac{\omega_m}{N}\right)\\
&= \underbrace{\alpha_1 -\frac{1}{N}\sum_{\text{even }m}\cos^2\left(\frac{\omega_m}{N}\right)\sin^{-4}\left(\frac{\omega_m}{N}\right)}_{=0}\\
&\quad +\frac{4pk+4ql}{N}\underbrace{\sum_{\text{even }m}\cos^2\left(\frac{\omega_m}{N}\right)\sin^{-2}\left(\frac{\omega_m}{N}\right)}_{=N(N-1)/2}\\
&\quad -\sum_{d=2}^{k+l}\sum_{t=\max(0,d-l)}^{\min(d,k)}\frac{(-4)^{d}p^tq^{d-t}}{N}\binom{k}{t}\binom{l}{d-t}\underbrace{\sum_{\text{even }m}\cos^2\left(\frac{\omega_m}{N}\right)\sin^{2d-4}\left(\frac{\omega_m}{N}\right)}_{=c_{d-2}N}.
\end{align*}
Therefore,
\begin{align*}
\mathrm{tr }(T_p^kT_q^l\,M) &= 2(N-1)(kp+lq) -\sum_{d=2}^{k+l}\sum_{t=\max(0,d-l)}^{\min(d,k)}(-4)^{d}c_{d-2}\binom{k}{t}\binom{l}{d-t}p^tq^{d-t}\\
&= 2(N-1)(kp+lq)  -\sum_{d=2}^{k+l}\sum_{t=\max(0,d-l)}^{\min(d,k)}(-1)^{d}\frac{4}{d-1} \binom{2d-4}{d-2}\binom{k}{t}\binom{l}{d-t}p^tq^{d-t}
\end{align*}
\end{proof}



Let us consider the behavior as $N \to \infty$. For fixed $m$, since $k = |i-j| \leq m$ the term $$
-4\sum_{t=2}^k\frac{(-1)^t}{t-1}\binom{k}{t} \binom{2t-4}{t-2}p^t $$
is upper bounded by a constant determined by $m$, denoted as $C_m$. Thus as $N$ grows this term is negligible. The same holds for $ -\sum_{d=2}^{k+l}\sum_{t=\max(0,d-l)}^{\min(d,k)}(-1)^{d}\frac{4}{d-1} \binom{2d-4}{d-2}\binom{k}{t}\binom{l}{d-t}p^tq^{d-t}$ since $k + l \leq 2m$. Thus we have when $t^*_m/m = t^*$, for all $i,j\in [m]$:
$$
\frac{N(N-1)}{2c_A^2m}\left(\mathcal{D}^{(2)}_{d_{MV}}\right)_{i,j} - \left(\mathcal{D}_Z\right)_{i,j} = 
\frac{N(N-1)}{2c_A^2m}\left(\mathcal{D}_{d^2_{MV}}\right)_{i,j} - \left(\mathcal{D}_Z\right)_{i,j} 
\leq \frac{C_m}{2m(N-1)} \sim O\left(\frac{1}{N}\right), 
$$ simply denote $D' := \frac{N(N-1)}{2c_A^2m}\left(\mathcal{D^A}^{(2)}_{d_{MV}}\right) = \frac{N(N-1)}{2c_A^2m}\left(\mathcal{D^A}_{d^2_{MV}}\right)$, we have
$$
\left(D'\right)_{i,j} - \left(\mathcal{D}_Z\right)_{i,j} \sim O\left(\frac{1}{N}\right),
$$
since every entry of both matrices are bounded by constant we also have
$$
\left(D'^{(2)}\right)_{i,j} - \left(\mathcal{D}^{(2)}_Z\right)_{i,j} \sim O\left(\frac{1}{N}\right).
$$

Let us denote the eigenvalues and the corresponding orthonormal eigenvectors of $-\frac{1}{2}H\mathcal{D}^{(2)}_ZH$ and $-\frac{1}{2}HD'^{(2)}H$ respectively by $\lambda_1, \lambda_2, \cdots, \lambda_m$ and $U_1,U_2,\cdots,U_m$; $\tilde{\lambda}_1, \tilde{\lambda}_2,\cdots, \tilde{\lambda}_m$ and $\tilde{U}_1,\tilde{U}_2,\cdots,\tilde{U}_m$.  
Note that for fixed $m$, both $\psi^1_{d^2_{MV}}$ and $\psi_Z$ are $m \times 1$ vectors as the first dimension of CMDS result of $\mathcal{D}_{d^2_{MV}}$ and $\mathcal{D}_Z$, and because $D' = \frac{N(N-1)}{2c_A^2m} \mathcal{D}_{d^2_{MV}}$,
we have:
$$
\frac{N(N-1)}{2c_A^2m}\psi^1_{d^2_{MV}} = \sqrt{\tilde{\lambda}_1} \tilde{U}_1, ~~~ \psi_Z = \sqrt{\lambda_1} U_1.
$$
Now we only need to show that: 
$$
\max_{i \in [m]}\bigg| \frac{N(N-1)}{2c_A^2m}\psi^1_{d^2_{MV}}(t_i) -  \psi_Z(t_i) \bigg| = \bigg\| \sqrt{\tilde{\lambda}_1} \tilde{U}_1 - \sqrt{\lambda_1} U_1\bigg\|_{2 \to \infty} \sim O\left(\frac{1}{N}\right).
$$
To show this we follow the proof in \cite{chen2024euclidean} Theorem 2.26. Note here $m$ is not growing as in \cite{chen2024euclidean} and it is treated as a constant. Denote  double centering matrix $H \in \RR^{m \times m} : = I -\frac{J}{m}$ where $I$ is the order $m$ identity matrix and $J$ is $m\times m$ all one matrix. Here we state several properties about $D'$ and $\mathcal{D}_Z$:
\begin{itemize}
    \item $\lambda_1 \sim \Theta(1)$ and $\lambda_i =0$ for $i \geq 2$; $ \left\Vert U_1 \right\Vert_{2 \to \infty} \sim \Theta(1)$.
    \item $\|D' -\mathcal{D}_Z\|_{\infty} \sim O\left(\frac{1}{N}\right)$, and $\|D'^{(2)} -\mathcal{D}_Z^{(2)}\|_{\infty} \sim O\left(\frac{1}{N}\right)$, $\|HD'^{(2)}H -H\mathcal{D}_Z^{(2)}H\|_{\infty} \sim O\left(\frac{1}{N}\right)$.
    \item $\left|\tilde{\lambda}_1-\lambda_1\right| \sim O\left(\frac{1}{N}\right)$,$\tilde{\lambda}_1 \sim \Theta(1)$, $\left|\sqrt{\tilde{\lambda}_1} -\sqrt{\lambda_1}\right| = \frac{\left|\tilde{\lambda}_1-\lambda_1\right|}{\sqrt{\tilde{\lambda}_1}+\sqrt{\lambda_1}} \sim \frac{O\left(\frac{1}{N}\right)}{\Theta\left(1\right)}  \sim O\left(\frac{1}{N}\right).$
\end{itemize}

Then because $\lambda_1 \sim \Theta(1)$ and  $\|HD'^{(2)}H -H\mathcal{D}_Z^{(2)}H\|_{\infty} \sim O\left(\frac{1}{N}\right)$, once $N$ is large enough, we will have $|\lambda_1|>4\left\Vert H\left(\mathcal{D}^{(2)}_{\varphi_m}-\mathcal{D}^{(2)}_Z\right)H\right\Vert_{\infty}$ and we can use Theorem 4.2 in \cite{cape2019two}. Thus there exists a $w=1$ or $w=-1$ such that:
\begin{align*}
&\left\Vert \sqrt{\tilde{\lambda}_1}\tilde{U}_1-w\sqrt{\lambda_1}U_1  \right\Vert_{2 \to \infty}  \leq \sqrt{ \tilde{\lambda}_1} \left\Vert \tilde{U}_1-wU_1\right\Vert_{2 \to \infty} + \left|\sqrt{\tilde{\lambda}_1} -\sqrt{\lambda_1}\right|\left\Vert U_1 \right\Vert_{2 \to \infty}  \\
&\leq  \sqrt{ \tilde{\lambda}_1}\frac{c\|HD'^{(2)}H -H\mathcal{D}_Z^{(2)}H\|_{\infty} \left\Vert U_1 \right\Vert_{2 \to \infty} }{\lambda_1}
+\left|\sqrt{\tilde{\lambda}_1} -\sqrt{\lambda_1}\right|\left\Vert U_1 \right\Vert_{2 \to \infty} \\
& \leq \left( \frac{ c \sqrt{ \tilde{\lambda}_1} \|HD'^{(2)}H -H\mathcal{D}_Z^{(2)}H\|_{\infty}}{\lambda_1}+\left|\sqrt{\tilde{\lambda}_1} -\sqrt{\lambda_1}\right| \right) \left\Vert U_1 \right\Vert_{2 \to \infty} \sim O\left(\frac{1}{N}\right). 
\end{align*}
Thus the result follows.

Now we consider CMDS on $\mathcal{D^A}_{\text{ind-}d_{MV}}$. Denote $a : = \frac{c_A^2}{6}\frac{N+1}{N-1}$, we note that $\mathcal{D^A}^{(2)}_{\text{ind-}d_{MV}} = a J - aI$ where $I$ is $m \times m$ identity matrix and $J$ is $m \times m$ all one matrix, then 
$$
-\frac{1}{2}H\mathcal{D^A}^{(2)}_{\text{ind-}d_{MV}}H = \frac{a}{2}H(I - J  )H = \frac{a}{2}\left( I - \frac{J}{m} \right)\left(I - J  \right) \left( I - \frac{J}{m} \right)    = \frac{a}{2}H.
$$
Note matrix $H= I - \frac{J}{m}$ is idempotent and has eigenvalues $1$ and $0$ with geometric multiplicity $m-1$ and $1$ respectively. Thus eigenvalues of $\frac{a}{2}H$ are $a/2 > 0$ and $0$ and the first $m-1$ eigenvectors correspond to the eigenvalue $a/2$. Then any orthonormal basis of $\RR^{m-1}$ can serve as these eigenvectors, denoted as $U_{m-1}$. The results follow.  


For $ \mathcal{D^A}_{\alpha -d_{MV}}^{(2)}$. Denote $B^{\mathcal{A}}_{d_{MV}} = -\frac{1}{2} H \mathcal{D^A}_{d_{MV}}^{(2)}  H$. Thus 
$$
-\frac{1}{2} H \mathcal{D^A}_{\alpha -d_{MV}}^{(2)} H = 
-\frac{1}{2} H \left( \alpha \mathcal{D^A}_{\text{ind-}d_{MV}}^{(2)} +(1-\alpha) \mathcal{D^A}_{d_{MV}}^{(2)}  \right) H
= (1-\alpha)B^{\mathcal{A}}_{d_{MV}}  + \frac{\alpha a }{2}H.
$$ 
%Further note that 
%$$
%B^{\mathcal{A}}_{d_{MV}} H = H B^{\mathcal{A}}_{d_{MV}}.
%$$
%Thus $B^{\mathcal{A}}_{d_{MV}}$ and $H$ commute and have the same orthonormal eigenvectors. 
Now denote the eigenvalues of $B^{\mathcal{A}}_{d_{MV}}$ as $\lambda_1>\lambda_2>\cdots>\lambda_m = 0$ and $U_1$,\ldots,$U_m$. Then for any orthonormal eigenvector of $B^{\mathcal{A}}_{d_{MV}}$ that is orthogonal to $\mathbf{1}$: $U_i$ such that $U^{\top}_i\mathbf{1} = 0$ because $B^{\mathcal{A}}_{d_{MV}} U_i  = \lambda_i U_i$ then
$$
\left(-\frac{1}{2} H \mathcal{D^A}_{\alpha -d_{MV}}^{(2)} H \right)U_i = \left( (1- \alpha) B^{\mathcal{A}}_{d_{MV}} + \frac{a\alpha}{2}H \right) U_i = \left((1-\alpha)\lambda_i + \frac{a \alpha}{2}\right)U_i.
$$ 
Those $U_i$ that are  orthogonal to $\mathbf{1}$ are also eigenvectors of $-\frac{1}{2} H \mathcal{D^A}_{\alpha -d_{MV}}^{(2)} H $ and correspond to the eigenvalue $(1-\alpha)\lambda_i + a \alpha/2$. So when $\lambda_1$ is the most positive eigenvalue of $B_{d_{\mathrm{MV}}}^{\mathcal{A}}$, $(1-\alpha)\lambda_1 + \frac{a \alpha}{2}$ will be the most positive eigenvalue of $-\frac{1}{2}H\mathcal{D}^{\mathcal{A}^{(2)}}_{\alpha-d_{\mathrm{MV}}}H.$ Thus 
$$
\psi^A_{d_{MV}} = \sqrt{\lambda_1} U_1, ~~ \psi^A_{\alpha -d_{MV}} = \sqrt{ (1-\alpha)\lambda_1 + \frac{a \alpha}{2} } U_1.
$$

For the $W_1$ distance, we know $X^A_{t_i} \sim u$ for all $i \in[m]$, thus $W_1(X^A_{t_i},X^A_{t_j}) = 0$ for all $i,j \in [m]$. 

For the expected average degree, since in Atlanta model, for any $i\in[m]$: $\EE[X_{t_i}] = \EE [u] = \frac{c_A}{2}$, the results follow. 
%$P = I - \frac{J}{m}$ is the projection matrix that project any $\RR^m$ vector into orthogonal complement of space spanned by $(1,1,...,1) \in \RR^m$. 

\end{proof}




\subsection{Proof of Theorem \ref{thm:W1_continuous}}
% \begin{proof}
% \textbf{Upper bound:} Let $\sigma:[n]\rightarrow[n]$ be the optimal permutation between $\{X_i\}, \{Y_j\}$, so that
% $$ W^1(\mu,\nu) = \frac{1}{n}\sum_{i=1}^n \|X_i - Y_{\sigma(i)}\|.$$ Then since this permutation is one possible choice of alignment between $\{\hat{X}_i\}, \{\hat{Y}_j\}$, but may not be optimal, we have
% \begin{align*} 
% W^1(\hat{\mu},\hat{\nu}) &\leq \frac{1}{n}\sum_{i=1}^n \|\hat{X}_i - \hat{Y}_{\sigma(i)}\|\\
% &\leq \frac{1}{n}\sum_{i=1}^n \left( \|\hat{X}_i - X_i\|+\|X_i-Y_{\sigma(i)}\|+\|Y_{\sigma(i)}-\hat{Y}_{\sigma(i)}\|\right)\\
% &\leq \frac{1}{n}\sum_{i=1}^n \left( \|\hat{X}-X\|_{2\rightarrow\infty}+ \|X_i-Y_{\sigma(i)}\|+\|Y-\hat{Y}\|_{2\rightarrow\infty}\right)\\
% &= W^1(\mu,\nu)+2\epsilon.
% \end{align*}

% \textbf{Lower bound:} We make use of the dual formulation for the optimal transport problem. Let $f:\RR^d\rightarrow\RR^d$ be a 1-Lipschitz function achieving the maximum value in the dual, so that 
% $$ W^1(\mu,\nu) = \frac{1}{n}\sum_{i=1}^n f(X_i) - f(Y_i).$$
% Since this is one possible 1-Lipschitz function, but may not be optimal for $\hat{\mu},\hat{\nu}$, we have
% \begin{align*}
% W^1(\hat{\mu},\hat{\nu})&\geq \frac{1}{n}\sum_{i=1}^n f(\hat{X}_i)-f(\hat{Y}_i)\\
% &= \frac{1}{n}\sum_{i=1}^n \left( f(X_i) + (f(\hat{X}_i)-f(X_i)) - f(Y_i) - (f(\hat{Y}_i)-f(Y_i))\right)\\
% &\geq \frac{1}{n}\sum_{i=1}^n \left( f(X_i) - f(Y_i) - |f(\hat{X}_i)-f(X_i)|- |f(\hat{Y}_i)-f(Y_i)|\right)\\
% &\geq \frac{1}{n}\sum_{i=1}^n \left( f(X_i)-f(Y_i) - \|\hat{X}_i-X_i\| - \|\hat{Y}_i-Y_i\|\right)\\
% &\geq \frac{1}{n}\sum_{i=1}^n \left( f(X_i)-f(Y_i) - \|\hat{X}-X\|_{2\rightarrow\infty} - \|\hat{Y}-Y\|_{2\rightarrow\infty}\right)\\
% &\geq W^1(\mu,\nu)-2\epsilon,
% \end{align*}
% making use of the 1-Lipschitz condition for $f$. 
% \end{proof}
% Note that this calculation actually showed the tighter lower bound 
% $$ W^1(\hat{\mu},\hat{\nu}) \geq W^1(\mu,\nu) - \frac{1}{n}\sum_{i=1}^n \|\hat{X}_i-X_i\| - \frac{1}{n}\sum_{i=1}^n \|\hat{Y}_i-Y_i\|.$$ The sum of the Euclidean norms of the rows of a matrix is a norm on the space of matrices, but is not induced since $\|I\|_{+,2}=n$. It does satisfy some sub-multiplicativity properties, though:
% $$\|AB\|_{+,2} = \sum_{i=1}^n \|e_i^T AB\|_2 \leq \sum_{i=1}^n \|e_i^T A\|_2 \|B\|_2 = \|A\|_{+,2}\|B\|_2.$$
% There might be more, idk.

\begin{proof}
Suppose 
\begin{equation}
\label{eq:wmin2inf}
\min_W \|\hat{\bX}-\bX W\|_{2\rightarrow\infty}\leq \epsilon, \min_W \|\hat{\bY}-\bY W\|_{2\rightarrow\infty}\leq \epsilon,
\end{equation}
and recall that 
$$\procwphat[1](\bX,\bY)=\min_W \min_\sigma \frac{1}{n}\sum_{i=1}^n \|X_i-W Y_{\sigma(i)}\|.$$
We will show that
$$ |\procwphat[1](\hat{\bX},\hat{\bY})-\procwphat[1](\bX,\bY)|\leq 2\epsilon.$$
Let $\sigma:[n]\rightarrow[n]$ be the optimal permutation and $W_{XY}$ be the optimal orthogonal matrix for $\bX,\bY$, so that
$$\procwphat[1](\bX,\bY)=\frac{1}{n}\sum_{i=1}^n \|X_i-W_{XY}Y_{\sigma(i)}\|.$$
Then since this permutation is one possible choice of alignment for $\hat{\bX},\hat{\bY}$, but may not be optimal, we have
\begin{align*}
\procwphat[1](\hat{\bX},\hat{\bY})&=\min_{\sigma'}\min_{W} \frac{1}{n}\sum_{i=1}^n \|\hat{X}_i-W\hat{Y}_{\sigma'(i)}\|\\
&\leq \min_W \frac{1}{n}\sum_{i=1}^n \|\hat{X}_i-W\hat{Y}_{\sigma(i)}\|\\
&\leq \min_{W,W_1,W_2} \frac{1}{n}\sum_{i=1}^n \left(\|\hat{X}_i-W_1X_i\|+\|W_1 X_i-W_2Y_{\sigma(i)}\|+\|W_2Y_{\sigma(i)}-W\hat{Y}_{\sigma(i)}\|\right).
\end{align*}
Let $W_1'$ achieve $\|\hat{\bX}-\bX W_1'\|_{2\rightarrow\infty}\leq \epsilon$, as guaranteed to exist by Equation~\ref{eq:wmin2inf}. Then we select $W_2'=(W_1')^\top W_{XY}$, and $W'=W_2' (W_3')^\top $, where $W_3'$ achieves $\|\hat{\bY}-\bY W_3'\|_{2\rightarrow\infty}\leq \epsilon$ (again by Equation~\ref{eq:wmin2inf}). Since these are three particular choices of orthogonal matrices, we have
\begin{align*}
\procwphat[1](\hat{\bX},\hat{\bY})&\leq \frac{1}{n}\sum_{i=1}^n \left(\|\hat{\bX}-\bX W_1'\|_{2\rightarrow\infty}+\|(W_1')^\top(X_i-W_{XY} Y_{\sigma(i)})\|+\|(\bY-\hat{\bY}W_3')(W_2')^\top\|_{2\rightarrow\infty}\right)\\
&\leq \procwphat[1](\bX,\bY)+2\epsilon.
\end{align*}
For any $\sigma$ and $W$, after choosing $W_1$, $W_2$ in light of Equation~\ref{eq:wmin2inf}, we have
\begin{align*}
\frac{1}{n}\sum_{i=1}^n \|\hat{X}_i-W\hat{Y}_{\sigma(i)}\|&\geq \frac{1}{n}\sum_{i=1}^n \left(\|W_1 X_i-WW_2Y_{\sigma(i)}\|-\|\hat{X}_i-W_1X_i\|-\|W(W_2Y_{\sigma(i)}-\hat{Y}_{\sigma(i)})\|\right)\\
&\geq \frac{1}{n}\sum_{i=1}^n \|W_1X_i-WW_2Y_{\sigma(i)}\|-2\epsilon\\
&\geq \procwphat[1](\bX,\bY)-2\epsilon.
\end{align*}
Now minimizing over $\sigma$ and $W$ gives 
$$
\procwphat[1](\hat{\bX},\hat{\bY})\geq \procwphat[1](\bX,\bY)-2\epsilon.
$$
\end{proof}




\subsection{Proof of Theorem \ref{thm:W1_as_converge}}
We will make use of the following lemma: 
\begin{lemma}
\label{lem:supsclose}
For a metric space $(\mathcal{M},d)$ and functions $L_1,L_2: \mathcal{M}\rightarrow\RR, $ if $|L_1(x)-L_2(x)|\leq \epsilon$ for all $x\in\mathcal{M}$, then 
$$\left| \sup_{x\in\mathcal{M}} L_1(x) - \sup_{y\in\mathcal{M}} L_2(y)\right|\leq \epsilon.$$
\end{lemma}
\begin{proof}
Indeed, let $\delta>0$, and let $x^*\in\mathcal{M}$ satisfy $L_1(x^*) \geq \sup_{x\in\mathcal{M}} L_1(x)-\delta$. Then 
\begin{align*}
\sup_{y\in \mathcal{M}} L_2(y) & \geq L_2(x^*)\\
&\geq L_1(x^*)-\epsilon\\
&\geq \sup_{x\in \mathcal{M}} L_1(x) - \epsilon - \delta.
\end{align*}
Since the choice of $\delta>0$ was arbitrary, we have
$$ \sup_{y\in \mathcal{M}} L_2(y) \geq \sup_{x\in \mathcal{M}}L_1(x)-\epsilon.$$ Now observe that the assumptions on $L_1$ and $L_2$ are symmetric, so we must also have 
$$ \sup_{x\in \mathcal{M}} L_1(x) \geq \sup_{y\in \mathcal{M}} L_2(y) - \epsilon,$$ which completes the proof of the lemma.
\end{proof}

% % \textcolor{red}{begin remove}
% % \begin{proof}
% % Now we prove the theorem. We must find a set of full measure $\Omega$ on which we may apply the following argument: Let $\epsilon>0$. We wish to apply Lemma~\ref{lem:supsclose} to the set 
% % $$\mathcal{F}_1 = \{f: B_d\rightarrow B_d: f\text{ is 1-Lipschitz}, f(0)=0\},$$
% % with the functions
% % \begin{align*} 
% % L_{1,n}(f) &= \EE_{x\sim \mu_n} f(x) - \EE_{y\sim \nu_n} f(y),\\
% % L_2(f) &= \EE_{x\sim \mu} f(x) - \EE_{y\sim \nu} f(y),
% % \end{align*}
% % where $n$ is a sufficiently large number that may depend on $\epsilon$. The choice of $f(0)=0$ is a convenient normalization, since adding a constant to $f$ affects neither the Lipschitz property nor the objective function value in the Kantorovich problem. Then since $W^1(\mu_n, \nu_n) = \sup_{f\in \mathcal{F}_1} L_{1,n}(f)$, and $W^1(\mu,\nu) = \sup_{f\in \mathcal{F}_1} L_2(f)$, the lemma shows that so long as 
% % $$\sup_{f\in\mathcal{F}_1} |L_{1,n}(f)-L_2(f)|\leq \epsilon,$$
% % then we have
% % $$ |W^1(\mu_n,\nu_n)-W^1(\mu,\nu)|= |\sup_{f\in \mathcal{F}_1} L_{1,n}(f) - \sup_{g\in \mathcal{F}_1} L_2(g)| \leq \epsilon.$$ Since $\epsilon>0$ was arbitrary, on $\Omega$ we have
% % $$ W^1(\mu_n,\nu_n)\rightarrow W^1(\mu,\nu),$$
% % as required by the theorem statement.

% % Note that for any $f\in \mathcal{F}_1$, we have that on a set of full measure $\Omega_f$, $L_{1,n}(f)\rightarrow L_2(f)$ as $n\rightarrow\infty$ by the ordinary strong law of large numbers, since all $X_i\overset{\text{iid}}{\sim} \mu$, $Y_j\overset{\text{iid}}{\sim} \nu$, and the Lipschitz condition on $f$, plus the compactness of $B_d$, are sufficient to guarantee that $\EE_{x\sim \mu} f(x), \EE_{y\sim \nu} f(y)$ are finite. The issue is that $\mathcal{F}_1$ is uncountable [since for any $x\in B_d$ we may define $f_x(y) = (\|x\|-\|y-x\|)\mathbf{1}_{B(x,\|x\|)}(y)$, any two of which are distinct since $f_x$ is the unique function among these maximized at $x$], so we cannot take a simple union bound. However, suppose we have a covering for $\mathcal{F}_1$ of the form $\bigcup_{i=1}^N B(f_i,\epsilon)$, where $f_i \in\mathcal{F}_1$, and $B(f_i, \epsilon) = \{g\in \mathcal{F}_1: \|f_i-g\|_{\infty}\leq \epsilon\}.$ The existence of such a covering is guaranteed by the Arzela-Ascoli theorem, since $\mathcal{F}_1$ is closed with respect to the sup-norm in the space of continuous functions; it is clearly equicontinuous; and it is bounded since $|f(z)|=|f(z)-f(0)|\leq \|z\|\leq 1$ for all $f\in \mathcal{F}_1$ and $z\in B_d$. The open cover $\{B(f,\epsilon); f\in\mathcal{F}_1\}$ then has a finite subcover by compactness, which gives a covering of the desired form. Then for any $n$, $g\in B(f,\epsilon)$, and $\{z_i\}_{i=1}^n\subset B_d,$ we have
% % $$
% % \left|\frac{1}{n}\sum_{i=1}^n g(z_i) - \frac{1}{n}\sum_{j=1}^n f(z_j)\right|\leq \frac{1}{n}\sum_{i=1}^n|g(z_i)-f(z_i)| \leq \epsilon,
% % $$
% % so $|L_{1,n}(g)-L_{1,n}(f)|\leq 2\epsilon,$ and $|L_2(g)-L_2(f)|\leq 2\epsilon$ by a similar argument. So for a given $\epsilon>0$, select a covering for $\mathcal{F}_1$ of the given type. For each $1\leq i\leq N(\epsilon)$, let $\Omega_i$ be the set of full measure for which $L_{1,n}(f_i)\rightarrow L_2(f_i)$ as $n\rightarrow\infty$. Then for each $i$, for any $\omega\in \Omega_i$, there is a sufficiently large $K_i=K_i(\epsilon,\omega)$ such that for all $k\geq K_i$, $|L_{1,k}(f_i)-L_2(f_i)|\leq \epsilon.$ So for any $\omega\in \Omega(\epsilon):=\bigcap_{i=1}^N \Omega_i$ and $k\geq K(\epsilon):=\max_{1\leq i\leq N} K_i(\epsilon,\omega)$, we have for $g\in B(f_i,\epsilon)$:
% % $$|L_{1,k}(g)-L_2(g)| \leq |L_{1,k}(g)-L_{1,k}(f_i)|+|L_{1,k}(f_i)-L_2(f_i)|+|L_2(f_i)-L_2(g)| \leq 5\epsilon. $$
% % But since $\{B(f_i,\epsilon)\}$ covers $\mathcal{F}_1$, this is true for all $g\in \mathcal{F}_1$. Put briefly, for $\omega\in\Omega(\epsilon)$, once $k\geq K(\epsilon)$, every $g\in\mathcal{F}_1$ satisfies $$|L_{1,k}(g)-L_2(g)| \leq 5\epsilon.$$ Consider the sequence $\epsilon_m=1/m$, and define $\Omega:= \bigcap_{m=1}^{\infty}\Omega(1/m)$, which is also a set of full measure. Then for any $\omega\in \Omega$ and $\epsilon>0$, letting $m>1/\epsilon$, $\omega \in \Omega(1/m)$, so there exists $K(1/m)$ so large that for any $k\geq K(1/m)$, every $g\in\mathcal{F}_1$ satisfies $$|L_{1,k}(g)-L_2(g)|\leq 5/m \leq 5\epsilon.$$ In other words, $\Omega$ is the desired set.
% % \end{proof}






% % \subsection{Proof of Theorem \ref{thm:W1-concentration}}

% % \begin{proof}
% % In \cite{GottliebCovering}, the following is Lemma 6:
% % \begin{lemma}
% % Let $\mathcal{F}_L$ be the collection of $L$-Lipschitz functions mapping the metric space $(\mathcal{X},\rho)$ to $[0,1]$. Then the covering number of $\mathcal{F}_L$ may be estimated in terms of the covering number of $\mathcal{X}$:
% % $$\mathcal{N}(\epsilon,\mathcal{F}_L,\|\cdot\|_{\infty})\leq \left(\frac{8}{\epsilon}\right)^{\mathcal{N}(\epsilon/8L,\mathcal{X},\rho)}.$$
% % \end{lemma}

% % In the present case, we have $\mathcal{X}=B_d$, $\rho=\|\cdot\|$, and all $f\in\mathcal{F}_1$ map into [0,1]. So we have
% % $$\mathcal{N}(\epsilon,\mathcal{F}_1,\|\cdot\|_{\infty})\leq \left(\frac{8}{\epsilon}\right)^{\mathcal{N}(\epsilon/8,B_d,\|\cdot\|)}\leq \left(\frac{8}{\epsilon}\right)^{(16/\epsilon+1)^d},$$ using the covering number bound from \cite{vershynin2018high}. 

% % Since $\mu,\nu$ are distributions on $B_d$, and $f$ is 1-Lipschitz, we may apply Hoeffding's inequality:
% % \begin{align*}
% % \PP\left[\left|\frac{1}{n}\sum_{i=1}^n f(X_i) - \EE_{X\sim\mu} f(X)\right|>t\right]&\leq 2\exp\left(\frac{-nt^2}{2}\right),\\
% % \PP\left[\left|\frac{1}{n}\sum_{j=1}^n f(Y_j) - \EE_{Y\sim\nu} f(Y)\right|>t\right]&\leq 2\exp\left(\frac{-nt^2}{2}\right).\\
% % \end{align*}
% % As such, for a given $f\in\mathcal{F}_1$, 
% % $$ \PP[|L_{1,n}(f)-L_2(f)|>n^{-1/(d+3)}]=4\exp(-n^{(d+1)/(d+3)}/2).$$ Consider the $\epsilon$-covering for $\mathcal{F}_1$ with $\epsilon=n^{-1/(d+3)}$. Then from the previous argument, we have that
% % $$\PP\left[\forall\;g\in\mathcal{F}_1\; |L_{1,n}(g)-L_2(g)|\leq 5n^{-1/(d+3)}\right]\geq 1- \left(\frac{8}{\epsilon}\right)^{(16/\epsilon+1)^d} \cdot [4\exp(-n^{(d+1)/(d+3)}/2)].$$

% % The log of the complementary probability on the right hand side is
% % \begin{align*}
% % \log \left[\left(\frac{8}{\epsilon}\right)^{(16/\epsilon+1)^d} \cdot [4\exp(-n^{(d+1)/(d+3)}/2)]\right]& = (16n^{1/(d+3)}+1)^d \log(8n^{1/(d+3)}) -n^{(d+1)/(d+3)}/2+\log(4)\\
% % &\leq C_d n^{d/(d+3)}\log(n)- C n^{(d+1)/(d+3)},
% % \end{align*}
% % which goes to $-\infty$ as $n\rightarrow\infty$. But on this event, we have that
% % $$|W^1(\mu_n,\nu_n)-W^1(\mu,\nu)|\leq 5n^{-1/(d+3)}.$$ Also 
% % \[
% % C_d\,n^{d/(d+3)}\log n \;\le\;\tfrac12\,C\,n^{(d+1)/(d+3)},
% % \]
% % so that
% % \[
% % C_d n^{d/(d+3)}\log n \;-\;C\,n^{(d+1)/(d+3)}
% % \;\le\;-\,\tfrac12\,C\,n^{(d+1)/(d+3)}.
% % \]
% % so we have the desired result.
% % \end{proof}

% % \textcolor{red}{End remove}



% In \cite{GottliebCovering}, the following is Lemma 6:
% \begin{lemma}
% Let $\mathcal{F}_L$ be the collection of $L$-Lipschitz functions mapping the metric space $(\mathcal{X},\rho)$ to $[0,1]$. Then the covering number of $\mathcal{F}_L$ may be estimated in terms of the covering number of $\mathcal{X}$:
% $$\mathcal{N}(\epsilon,\mathcal{F}_L,\|\cdot\|_{\infty})\leq \left(\frac{8}{\epsilon}\right)^{\mathcal{N}(\epsilon/8L,\mathcal{X},\rho)}.$$
% \label{lem:f1coveringnumber}
% \end{lemma}

In \cite{chewi2025statistical}, the following is Lemma 2.7:
\begin{lemma}
\label{lem:f1coveringnumber}
    Let $\mathcal{F}_1$ be the collection of 1-Lipschitz functions from $[0,1]^d\rightarrow\RR$ with $f(0)=0$. Then $$\mathcal{N}(\epsilon,\mathcal{F}_1,\|\cdot\|_{\infty})\leq \exp\left[C\left(\frac{4\sqrt{d}}{\epsilon}\right)^d\right].$$
\end{lemma}



% In the present case, we have $\mathcal{X}=B_d$, $\rho=\|\cdot\|$, and all $f\in\mathcal{F}_1$ map into [0,1]. So we have
% $$\mathcal{N}(\epsilon,\mathcal{F}_1,\|\cdot\|_{\infty})\leq \left(\frac{8}{\epsilon}\right)^{\mathcal{N}(\epsilon/8,B_d,\|\cdot\|)}\leq \left(\frac{8}{\epsilon}\right)^{(16/\epsilon+1)^d},$$ using the covering number bound for $B_d$ from \cite{vershynin2018high}. 
\begin{proof}
For $n\geq 1$, $f\in \mathcal{F}_1$, and some real orthogonal matrix $W$, let
\begin{align*}
L_{1,n}(f,W)&= \EE_{x\sim \mu_\bX} f(x) - \EE_{y\sim \mu_\bY} f(Wy),\\
L_2(f,W) &= \EE_{x\sim \mu} f(x) - \EE_{y\sim \nu} f(Wy).
\end{align*}
Recall that by the Kantorovich dual formulation of the Wasserstein-1 problem, 
\begin{align*}
\procwphat[1](\bX,\bY)&=\procwp[1](\mu_{\bX},\mu_{\bY})=\min_W W_1(\mu_\bX,\mu_{\bY W^\top})= \min_W \sup_{f\in\mathcal{F}_1} L_{1,n}(f,W),\\
\procwp[1](\mu_X,\mu_Y)&=\min_W W_1(\mu_X,\mu_{Y}\circ W^\top)=\min_W\sup_{f\in\mathcal{F}_1}L_2(f,W).
\end{align*}
Since $L_{1,n}(f,W)=L_{1,n}(f+c,W)$ and $L_2(f,W)=L_2(f+c,W)$ for any constant $c$, there is no loss of generality in taking $f(0)=0$. On a set of full measure $\Omega_{f,W}$, $L_{1,n}(f,W)\rightarrow L_2(f,W)$ as $n\rightarrow\infty$ by the ordinary strong law of large numbers, since all $X_i\overset{\text{iid}}{\sim} \mu_X$, $Y_j\overset{\text{iid}}{\sim} \mu_Y$, and the Lipschitz condition on $f$, plus the compactness of $B_d$, are sufficient to guarantee that $\EE_{x\sim \mu} f(x), \EE_{y\sim \nu} f(Wy)$ are finite. We want to show that there is a set of full measure on which this convergence takes place uniformly in $(f,W)$, which we will accomplish using a covering for $\mathcal{F}_1\times \mathcal{O}^{d}$. 

Consider a covering for $\mathcal{F}_1\times \mathcal{O}^d$ of the form $\bigcup_{i=1}^N B((f_i,W_i),\epsilon)$, where $f_i \in\mathcal{F}_1$, $W_i\in \mathcal{O}^d$, and $B((f_i,W_i), \epsilon) = \{(g,W)\in \mathcal{F}_1\times \mathcal{O}^d: 2\|g-f_i\|_{\infty}+\|W-W_i\| \leq \epsilon\}.$ The existence of such a covering is guaranteed by compactness of $\mathcal{F}_1\times \mathcal{O}^d$, where the compactness of $\mathcal{F}_1$ is supplied by Lemma~\ref{lem:f1coveringnumber}. The open cover $\{B((f,W),\epsilon); (f,W)\in\mathcal{F}_1\times \mathcal{O}^d\}$ then has a finite subcover by compactness, which gives a covering of the desired form. Then for any $n$, $(g,R)\in B((f,W),\epsilon)$, and $\{x_i\}_{i=1}^n, \{y_i\}_{i=1}^n\subset B_d,$ we have
\begin{align*}
|L_{1,n}(g,R)-L_{1,n}(f,W)|&=\left|\frac{1}{n}\sum_{i=1}^n [g(x_i)-g(Ry_i)] - \frac{1}{n}\sum_{j=1}^n [f(x_j)-f(Wy_j)]\right|\\
&\leq \frac{1}{n}\sum_{i=1}^n|g(x_i)-f(x_i)|+|g(Ry_i)-g(Wy_i)|+|g(Wy_i)-f(Wy_i)|\\
&\leq \frac{1}{n}\sum_{i=1}^n(2\|g-f\|_{\infty}+\|Ry_i-Wy_i\|)\\
&\leq 2\|g-f\|_{\infty}+\|R-W\|\leq \epsilon,
\end{align*}
and $|L_2(g,R)-L_2(f,W)|\leq \epsilon$ by a similar argument. So for a given $\epsilon>0$, select a covering for $\mathcal{F}_1\times \mathcal{O}^d$ of the given type. For each $1\leq i\leq N(\epsilon)$, let $\Omega_i$ be the set of full measure for which $L_{1,n}(f_i,W_i)\rightarrow L_2(f_i,W_i)$ as $n\rightarrow\infty$. Then for each $i$, for any $\omega\in \Omega_i$, there is a sufficiently large $K_i=K_i(\epsilon,\omega)$ such that for all $k\geq K_i$, $|L_{1,k}(f_i,W_i)-L_2(f_i,W_i)|\leq \epsilon.$ So for any $\omega\in \Omega(\epsilon):=\bigcap_{i=1}^N \Omega_i$ and $k\geq K(\epsilon):=\max_{1\leq i\leq N} K_i(\epsilon,\omega)$, we have for $(g,R)\in B((f_i,W_i),\epsilon)$:
$$|L_{1,k}(g,R)-L_2(g,R)| \leq |L_{1,k}(g,R)-L_{1,k}(f_i,W_i)|+|L_{1,k}(f_i,W_i)-L_2(f_i,W_i)|+|L_2(f_i,W_i)-L_2(g,R)| \leq 3\epsilon. $$
But since $\{B((f_i,W_i),\epsilon)\}$ covers $\mathcal{F}_1\times \mathcal{O}^d$, this bound holds for all $(g,R)\in \mathcal{F}_1\times \mathcal{O}^d$. Put briefly, for $\omega\in\Omega(\epsilon)$, once $k\geq K(\epsilon)$, every $(g,R)\in\mathcal{F}_1\times \mathcal{O}^d$ satisfies $$|L_{1,k}(g,R)-L_2(g,R)| \leq 3\epsilon.$$ Consider the sequence $\epsilon_m=1/m$, and define $\Omega:= \bigcap_{m=1}^{\infty}\Omega(1/m)$, which is also a set of full measure. Then for any $\omega\in \Omega$ and $\epsilon>0$, letting $m>1/\epsilon$, $\omega \in \Omega(1/m)$, so there exists $K(1/m)$ so large that for any $k\geq K(1/m)$, every $(g,R)\in\mathcal{F}_1\times \mathcal{O}^d$ satisfies $$|L_{1,k}(g,R)-L_2(g,R)|\leq 3/m \leq 3\epsilon.$$ We apply Lemma~\ref{lem:supsclose} to get 
$$ \left|\sup_{f\in\mathcal{F}_1}L_{1,n}(f,W)-\sup_{g\in\mathcal{F}_1}L_2(g,W)\right|\leq 3\epsilon,$$ and apply Lemma~\ref{lem:supsclose} again to get
$$|\procwphat[1](\bX,\bY)-\procwp[1](\mu_X,\mu_Y)|=\left|\inf_W \sup_{f\in \mathcal{F}_1} f_{1,n}(f,W) - \inf_R \sup_{g\in\mathcal{F}_1} L_2(g,R)\right|\leq 3\epsilon. $$ Since $\epsilon>0$ was arbitrary, on $\Omega$ we have $$\procwphat[1](\bX,\bY)\rightarrow \procwp[1](\mu_X,\mu_Y),$$
as required.

Now we wish to prove the concentration result. Since $\mu_X, \mu_Y$ are distributions on $B_d^+$, $W\in\mathcal{O}^d$, and $f$ is 1-Lipschitz, we may apply Hoeffding's inequality:
\begin{align*}
\PP\left[\left|\frac{1}{n}\sum_{i=1}^n f(X_i) - \EE_{X\sim\mu_X} f(X)\right|>t\right]&\leq 2\exp\left(\frac{-nt^2}{2}\right),\\
\PP\left[\left|\frac{1}{n}\sum_{j=1}^n f(WY_j) - \EE_{Y\sim\mu_Y} f(WY)\right|>t\right]&\leq 2\exp\left(\frac{-nt^2}{2}\right).\\
\end{align*}
As such, for a given $(f,W)\in\mathcal{F}_1\times \mathcal{O}^d$, 
$$ \PP[|L_{1,n}(f,W)-L_2(f,W)|>t]\leq 4\exp(-nt^2/2).$$ Consider the $\epsilon$-covering for $\mathcal{F}_1\times \mathcal{O}^d$ with $\epsilon=t$. Then from the consistency argument, we have that
\begin{equation}
\label{eq:procwassprobbound}
\PP\left[\forall\;(g,R)\in\mathcal{F}_1\times\mathcal{O}^d\; |L_{1,n}(g,R)-L_2(g,R)|\leq 3t\right]\geq 1- \mathcal{N}(t,\mathcal{F}_1\times\mathcal{O}^d,d)[4\exp(-nt^2/2)].\end{equation}

Using Lemma~\ref{lem:f1coveringnumber}, we may bound 
\begin{align*}
\mathcal{N}(t,\mathcal{F}_1\times\mathcal{O}^d,d) &\leq \mathcal{N}(t/4,\mathcal{F}_1,\|\cdot\|_{\infty})\mathcal{N}(t/2,\mathcal{O}^d,\|\cdot\|)\\
&\leq \exp\left[C\left(\frac{4\sqrt{d}}{t}\right)^d\right] \left(1+\frac{4\sqrt{d}}{t}\right)^{d^2},
\end{align*}
where we bounded the latter covering number as follows:
$$ \mathcal{N}(\epsilon,\mathcal{O}^d,\|\cdot\|)\leq \mathcal{N}(\epsilon,\mathcal{O}^d,\|\cdot\|_F)\leq \mathcal{N}(\epsilon/\sqrt{d},B_d,\|\cdot\|)^d,$$ since any real orthogonal matrix has $d$ columns in $B_d$.

So the log of the complementary probability on the right hand side of Equation~\ref{eq:procwassprobbound} when $t=C_d n^{-1/(d+2)}$ for a sufficiently large constant $C_d$ is  
\begin{align*}
\log\left[\mathcal{N}(t,\mathcal{F}_1\times\mathcal{O}^d,d)[4\exp(-nt^2/2)]\right] &\leq C\left(\frac{4\sqrt{d}}{t}\right)^d + d^2\log\left(1+\frac{4\sqrt{d}}{t}\right)-\frac{nt^2}{2}+\log(4)\\
&\sim (C(4\sqrt{d}/C_d)^d+1) n^{d/(d+2)} - C_d n^{d/(d+2)}/2\leq -n^{d/(d+2)}.
\end{align*}

Theorem 1 from \cite{rubin2022statistical} proves convergence of the ASE estimates of the latent position matrices to the true ones with respect to the $\|\cdot\|_{2\rightarrow\infty}$ norm. They suppose that scaled latent positions $\{\xi_i\}_{i=1}^n$ are an i.i.d. sample drawn from some latent position distribution $F$ on $\RR^d$, where $F$ has a second moment matrix $\EE_F[\xi\xi^\top]$ that is positive definite with rank $d$, and that the latent positions are defined as $X_i=\sqrt{\rho_n}\xi_i$ for some quantity $\rho_n\in(0,1]$. We will consider the case that $\rho_n\equiv 1$. They further suppose that $A_X\in\{0,1\}^{n\times n}$ is drawn from a generalized random dot product graph, so that $\EE[A_X]= \bX I_{p,q} \bX^\top$; we will consider the special case that $p=d$ and $q=0$. They define a universal constant $c>1$ which gives a power of the logarithmic factor in their assumptions, as well as in the upper bound for the $\|\cdot\|_{2\rightarrow\infty}$ norm. With our simplifying assumptions, the theorem then reads that for any $a>0$, there is a constant $C_a>0$ such that for sufficiently large $n$, with probability at least $1-n^{-a}$,
$$
\min_{W\in\mathcal{O}^d}\|\hat{X}-XW\|_{2,\infty} \leq C_a \frac{\log^c(n)}{\sqrt{n}}.
$$
We note that a more stringent bound might be obtained from Corollary 4.1 of \cite{xie2024entrywise}, but either bound reveals that this is a lower-order term in the present setting.

As such, given $A\sim\mathrm{RDPG}(\bX), B\sim\mathrm{RDPG}(\bY)$ with ASEs $\hat{\bX}, \hat{\bY}$, there are events $\mathcal{A}_n,\mathcal{B}_n$ each with probability at least $1-n^{-a}$ such that
\begin{align*}
\text{On }\mathcal{A}_n: &\min_{W\in\mathcal{O}^d}\|\hat{\bX}-\bX W\|_{2\rightarrow\infty} \leq C_a \frac{\log^c(n)}{\sqrt{n}},\\
\text{On }\mathcal{B}_n: &\min_{W\in\mathcal{O}^d}\|\hat{\bY}-\bY W\|_{2\rightarrow\infty} \leq C_a' \frac{\log^c(n)}{\sqrt{n}}.\\
\end{align*}
So on $\mathcal{A}_n\cap \mathcal{B}_n$, Theorem~\ref{thm:W1_continuous} gives us
$$|\procwp[1](\hat{\mu}_n,\hat{\nu}_n)-\procwp[1](\mu_n,\nu_n)|\leq C_a'' \frac{\log^c(n)}{\sqrt{n}}.$$


From the first part of this theorem, there is an event $\mathcal{C}_n$ with probability at least $1-\exp(-n^{d/(d+2)})$ such that
$$
|\procwp[1](\mu_n,\nu_n)-\procwp[1](\mu,\nu)|\leq C_d n^{-1/(d+2)}.
$$
So with probability at least $1-3n^{-a}$,
$$
|\procwp[1](\hat{\mu}_n,\hat{\nu}_n)-\procwp[1](\mu,\nu)|\leq (C_d+1)n^{-1/(d+2)}.
$$



% Alternatively, Corollary 4.1 from \cite{xie2024entrywise} offers the bound
% \begin{align*}
% \min_{W\in\mathcal{O}^d} \|\hat{X}-XW\|_{2\rightarrow\infty} &\leq \frac{\|U\|_{2\rightarrow\infty}}{\sqrt{n}\lambda_d(\Delta_n)^2}\max\left\{\frac{\log^{1/2}(n)}{\lambda_d(\Delta_n)^2},\frac{\kappa(\Delta_n)}{\lambda_d(\Delta_n)^2},\log(n)\right\}\\
% &\quad+\frac{\log^{1/2}(n)\|U\|_{2\rightarrow\infty}}{\lambda_d(\Delta_n)^{1/2}},
% \end{align*}
% where $\Delta_n=\frac{1}{n}X^\top X.$ \textcolor{red}{This bound applies under 5 conditions that you can find in the paper, but I need to see how that translates into LPP world.}


\end{proof}

% \subsection{Proof of Theorem \ref{thm:A-L-W1} }
% \begin{proof}
% For London model, assume $i = j-1$, then $X^L_{t_j} = X^L_{t_i} + Z$ where $Z \sim \text{Bernoulli}(p)$ if $j \le t^*_m$ and $Z \sim \text{Bernoulli}(q)$  if $j \ge t^*_m +1$. Either case, $X_{t_j}$ stochastically dominants $X^L_{t_i}$ because $Z \ge 0$. Similarly we can argue for any $i,j \in[m]$, if $ i < j $, then $X^L_{t_j}$ stochastically dominants $X^L_{t_i}$. 
% Thus by \cite{de20211}, 
% $$
% W_1(X^L_{t_i},X^L_{t_j}) = |\EE [X^L_{t_i}] - \EE[X^L_{t_j}]|.
% $$
% For Atlanta model, we know $X^A_{t_i} \sim u$ for all $i \in[m]$, thus $W_1(X^A_{t_i},X^A_{t_j}) = 0$ for all $i,j \in [m]$. 
% \end{proof}























\bibliographystyle{plain}
\begin{thebibliography}{10}

\bibitem{agterberg2022joint}
Joshua Agterberg, Zachary Lubberts, and Jes{\'u}s Arroyo.
\newblock Joint spectral clustering in multilayer degree-corrected stochastic blockmodels.
\newblock {\em arXiv preprint arXiv:2212.05053}, 2022.

\bibitem{agterberg2022entrywise}
Joshua Agterberg, Zachary Lubberts, and Carey~E Priebe.
\newblock Entrywise estimation of singular vectors of low-rank matrices with heteroskedasticity and dependence.
\newblock {\em IEEE Transactions on Information Theory}, 68(7):4618--4650, 2022.

\bibitem{arroyo2021graph}
Jes{\'u}s Arroyo, Carey~E Priebe, and Vince Lyzinski.
\newblock Graph matching between bipartite and unipartite networks: to collapse, or not to collapse, that is the question.
\newblock {\em IEEE transactions on network science and engineering}, 8(4):3019--3033, 2021.

\bibitem{arroyo2021maximum}
Jes{\'u}s Arroyo, Daniel~L Sussman, Carey~E Priebe, and Vince Lyzinski.
\newblock Maximum likelihood estimation and graph matching in errorfully observed networks.
\newblock {\em Journal of Computational and Graphical Statistics}, 30(4):1111--1123, 2021.

\bibitem{athreya2025euclidean}
Avanti Athreya, Zachary Lubberts, Youngser Park, and Carey Priebe.
\newblock Euclidean mirrors and dynamics in network time series.
\newblock {\em Journal of the American Statistical Association}, pages 1--12, 2025.

\bibitem{BergmannHerzogJasa:2024}
Ronny Bergmann, Roland Herzog, and Hajg Jasa.
\newblock The riemannian convex bundle method.
\newblock {\em Preprint}, 2024.

\bibitem{box2015time}
George~EP Box, Gwilym~M Jenkins, Gregory~C Reinsel, and Greta~M Ljung.
\newblock {\em Time series analysis: forecasting and control}.
\newblock John Wiley \& Sons, 2015.

\bibitem{cai2021optimal}
Tony Cai, Hongzhe Li, and Rong Ma.
\newblock Optimal structured principal subspace estimation: Metric entropy and minimax rates.
\newblock {\em Journal of machine learning research}, 22(46):1--45, 2021.

\bibitem{cape2019signal}
Joshua Cape, Minh Tang, and Carey~E Priebe.
\newblock Signal-plus-noise matrix models: eigenvector deviations and fluctuations.
\newblock {\em Biometrika}, 106(1):243--250, 2019.

\bibitem{cape2019two}
Joshua Cape, Minh Tang, and Carey~E Priebe.
\newblock The two-to-infinity norm and singular subspace geometry with applications to high-dimensional statistics.
\newblock {\em Annual of Statistics}, 2019.

\bibitem{chai2024efficient}
Shuwen Chai and Mikl{\'o}s~Z R{\'a}cz.
\newblock Efficient graph matching for correlated stochastic block models.
\newblock {\em Advances in Neural Information Processing Systems}, 37:116388--116461, 2024.

\bibitem{chen_worm}
L.~Chen, J.T. Vogelstein, V.~L. Lyzinski, and C.~E. Priebe.
\newblock A joint graph inference case study: the {C. Elegans} chemical and electrical connectomes.
\newblock {\em Worm}, 5:e1142041, 2016.

\bibitem{chen2022analysis}
Rong Chen, Yuefeng Han, Zebang Li, Han Xiao, Dan Yang, and Ruofan Yu.
\newblock Analysis of tensor time series: tensorts.
\newblock {\em Journal of Statistical Software}, 323, 2022.

\bibitem{chen2024euclidean}
Tianyi Chen, Zachary Lubberts, Avanti Athreya, Youngser Park, and Carey~E Priebe.
\newblock Euclidean mirrors and first-order changepoints in network time series.
\newblock {\em arXiv preprint arXiv:2405.11111}, 2024.

\bibitem{chen2023discovering}
Tianyi Chen, Youngser Park, Ali Saad-Eldin, Zachary Lubberts, Avanti Athreya, Benjamin~D Pedigo, Joshua~T Vogelstein, Francesca Puppo, Gabriel~A Silva, Alysson~R Muotri, et~al.
\newblock Discovering a change point and piecewise linear structure in a time series of organoid networks via the iso-mirror.
\newblock {\em Applied Network Science}, 8(1):45, 2023.

\bibitem{chewi2025statistical}
Sinho Chewi, Jonathan Niles-Weed, and Philippe Rigollet.
\newblock {\em Statistical optimal transport}.
\newblock Springer, 2025.

\bibitem{conte04:_thirt}
D.~Conte, P.~Foggia, C.~Sansone, and M.~Vento.
\newblock Thirty years of graph matching in pattern recognition.
\newblock {\em International Journal of Pattern Recognition and Artificial Intelligence}, 18:265--298, 2004.

\bibitem{gms1}
Donatello Conte, Pasquale Foggia, Carlo Sansone, and Mario Vento.
\newblock Thirty years of graph matching in pattern recognition.
\newblock {\em International journal of pattern recognition and artificial intelligence}, 18(03):265--298, 2004.

\bibitem{de20211}
Marco De~Angelis and Ander Gray.
\newblock Why the 1-wasserstein distance is the area between the two marginal cdfs.
\newblock {\em arXiv preprint arXiv:2111.03570}, 2021.

\bibitem{fishkind2019seeded}
Donniell~E Fishkind, Sancar Adali, Heather~G Patsolic, Lingyao Meng, Digvijay Singh, Vince Lyzinski, and Carey~E Priebe.
\newblock Seeded graph matching.
\newblock {\em Pattern recognition}, 87:203--215, 2019.

\bibitem{gallagher2021spectral}
Ian Gallagher, Andrew Jones, and Patrick Rubin-Delanchy.
\newblock Spectral embedding for dynamic networks with stability guarantees.
\newblock {\em Advances in Neural Information Processing Systems}, 34:10158--10170, 2021.

\bibitem{racz2}
Julia Gaudio, Miklos~Z Racz, and Anirudh Sridhar.
\newblock Exact community recovery in correlated stochastic block models.
\newblock In {\em Conference on Learning Theory}, pages 2183--2241. PMLR, 2022.

\bibitem{hindes2024swarming}
Jason Hindes, Kevin Daley, George Stantchev, and Ira~B Schwartz.
\newblock Swarming network inference with importance clustering of relative interactions.
\newblock {\em Journal of Physics: Complexity}, 5(4):045009, 2024.

\bibitem{jasa2025procrustes}
Hajg Jasa, Ronny Bergmann, Christian K{\"u}mmerle, Avanti Athreya, and Zachary Lubberts.
\newblock Procrustes problems on random matrices.
\newblock {\em arXiv preprint arXiv:2510.05182}, 2025.

\bibitem{jing2021community}
Bing-Yi Jing, Ting Li, Zhongyuan Lyu, and Dong Xia.
\newblock Community detection on mixture multilayer networks via regularized tensor decomposition.
\newblock {\em The Annals of Statistics}, 49(6):3181--3205, 2021.

\bibitem{kedem2005regression}
Benjamin Kedem and Konstantinos Fokianos.
\newblock {\em Regression models for time series analysis}.
\newblock John Wiley \& Sons, 2005.

\bibitem{kivela2014multilayer}
M.~Kivel{\"a}, A.~Arenas, M.~Barthelemy, J.~P. Gleeson, Y.~Moreno, and M.~A. Porter.
\newblock Multilayer networks.
\newblock {\em Journal of complex networks}, 2(3):203--271, 2014.

\bibitem{kouachi2006eigenvalues}
Said Kouachi.
\newblock Eigenvalues and eigenvectors of tridiagonal matrices.
\newblock {\em The Electronic Journal of Linear Algebra}, 15:115--133, 2006.

\bibitem{lei2024computational}
Jing Lei, Anru~R Zhang, and Zihan Zhu.
\newblock Computational and statistical thresholds in multi-layer stochastic block models.
\newblock {\em The Annals of Statistics}, 52(5):2431--2455, 2024.

\bibitem{lyzinski16:_infoGM}
V.~Lyzinski.
\newblock Information recovery in shuffled graphs via graph matching.
\newblock {\em IEEE Transactions on Information Theory}, 64(5):3254--3273, 2018.

\bibitem{sgm_jofc}
V.~Lyzinski, S.~Adali, J.~T. Vogelstein, Y.~Park, and C.~E. Priebe.
\newblock Seeded graph matching via joint optimization of fidelity and commensurability.
\newblock Arxiv preprint at \url{http://arxiv.org/abs/1401.3813}, 2014.

\bibitem{lyzinski15:_relax}
V.~Lyzinski, D.~E. Fishkind, M.~Fiori, J.~T. Vogelstein, C.~E. Priebe, and G.~Sapiro.
\newblock Graph matching: Relax at your own risk.
\newblock {\em IEEE Transactions on Pattern Analysis and Machine Intelligence}, 38:60--73, 2016.

\bibitem{macdonald2022latent}
Peter~W MacDonald, Elizaveta Levina, and Ji~Zhu.
\newblock Latent space models for multiplex networks with shared structure.
\newblock {\em Biometrika}, 109(3):683--706, 2022.

\bibitem{segmented}
Vito~M.R. Muggeo.
\newblock segmented: an r package to fit regression models with broken-line relationships.
\newblock {\em R News}, 8(1):20--25, 2008.

\bibitem{muggeo2017interval}
Vito~MR Muggeo.
\newblock Interval estimation for the breakpoint in segmented regression: A smoothed score-based approach.
\newblock {\em Australian \& New Zealand Journal of Statistics}, 59(3):311--322, 2017.

\bibitem{padilla2022change}
Oscar Hernan~Madrid Padilla, Yi~Yu, and Carey~E Priebe.
\newblock Change point localization in dependent dynamic nonparametric random dot product graphs.
\newblock {\em The Journal of Machine Learning Research}, 23(1):10661--10719, 2022.

\bibitem{padilla2021optimal}
Oscar Hernan~Madrid Padilla, Yi~Yu, Daren Wang, and Alessandro Rinaldo.
\newblock Optimal nonparametric multivariate change point detection and localization.
\newblock {\em IEEE Transactions on Information Theory}, 68(3):1922--1944, 2021.

\bibitem{pantazis2022importance}
Konstantinos Pantazis, Avanti Athreya, Jes{\'u}s Arroyo, William~N Frost, Evan~S Hill, and Vince Lyzinski.
\newblock The importance of being correlated: Implications of dependence in joint spectral inference across multiple networks.
\newblock {\em Journal of Machine Learning Research}, 23(141):1--77, 2022.

\bibitem{paul2020spectral}
S.~Paul and Y.~Chen.
\newblock Spectral and matrix factorization methods for consistent community detection in multi-layer networks.
\newblock {\em The Annals of Statistics}, 48(1):230--250, 2020.

\bibitem{pensky2025signed}
Marianna Pensky.
\newblock Signed diverse multiplex networks: clustering and inference.
\newblock {\em IEEE Transactions on Information Theory}, 2025.

\bibitem{priebe2015statistical}
Carey~E Priebe, Daniel~L Sussman, Minh Tang, and Joshua~T Vogelstein.
\newblock Statistical inference on errorfully observed graphs.
\newblock {\em Journal of Computational and Graphical Statistics}, 24(4):930--953, 2015.

\bibitem{qi2025asymptotically}
Tong Qi, Vera Andersson, Peter Viechnicki, and Vince Lyzinski.
\newblock Asymptotically perfect seeded graph matching without edge correlation (and applications to inference).
\newblock {\em arXiv preprint arXiv:2506.02825}, 2025.

\bibitem{qiao2021igraphmatch}
Zihuan Qiao and Daniel Sussman.
\newblock igraphmatch: an r package for the analysis of graph matching.
\newblock {\em arXiv preprint arXiv:2112.09212}, 2021.

\bibitem{racz2021correlated}
Miklos Racz and Anirudh Sridhar.
\newblock Correlated stochastic block models: Exact graph matching with applications to recovering communities.
\newblock {\em Advances in Neural Information Processing Systems}, 34:22259--22273, 2021.

\bibitem{rubin2022statistical}
Patrick Rubin-Delanchy, Joshua Cape, Minh Tang, and Carey~E Priebe.
\newblock A statistical interpretation of spectral embedding: The generalised random dot product graph.
\newblock {\em Journal of the Royal Statistical Society Series B: Statistical Methodology}, 84(4):1446--1473, 2022.

\bibitem{saad2021graph}
Ali Saad-Eldin, Benjamin~D Pedigo, Carey~E Priebe, and Joshua~T Vogelstein.
\newblock Graph matching via optimal transport.
\newblock {\em arXiv preprint arXiv:2111.05366}, 2021.

\bibitem{saxena2025lost}
Ayushi Saxena and Vince Lyzinski.
\newblock Lost in the shuffle: Testing power in the presence of errorful network vertex labels.
\newblock {\em Computational Statistics \& Data Analysis}, 204:108091, 2025.

\bibitem{shumway2000time}
Robert~H Shumway, David~S Stoffer, and David~S Stoffer.
\newblock {\em Time series analysis and its applications}, volume~3.
\newblock Springer, 2000.

\bibitem{gms3}
Hui Sun, Wenju Zhou, and Minrui Fei.
\newblock A survey on graph matching in computer vision.
\newblock In {\em 2020 13th International Congress on Image and Signal Processing, BioMedical Engineering and Informatics (CISP-BMEI)}, pages 225--230. IEEE, 2020.

\bibitem{STFP-2011}
D.~L. Sussman, M.~Tang, D.~E. Fishkind, and C.~E. Priebe.
\newblock A consistent adjacency spectral embedding for stochastic blockmodel graphs.
\newblock {\em Journal of the American Statistical Association}, 107:1119--1128, 2012.

\bibitem{tang2025eigenvector}
Minh Tang and Joshua~R Cape.
\newblock Eigenvector fluctuations and limit results for random graphs with infinite rank kernels.
\newblock {\em arXiv preprint arXiv:2501.15725}, 2025.

\bibitem{isomap_science}
J.~B. Tenenbaum, V.~de~Silva, and J.C. Langford.
\newblock A global geometric framework for nonlinear dimensionality reduction.
\newblock {\em Science}, 290:2319---2323, 2000.

\bibitem{verzelen2023optimal}
Nicolas Verzelen, Magalie Fromont, Matthieu Lerasle, and Patricia Reynaud-Bouret.
\newblock Optimal change-point detection and localization.
\newblock {\em The Annals of Statistics}, 51(4):1586--1610, 2023.

\bibitem{vogelstein2015shuffled}
Joshua~T Vogelstein and Carey~E Priebe.
\newblock Shuffled graph classification: Theory and connectome applications.
\newblock {\em Journal of Classification}, 32(1):3--20, 2015.

\bibitem{wang2023multilayer}
Fan Wang, Wanshan Li, Oscar Hernan~Madrid Padilla, Yi~Yu, and Alessandro Rinaldo.
\newblock Multilayer random dot product graphs: Estimation and online change point detection.
\newblock {\em arXiv preprint arXiv:2306.15286}, 2023.

\bibitem{xie2024entrywise}
Fangzheng Xie.
\newblock Entrywise limit theorems for eigenvectors of signal-plus-noise matrix models with weak signals.
\newblock {\em Bernoulli}, 30(1):388--418, 2024.

\bibitem{gms2}
Junchi Yan, Xu-Cheng Yin, Weiyao Lin, Cheng Deng, Hongyuan Zha, and Xiaokang Yang.
\newblock A short survey of recent advances in graph matching.
\newblock In {\em Proceedings of the 2016 ACM on international conference on multimedia retrieval}, pages 167--174, 2016.

\bibitem{zhang2018tensor}
Anru Zhang and Dong Xia.
\newblock Tensor svd: Statistical and computational limits.
\newblock {\em IEEE Transactions on Information Theory}, 64(11):7311--7338, 2018.

\bibitem{zhu2006automatic}
Mu~Zhu and Ali Ghodsi.
\newblock Automatic dimensionality selection from the scree plot via the use of profile likelihood.
\newblock {\em Computational Statistics \& Data Analysis}, 51(2):918--930, 2006.

\bibitem{zuzul2021dynamic}
Tiona Zuzul, Emily~Cox Pahnke, Jonathan Larson, Patrick Bourke, Nicholas Caurvina, Neha~Parikh Shah, Fereshteh Amini, Jeffrey Weston, Youngser Park, Joshua Vogelstein, et~al.
\newblock Dynamic silos: Increased modularity in intra-organizational communication networks during the covid-19 pandemic.
\newblock {\em arXiv preprint arXiv:2104.00641}, 2021.

\end{thebibliography}



\end{document}
